% ============================================================
% RVCE Major Project Report
% Title: Quantifying and Reducing Semantic Drift in Text-to-SQL Systems
% Template: ecproject.sty v3.9
% ============================================================

\documentclass[12pt,a4paper]{report}

% ── Core template ────────────────────────────────────────────────────────────
\usepackage{ecproject}

% ── Typography ───────────────────────────────────────────────────────────────
\usepackage{times}
\usepackage[T1]{fontenc}
\usepackage[utf8]{inputenc}
\usepackage{setspace}
\usepackage{microtype}

% ── Page layout ──────────────────────────────────────────────────────────────
\usepackage[top=1in,bottom=1in,left=1.25in,right=1in]{geometry}

% ── Mathematics ──────────────────────────────────────────────────────────────
\usepackage{amsmath}
\usepackage{amssymb}

% ── Tables ───────────────────────────────────────────────────────────────────
\usepackage{booktabs}
\usepackage{tabularx}
\usepackage{multirow}
\usepackage{longtable}
\usepackage{array}

% ── Figures ──────────────────────────────────────────────────────────────────
\usepackage{graphicx}
\usepackage{float}
\usepackage{caption}
\usepackage{subcaption}

% ── Listings and code ────────────────────────────────────────────────────────
\usepackage{tcolorbox}
\tcbuselibrary{listings,skins,breakable}
\usepackage{listings}
\usepackage{xcolor}

% ── Bibliography ─────────────────────────────────────────────────────────────
\usepackage{cite}

% ── Abbreviations ────────────────────────────────────────────────────────────
\usepackage{acronym}

% ── Miscellaneous ────────────────────────────────────────────────────────────
\usepackage{enumitem}
\usepackage{parskip}

% ── TiKZ (required by ecproject and Abstract) ────────────────────────────────
\usepackage{tikz}
\usetikzlibrary{calc}

% ── Cross-references (loaded last) ───────────────────────────────────────────
\usepackage{hyperref}
\hypersetup{
  colorlinks=true,
  linkcolor=black,
  citecolor=black,
  urlcolor=blue,
  pdftitle={Quantifying and Reducing Semantic Drift in Text-to-SQL Systems},
  pdfauthor={RVCE}
}

% ── Listing style ────────────────────────────────────────────────────────────
\definecolor{codebg}{gray}{0.95}
\definecolor{codeframe}{gray}{0.7}
\definecolor{kwcolor}{HTML}{0000AA}
\definecolor{cmtcolor}{HTML}{555555}
\definecolor{strcolor}{HTML}{007700}

\lstset{
  basicstyle=\small\ttfamily,
  keywordstyle=\color{kwcolor}\bfseries,
  commentstyle=\color{cmtcolor}\itshape,
  stringstyle=\color{strcolor},
  breaklines=true,
  showstringspaces=false,
  numbers=left,
  numberstyle=\tiny\color{gray},
  frame=single,
  backgroundcolor=\color{codebg},
  captionpos=b
}

% ── Project metadata (required by ecproject.sty) ─────────────────────────────
\title{Quantifying and Reducing Semantic Drift in Text-to-SQL Systems}
\subCode{22MCA}
\academicYear{2025--26}
\Department[CSE]{Computer Science and Engineering}

% ── Spacing ──────────────────────────────────────────────────────────────────
\onehalfspacing

% =============================================================================
\begin{document}
% =============================================================================

% ── Abstract ─────────────────────────────────────────────────────────────────
\addcontentsline{toc}{chapter}{Abstract}\vspace{-1cm}
%Border
\begin{tikzpicture}[remember picture, overlay]
  \draw[line width = 4pt] ($(current page.north west) + (0.75in,-0.75in)$) rectangle ($(current page.south east) + (-0.75in,0.75in)$);
\end{tikzpicture}
Abstract is essentially a complete overview or the Synopsis of the project in one page, should be precise, concise, and clear. It has to be written in past tense as the work is already completed. If abbreviations or acronyms need to be included in the abstract, the first time it is used, it should be written out in the full form and put the abbreviation in brackets. e.g. Magnetic Resonance Imaging (MRI). The abstract may have about 3 paragraphs covering the outline of the work as below:

\textbf{First para:} Brief introduction about the project which should include background, Areas/type of applications, market size (if available) and hence the importance of the domain. Key developments in the domain, unresolved issues or/and new emerging opportunities in the domain areas and hence the problem or the issue that is proposed to be addressed in the project or the objective. If it is theory or simulation based add all the existing theory/algorithms names and what is it that the project is based on, their essentially advantage or drawbacks in brief and hence the motivation to work on the project.

\textbf{Second Para:} Objectives, scope, Methodology of your project including the key theory, design tools or software tools or / and the experimental tools that will be used, the key specifications including constrains or boundary conditions or approximations or assumptions that have been considered for the work. The sequence in which the work has been carried out.

\textbf{Third para:} Key findings, results or outcome of the work including experimental data of the out come or / and percentage in terms of working efficiency of the developed product or system.

In brief the Abstract should be
\begin{itemize}
	\item Complete — it covers all aspects of the project
	\item Concise — it contains no excess wordiness or unnecessary information.
	\item Clear — it is readable, well organized, and not too jargon-laden.
	\item Cohesive — it flows smoothly between the parts.
\end{itemize}

\pagebreak

% ── Front matter lists ───────────────────────────────────────────────────────
\tableofcontents
\listoffigures
\listoftables

% ── List of Abbreviations ────────────────────────────────────────────────────
\chapter*{List of Abbreviations}
\addcontentsline{toc}{chapter}{List of Abbreviations}
\begin{acronym}[NLIDB]
  \acro{NL}{Natural Language}
  \acro{NLIDB}{Natural Language Interface to Databases}
  \acro{SQL}{Structured Query Language}
  \acro{NLP}{Natural Language Processing}
  \acro{LLM}{Large Language Model}
  \acro{API}{Application Programming Interface}
  \acro{REST}{Representational State Transfer}
  \acro{SDG}{Sustainable Development Goal}
  \acro{ORM}{Object-Relational Mapping}
  \acro{UI}{User Interface}
  \acro{AST}{Abstract Syntax Tree}
  \acro{BERT}{Bidirectional Encoder Representations from Transformers}
  \acro{RL}{Reinforcement Learning}
  \acro{YoY}{Year-over-Year}
  \acro{PII}{Personally Identifiable Information}
  \acro{KPI}{Key Performance Indicator}
  \acro{DAG}{Directed Acyclic Graph}
\end{acronym}

% =============================================================================
% CHAPTER 1 — INTRODUCTION
% =============================================================================
\chapter{Introduction}
\label{ch:introduction}

\section{Background and Motivation}
\label{sec:background}

The ability to query a database using everyday language has been a long-standing goal in computer science. Organisations generate enormous volumes of structured data, yet the skill required to write even moderately complex \ac{SQL} remains a barrier that keeps most of that data out of reach for non-technical decision-makers. Over the past decade, advances in \ac{NLP} and the availability of large pre-trained language models have revived interest in systems that translate natural language questions directly into executable \ac{SQL} — a task broadly referred to as Text-to-SQL.

Early approaches to the problem relied on hand-crafted grammars and keyword mapping \cite{seq2sql}. These rule-based systems worked well within carefully defined domains but broke down quickly when confronted with phrasing variations or schema diversity. Sequence-to-sequence architectures offered a more flexible alternative \cite{sqlnet}, learning patterns from paired (question, SQL) training corpora. The release of the Spider benchmark in 2018 \cite{spider} established a cross-domain evaluation standard and accelerated model development significantly: top-performing models on Spider now achieve execution accuracies above 85\% on well-posed questions drawn from the test distribution.

Despite this progress, production deployments of Text-to-SQL systems continue to produce unexpected outputs. A query phrased as ``show me total revenue by region excluding taxes'' may yield a result that ignores the tax-exclusion clause, groups by city rather than region, or applies an incorrect date window. Each of these failures represents a divergence between the user's original intent and the SQL actually executed — a phenomenon this project calls \textit{semantic drift}. Unlike a hard parsing error, semantic drift often goes undetected: the query runs, numbers are returned, and a decision may be made on data that does not answer the original question.

The commercial stakes are substantial. Market research by Gartner and McKinsey consistently estimates that poor data quality costs organisations between \$9.7 million and \$15 million per year on average, with misinterpreted query results being among the most insidious causes. Small errors in filter conditions or aggregation logic compound rapidly when rolled up into dashboards and executive reports. The problem is especially acute in e-commerce and financial reporting, where revenue, margin, and return-rate figures directly inform pricing, inventory, and headcount decisions.

The insidiousness of semantic drift — as opposed to hard query failure — is that it produces confident-looking results. A query that returns zero rows, raises an exception, or generates a syntax error is immediately recognisable as wrong. A query that returns four rows of revenue figures, grouped by region, with a plausible total around the expected order of magnitude, looks entirely correct. A business analyst who receives such a result has no obvious signal to check further; the error propagates silently into the report, the slide deck, and the decision. This is the failure mode that motivates the project: not the noisy, visible failure, but the quiet, confident, wrong answer.

The technical challenge of detecting this failure mode is substantial. Unlike a corrupted data record, which can in principle be flagged by a schema constraint or a null check, a semantically drifted query result is internally consistent: the rows are valid, the types match, the aggregation executes without error. Detection requires external reference information — either a ground-truth SQL (unavailable in production), domain knowledge about what results should look like (partially available through historical baselines), or a formal decomposition of the user's intent and the system's actions (the approach taken in this project).

The volume of structured data generated by modern organisations has grown at a pace that outstrips the supply of data analysts qualified to query it. According to industry surveys, less than 20\% of employees at data-driven organisations have sufficient SQL proficiency to write queries independently; the remainder depend on a small pool of analysts and data engineers to translate their information needs into database queries, creating bottlenecks and delays. Natural language database interfaces address this directly by allowing any user to express an information need in the same way they would ask a colleague: in everyday language, without schema knowledge or query syntax.

The industry shift from batch ETL pipelines to conversational data access has accelerated this concern. Products such as Tableau Ask Data, Microsoft Power BI Q\&A, and Amazon QuickSight Q\&A are now embedded in widely deployed analytics platforms, meaning that millions of business users interact with relational databases through natural language daily. At this scale, even a 5\% rate of semantically incorrect queries — each of which may silently influence a business decision — represents a substantial quality risk. The challenge is compounded by the fact that \ac{LLM}-era Text-to-SQL models sometimes produce plausible-looking \ac{SQL} that contains hallucinated column names, ignores negation words such as ``excluding'' or ``except,'' or applies the wrong aggregation function when multiple metrics share similar surface forms.

This project was motivated by the observation that existing Text-to-SQL frameworks provide no principled mechanism for detecting or recovering from these silent failures. Most systems return a result or an error — there is no middle ground that says ``this result is plausible but semantically suspect.'' Introducing such a mechanism, and wiring it into a feedback loop that attempts to correct the query before the result is served, is the central contribution of this work.

The gap between benchmark accuracy and production reliability can be quantified roughly as follows. A Text-to-SQL model with 85\% exact-match accuracy on Spider will produce incorrect SQL on approximately 15\% of queries. In a Spider evaluation context, these failures are counted, analysed, and used to drive model improvement. In a production context, deployed to a user base of 1,000 business analysts each submitting 5 queries per day, the same error rate means approximately 750 incorrect results per day — each of which could silently influence a business decision before anyone checks the underlying data. At this volume, even a 1\% reduction in the silent failure rate has material business value. The drift metric and critic loop developed in this project are designed precisely for this operational context: reducing the rate of undetected incorrect results in a deployed system, even when the underlying model cannot be retrained.

\section{Problem Statement}
\label{sec:problem}

Given a natural language question $q$ from a business user and a relational database schema $\mathcal{S}$, a Text-to-SQL system $f$ generates a \ac{SQL} query $Q = f(q, \mathcal{S})$ and executes it to obtain a result set $R$. The core problem is that $Q$ may not faithfully represent the user's intent $I(q)$ even when $R$ is non-empty and syntactically valid.

Three categories of drift were identified empirically during development:

\begin{itemize}[leftmargin=2em]
  \item \textbf{Entity drift} — the system maps a user's mention of ``clients'' to the \texttt{orders} table rather than the \texttt{customers} table, returning a count of transactions instead of a count of people.
  \item \textbf{Constraint drift} — a filter condition such as ``excluding taxes'' or ``active products only'' is recognised by the intent parser but omitted from the generated \ac{SQL}, silently inflating the result.
  \item \textbf{Plausibility drift} — the returned aggregate value is statistically inconsistent with historical baselines (e.g., a total revenue figure three standard deviations above the mean for an identical query run the previous week), suggesting a join explosion or a misconfigured \texttt{WHERE} clause.
\end{itemize}

To make the problem concrete, consider a specific instance. The query ``show me total revenue by region excluding taxes'' submitted to a single-pass pattern-matching system might produce a result with $a_i = 0.73$ (the region grouping was correctly identified but the tax-exclusion flag was missed), $a_c = 0.67$ (2 of 3 applicable constraints satisfied), and $a_p = 0.80$ (result values are plausible). These component scores yield a composite drift of $D = 0.4(1-0.73) + 0.3(1-0.67) + 0.3(1-0.80) = 0.108 + 0.099 + 0.060 = 0.267$. At this level of drift, the query is not yet safe to serve; the critic loop would re-invoke the pipeline with explicit feedback that the TAX\_EXCLUSION constraint was not satisfied, prompting a revised \ac{SQL} that subtracts \texttt{o.tax\_amount} from the aggregation.

Table~\ref{tab:drift-failure-modes} illustrates the three categories of drift with prototypical examples.

\begin{table}[H]
  \centering
  \caption{Semantic Drift Failure Mode Categories with Examples}
  \label{tab:drift-failure-modes}
  \renewcommand{\arraystretch}{1.3}
  \begin{tabularx}{\textwidth}{@{}lXX@{}}
    \toprule
    \textbf{Drift Type} & \textbf{User Intent} & \textbf{Actual SQL Behaviour} \\
    \midrule
    Entity drift       & ``list clients'' → COUNT customers & COUNT orders returned instead \\
    Constraint drift   & ``excluding taxes'' → net revenue  & Gross revenue returned, tax not subtracted \\
    Plausibility drift & ``revenue last month'' → \$2.1M   & \$21M returned (join explosion with order\_items) \\
    \bottomrule
  \end{tabularx}
\end{table}

To make the problem tractable, a composite \textit{semantic drift score} $D \in [0,1]$ was defined, where $D = 0$ represents perfect fidelity to the user's intent and $D = 1$ represents complete misalignment. A query result is considered safe to serve when $D < 0.15$. The derivation of this score and its three component dimensions form the primary theoretical contribution of this project.

\section{Objectives}
\label{sec:objectives}

The project set out to achieve the following objectives:

\begin{enumerate}[leftmargin=2em,label=\arabic*.]
  \item Design and formalise a composite semantic drift metric that decomposes intent fidelity into measurable, independently computable components — specifically intent alignment, constraint adherence, and result plausibility — and assign weights to each component that reflect their relative contribution to downstream query failure.
  \item Build a business ontology for an e-commerce domain that encodes entities, relationships, metric definitions, and constraint rules in a machine-navigable graph using \texttt{NetworkX}, enabling shortest-path join inference and deterministic constraint lookup.
  \item Implement a five-agent pipeline — IntentParser, OntologyMapper, ConstraintValidator, ExecutionPlanner, ResultVerifier — each responsible for a well-defined stage of query translation, with a shared \texttt{QueryState} typed dataclass as the communication medium.
  \item Develop a critic-loop orchestrator that iterates the pipeline, recomputing drift after each pass and refining the query until either $D < 0.15$ or a maximum of five iterations is exhausted, using component-specific feedback messages to guide each retry.
  \item Expose the system through a \ac{REST} \ac{API} (eight endpoints) and a real-time monitoring dashboard, demonstrating integration readiness for production environments and providing full execution traces in all responses.
  \item Evaluate the system on a representative set of ten business queries, quantifying reduction in drift across iterations and comparing against a single-pass baseline, using a synthetic e-commerce dataset at realistic production scale.
\end{enumerate}

Each objective corresponds to a measurable deliverable: the drift metric is specified in Equations~\ref{eq:intent-alignment}--\ref{eq:convergence}; the ontology is documented in Tables~\ref{tab:entities}--\ref{tab:join-rules}; the pipeline is described in Section~\ref{sec:pipeline-design}; the critic loop is described in Section~\ref{sec:critic-loop}; the API and dashboard are documented in Sections~\ref{sec:api-design} and~\ref{sec:dashboard-design}; and the evaluation results are presented in Chapter~\ref{ch:results}.

\section{Scope and Limitations}
\label{sec:scope}

The system targets analytical (read-only) queries over a synthetic e-commerce dataset spanning five relational databases — sales, inventory, analytics, HR, and finance — totalling approximately 1.5 million rows. The ontology covers eleven business entities, twelve aggregate metrics, and nine constraint rules, which is sufficient for a representative range of business intelligence questions but does not exhaust the combinatorial space of possible queries.

Several deliberate boundaries were maintained throughout the project. The system currently uses \ac{SQL}ite as its production backend; a PostgreSQL dialect is implemented and tested but not deployed, since running a full PostgreSQL server in the evaluation environment was outside scope. Semantic similarity in the OntologyMapper relies on Python's built-in \texttt{SequenceMatcher} rather than learned embeddings; the architecture includes a ChromaDB integration path for future replacement. Natural language understanding uses pattern matching across 8 entity patterns, 12 metric patterns, and more than 15 filter detectors — a design choice that makes the system deterministic and auditable but less flexible than a neural parser on out-of-distribution phrasing.

Three categories of query are explicitly out of scope. First, data definition language (DDL) statements — CREATE, ALTER, DROP — are not supported; the system is read-only by design and provides no mechanism for schema modification. The read-only constraint was enforced at the \ac{API} layer by validating that the generated SQL does not contain any DDL keywords before execution; any attempt to submit a query that would produce DDL SQL is rejected at the ConstraintValidatorAgent stage. Second, multi-hop federated joins that cross the boundaries between databases owned by different schema authorities are beyond the current system's join-inference logic, which assumes a single consolidated schema. Third, real-time streaming data is not addressed: the synthetic dataset is static, and the plausibility baselines are computed from historical snapshots rather than live distribution estimates. Each of these exclusions was a deliberate trade-off that kept the evaluation scope manageable within a single academic project cycle.

A practical consequence of the read-only constraint is that the system cannot be used to generate INSERT or UPDATE statements for data entry workflows, even though natural language data entry is a common enterprise use case. This constraint was included in the system's documentation and surfaced in the API's \texttt{/rules} endpoint response, so that client applications could inform users of the limitation. Future work could add a separate ``data mutation'' pipeline with appropriate authorisation controls, operating in parallel with the query pipeline, but this would constitute a substantially different system with different safety requirements.

\section{Alignment with SDG~9}
\label{sec:sdg}

\ac{SDG}~9 calls for building resilient infrastructure, promoting inclusive and sustainable industrialisation, and fostering innovation. The \ac{NLIDB} system developed in this project contributes directly to the ``inclusive innovation'' strand: by removing the \ac{SQL} skill barrier, it enables smaller enterprises and non-technical teams to extract value from their operational data without depending on specialist database engineers. The drift-detection mechanism further improves the reliability of that access, reducing the risk that a misinterpreted query propagates a bad decision into business operations. To a lesser extent, the project also touches \ac{SDG}~12 (responsible consumption and production) by enabling more precise inventory and supply chain analytics.

The connection to \ac{SDG}~12 is worth elaborating. Responsible consumption, in the context of a product-selling enterprise, depends critically on accurate demand forecasting, inventory management, and supplier performance monitoring. Each of these activities relies on database queries whose accuracy is difficult to verify manually at scale. A system that can flag its own uncertain results — rather than silently serving a number derived from a mis-scoped query — contributes to better-informed procurement and stock management decisions. While the present system was evaluated on synthetic e-commerce data rather than a live deployment, the inventory and supply chain analytics queries in the test set demonstrated that the drift metric is sensitive to the kinds of filter omissions most likely to affect stock-level calculations.

The contribution to \ac{SDG}~9 extends beyond the removal of the SQL skill barrier. The system's architecture — a five-agent pipeline orchestrated by a typed state graph — is itself an example of resilient, modular software infrastructure. By decomposing a complex NLP task into independently deployable agents with well-defined interfaces, the design allows individual components to be upgraded or replaced as better algorithms become available, without disrupting the rest of the system. This incremental upgradeability is a property that characterises resilient infrastructure in the sustainable development sense: it does not require a full rebuild when circumstances change, reducing the long-term maintenance burden and enabling sustained innovation on a stable foundation.

At the institutional level, systems like the one developed in this project contribute to the democratisation of data infrastructure. When a mid-size enterprise in an emerging economy gains reliable natural language access to its operational data — without the budget for a large data engineering team — it is able to compete more effectively with larger enterprises that can afford dedicated SQL expertise. This reduction in the data-access cost barrier is a form of inclusive innovation that aligns with \ac{SDG}~9's emphasis on making industrial capabilities available to smaller actors.

\section{Organisation of the Report}
\label{sec:organisation}

Chapter~\ref{ch:literature} surveys the evolution of Text-to-SQL systems from the 1970s through to the current \ac{LLM} era, discusses prior work on semantic drift and related faithfulness evaluation, and identifies the three gaps that motivated this project's specific design choices. Chapter~\ref{ch:design} presents the full system architecture, ontology design, agent pipeline design, semantic drift metric formulation, and technology selection rationale. Chapter~\ref{ch:implementation} documents the implementation in detail, with references to specific source files, function names, and line numbers. Chapter~\ref{ch:results} presents experimental results across ten test queries, including a full worked end-to-end trace and quantitative comparisons between the critic-loop pipeline and the single-pass baseline. Chapter~\ref{ch:conclusion} summarises the four contributions, discusses key findings, and outlines six directions for future work.

The report is structured to be read sequentially, but each chapter is also self-contained to the extent possible: readers primarily interested in the drift metric formulation can read Sections~\ref{sec:drift-metric} and~\ref{sec:drift-impl} independently; readers interested in evaluation methodology and results can begin at Chapter~\ref{ch:results} and refer back to Sections~\ref{sec:pipeline-design} and~\ref{sec:ontology-design} as needed.

% =============================================================================
% CHAPTER 2 — LITERATURE REVIEW
% =============================================================================
\chapter{Literature Review}
\label{ch:literature}

\section{Natural Language Interfaces to Databases --- A Historical Perspective}
\label{sec:nlidb-history}

Research into \acp{NLIDB} stretches back to the 1970s. LUNAR, developed at Bolt Beranek and Newman, was among the first operational systems, answering queries about Apollo moon rock samples using a hand-coded semantic grammar. BASEBALL and STUDENT followed a similar pattern: narrow domain, expert-authored rules, brittle transfer. These early systems demonstrated the conceptual feasibility of natural language database access but made clear that scalability was the fundamental challenge.

Through the 1980s and 1990s, researchers shifted toward more principled linguistic representations. CHAT-80, ELF, and MASQUE were built around compositional semantics and lambda calculus, translating questions into formal logical expressions before mapping to \ac{SQL}. The advantage was that the same logical intermediate could, in principle, target multiple database backends. The disadvantage was that building and maintaining the lexical knowledge bases required significant expert effort. Systems of this era typically handled a few hundred question types per domain; moving to a new schema required months of re-engineering.

MASQUE (Modular Answering System using Quantified English) deserves specific mention because its architecture bears a family resemblance to the five-agent pipeline in this project. MASQUE decomposed the translation process into a sequence of distinct stages: a morphological analyser, a syntactic parser, a semantic interpreter, and a database interface. Each stage had a well-defined input and output representation, enabling debugging by inspection at each boundary. The primary difference from the present work is that MASQUE's stages were designed for sequential linguistic processing and carried no quality metric — there was no mechanism for measuring how faithfully the final database query reflected the original question. The five-agent pipeline in this project can be seen as a modern variant of this compositional design, augmented with a continuous quality signal and a feedback loop.

ELF (English Language Front-end) and similar commercial interfaces of the late 1980s demonstrated that a limited but reliable natural language interface was commercially viable for narrow-domain applications. ELF shipped with IBM's DB2 and provided a constrained-English query interface that allowed users to ask questions without learning SQL syntax, provided they phrased their questions in patterns the system recognised. Its success in narrow domains illustrated that the scalability problem, not the conceptual feasibility, was the primary obstacle to broad deployment — a conclusion that remains relevant today.

The arrival of corpus-based \ac{NLP} in the late 1990s and early 2000s opened the possibility of learning the translation mapping from annotated data rather than encoding it by hand. However, progress was slow: annotating (question, \ac{SQL}) pairs at scale is expensive, and early datasets such as ATIS and GeoQuery, while valuable, covered only a single domain each. The field remained a niche academic pursuit until the introduction of large-scale benchmarks in the late 2010s.

Commercial \ac{NLIDB} products also appeared during the 2000s with varying degrees of success. EasyAsk and BusinessObjects Natural Language Query offered English-language query interfaces for enterprise databases, primarily targeting business users who were comfortable phrasing questions in structured English but not in \ac{SQL}. These products achieved moderate adoption in corporate environments but were ultimately limited by the same ontology-maintenance burden that affected academic systems: each new data domain required significant configuration by domain experts. The inability to generalise across domains without re-configuration was the defining weakness of all pre-deep-learning \ac{NLIDB} approaches.

The adoption trajectories of these commercial products offer a lesson for the present work. EasyAsk was acquired and discontinued; BusinessObjects NLQ survived through integration into SAP's analytics platform but remained a niche feature rather than a widely used capability. The pattern is consistent: narrow-domain NLIDBs achieved early adoption but failed to sustain it as user expectations grew beyond the system's configured scope. The failure mode was typically a gradual accumulation of unanswered questions — queries that the system did not recognise and returned an error for — rather than a dramatic single point of failure. Over time, users learned to avoid the system, working around it with SQL or standard reports, negating its original purpose.

This adoption failure mode is relevant to the present project because it motivates the design choices around graceful degradation and drift scoring. A system that confidently returns wrong answers (like pre-deep-learning NLIDBs when queries stray from their configured scope) will eventually be trusted less than a system that transparently communicates its uncertainty. By returning a drift score alongside every result, the system allows users to make calibrated trust decisions: queries with low drift scores can be acted on directly; queries with high drift scores warrant manual verification. This creates a more sustainable human-system collaboration than either ``always trust'' or ``never trust'' postures.

The deep learning era, beginning around 2016, fundamentally changed the economics of NLIDB deployment. For the first time, a single model could be trained once on a large annotated dataset and achieve reasonable accuracy across many domains, without per-domain ontology configuration. This dramatically reduced the deployment barrier. However, as argued in Section~\ref{sec:literature-gaps}, the accuracy improvements came with a new failure mode — silent semantic drift — that the previous era's explicit constraint systems actually handled better in some respects. A hybrid approach that combines a neural model's broad coverage with an ontology system's constraint enforcement and quality measurement is the natural synthesis; the architecture in this project is one instantiation of that hybrid, weighted toward the ontology side given the auditability constraint.

\section{Evolution of Text-to-SQL Systems}
\label{sec:text-to-sql-evolution}

The modern era of Text-to-SQL begins, arguably, with the release of Seq2SQL in 2017 \cite{seq2sql}. Zhong et al.\ framed the task as a sequence-to-sequence generation problem with reinforcement learning: the model was trained to maximise execution reward on the WikiSQL dataset, producing a \texttt{SELECT}, \texttt{WHERE}, and \texttt{AGG} triplet. WikiSQL's 80,654 question-SQL pairs were unprecedented in scale, and Seq2SQL demonstrated for the first time that end-to-end neural models could achieve competitive accuracy without hand-written grammars.

Xu et al.\ followed in the same year with SQLNet \cite{sqlnet}, replacing the \ac{RL} training signal with a sketch-based generation approach. Rather than decoding \ac{SQL} token by token — a process prone to exposure bias — SQLNet predicted each slot of a pre-defined query template independently. This decomposition reduced the search space and improved accuracy on WikiSQL to 68.0\% (execution match), with no \ac{RL} required. The slot-filling paradigm influenced many subsequent architectures.

TypeSQL and SyntaxSQLNet extended SQLNet's slot-filling approach by adding type information (column types, schema-linked entities) and structural grammar constraints respectively. TypeSQL used a type system to restrict candidate column predictions: when a slot is expected to hold a numeric column, non-numeric columns are excluded from the candidate set, substantially reducing the search space. SyntaxSQLNet encoded the full SQL grammar as a tree structure, generating SQL in a top-down tree traversal order that enforced syntactic validity by construction. Both approaches improved on WikiSQL but still suffered from the single-table limitation: WikiSQL queries involve at most one table join, and the architectures were not designed to generalise to the multi-table, multi-join queries that characterise real business intelligence scenarios.

The Spider benchmark \cite{spider} represented a qualitative step change in 2018. Unlike WikiSQL, Spider spans 200 databases across 138 domains and requires models to handle cross-domain generalisation, multi-table joins, nested subqueries, and set operations. Early models achieved only 12–14\% exact match on Spider's test set, highlighting how far the field was from general-purpose natural language database access. The benchmark catalysed a wave of research into schema linking, execution-guided decoding, and pre-trained language model fine-tuning.

Spider's evaluation methodology — comparing generated SQL to a human-written reference using exact-match or execution-match criteria — was a substantial practical advance over earlier evaluations. However, it is only applicable when a reference SQL exists. In production deployment, queries are ad hoc, and maintaining a reference SQL for every possible question a business user might ask is impractical. The failure cases not caught by Spider-style evaluation — queries that produce a result but reflect the wrong intent — are precisely the cases this project's drift metric is designed to detect.

\ac{BERT}-based models \cite{spider} closed the gap significantly. IRNet, RAT-SQL, and SDSQL all introduced mechanisms for encoding the question and schema jointly, enabling the model to align question tokens with column and table names. RAT-SQL in particular introduced a relation-aware transformer that modelled question-to-column, question-to-table, and column-to-column relationships explicitly through a graph attention mechanism, reaching 69.7\% exact match on Spider. BRIDGE followed a different path, serialising the schema as a special input sequence and allowing the encoder to attend across question tokens and schema tokens simultaneously.

By 2021, T5-based fine-tuned models exceeded 70\% exact match on Spider. The PICARD approach introduced constrained decoding, preventing the model from generating syntactically invalid \ac{SQL} tokens, which improved both accuracy and reliability. More recent work with GPT-4 and Claude-class \acp{LLM}, prompted with few-shot examples and schema descriptions, has pushed above 85\%. Prompting strategies such as DAIL-SQL and DIN-SQL use decomposition: the model first identifies relevant tables, then selects columns, then constructs the full query — a chain-of-thought approach that reduces the incidence of hallucinated column names. At these accuracy levels, production deployment appears within reach — yet silent semantic drift, particularly on constrained or filtered queries, remains largely unaddressed.

\section{Semantic Parsing Approaches}
\label{sec:semantic-parsing}

Semantic parsing — converting natural language into a formal meaning representation — provides the theoretical foundations on which Text-to-SQL is built. Early semantic parsers targeted logical forms such as lambda calculus or dependency-based compositional semantics \cite{seq2sql}. These representations are expressive but require significant supervision to learn from data.

Task-oriented dialogue parsers, developed for systems like Alexa and Google Assistant, take a shallower approach: they predict a flat frame consisting of an intent label and a set of slot-value pairs. This design is well-suited to narrow-domain commands (``set a timer for five minutes'') but struggles with the relational complexity of database queries, where a single question may involve multiple table joins, aggregations, and filter chains.

The sketch-based approach in SQLNet is, in some respects, a task-specific semantic frame: it predicts which aggregate function to use, which columns to select, and which conditions to apply, treating each as an independent classification decision. Later models such as IGSQL and ShadowGNN incorporated graph neural networks to encode the relational structure of the schema, allowing multi-hop joins to be predicted without exhaustive enumeration. The architecture in this project draws on sketch-filling ideas for SQL generation but replaces neural classification with deterministic pattern matching and ontology lookup, trading flexibility for full auditability.

Combinatory Categorial Grammar (CCG) parsing offers a theoretically cleaner approach: CCG assigns a syntactic category to each word that specifies both its syntactic role and its semantic contribution, allowing a fully compositional meaning representation to be constructed in parallel with the syntactic parse. CCG parsers have been applied to ATIS with strong results, but their construction requires a carefully curated CCG lexicon for each new domain — a labour-intensive process that limits practical deployment. Abstract Meaning Representation (AMR) takes a different path, encoding sentence meaning as a directed acyclic graph of concepts and roles that is more abstracted than a raw parse tree but less tied to specific database schemas than frame-based approaches. Neither CCG nor AMR has seen widespread adoption in Text-to-SQL, primarily because the mapping from AMR to executable \ac{SQL} requires an additional schema-grounding step that effectively recreates the problem it was meant to solve.

The sketch-and-fill paradigm initiated by SQLNet can be seen as a practical compromise between the full compositionality of CCG and the flat frame of task-oriented dialogue. By committing to a fixed sketch (SELECT [cols] FROM [tables] WHERE [conditions] GROUP BY [cols] ORDER BY [cols] LIMIT [n]) and learning to fill each slot independently, sketch-based models avoid the need for a formal semantic grammar while still constraining the output to a predictable structural form. The ontology-grounded approach in this project is a further specialisation of the sketch paradigm: the sketch structure is fixed in the \texttt{GROUP\_BY\_CONFIG} templates, and the slot-filling is performed by deterministic lookup rather than neural classification, making the process transparent and auditable at every step.

\section{Semantic Drift --- Prior Work and Definitions}
\label{sec:drift-prior-work}

The term ``semantic drift'' appears in the linguistics and cognitive science literature to describe the gradual shift in meaning that words undergo over time. In the context of information retrieval, drift has been used to describe query expansion gone wrong: a search that begins with ``apple pie recipe'' and, after successive expansion steps, retrieves results about Apple Inc. In the Text-to-SQL context, the term does not appear to have been formally defined in the literature reviewed for this project; the closest antecedent is the concept of \textit{faithfulness} from the natural language generation evaluation literature.

A related concept from the knowledge representation community is \textit{ontological drift} — the gradual divergence between a deployed ontology and the evolving domain it is intended to model. As business processes change (new product categories are introduced, tax rules are updated, regional boundaries shift), an ontology that is not actively maintained will accumulate a growing number of incorrect entity definitions and stale constraint rules. This form of drift is distinct from the query-level drift measured in this project, but the two are compounding: an outdated ontology leads to incorrect constraint evaluations, which in turn raise the constraint adherence component of the drift score, potentially causing converged queries to begin failing after an ontology update. Monitoring ontology staleness is therefore an important operational concern in any long-running deployment of a system of this kind.

The information retrieval literature also treats drift in the context of relevance feedback: a user refines a search query iteratively, and the system must not drift towards retrieving only documents similar to those already seen. The critic loop in this project faces an analogous risk — an iterated refinement process that converges to the same incorrect query on each pass, trapped in a local minimum of the drift function. In the evaluation, this was not observed: each iteration that changed the metric mapping also changed the SQL sufficiently to shift the plausibility score. However, it remains a theoretical failure mode that a more adversarial test set might expose.

Faithfulness metrics — used to evaluate whether a generated sentence accurately reflects a structured input — are the closest formal parallel. Work on data-to-text faithfulness by Maynez et al.\ (2020) proposed counting the fraction of facts in a generated text that are verifiable against the source. Transposing this to Text-to-SQL, one might count the fraction of intent elements (entities, metrics, filters) that are correctly reflected in the generated \ac{SQL}. This project adopts a similar decomposition, treating intent alignment, constraint adherence, and result plausibility as three independently measurable faithfulness dimensions.

From the database community, work on query repair and approximate query processing is also relevant. Systems such as QUOC and NaLIR detect syntactic errors in user-formulated \ac{SQL} and suggest corrections, but they operate on \ac{SQL} input rather than natural language. The gap this project fills is earlier in the pipeline: detecting when the natural language-to-SQL translation has introduced a semantic error that is invisible to the \ac{SQL} engine.

The distinction between syntactic and semantic correctness recurs throughout the literature and is central to this project. A syntactically correct \ac{SQL} query is one that the database engine can parse and execute without error; a semantically correct query is one that computes the answer the user intended. SQL engines are, by design, agnostic to semantic correctness: they execute whatever query they receive. Schema constraints and foreign key checks can catch certain classes of data integrity violations, but they cannot determine whether the query's GROUP BY clause correctly reflects the user's intended level of aggregation or whether the \texttt{WHERE} clause correctly represents the user's filter intent. These semantic properties live above the SQL engine's remit, in the space between the user's mental model and the query that was actually executed — precisely the space the drift metric is designed to measure.

The concept of ``query intent'' itself requires careful definition. In the context of this project, intent is operationalised as the tuple (entities, metrics, filters): the set of database tables the user wishes to aggregate over, the aggregation formula they want applied, and the filter conditions they want enforced. This operationalisation is admittedly a simplification — human query intent is richer, more context-dependent, and often includes unstated background assumptions. But as a practical approximation for the class of business intelligence queries in scope, it is sufficiently precise to yield a drift score that correlates with human judgements of query quality, as informally confirmed by the worked example in Section~\ref{sec:worked-example}.

Faithfulness evaluation in natural language generation has produced several reference-based metrics that are conceptually related to semantic drift. BARTScore measures faithfulness by computing the log-likelihood of a hypothesis given a reference using BART's encoder-decoder, yielding a continuous score that correlates well with human judgments. QuestEval uses a question-generation and question-answering approach: it generates questions from the reference, answers them using the hypothesis, and scores based on answer agreement. The PARENT metric, designed for table-to-text evaluation, computes precision and recall separately against both the table and the reference text. The common limitation of all three approaches is that they require a reference output for comparison — exactly the resource that is unavailable in production Text-to-SQL deployment. The drift metric in this project was designed specifically to avoid this requirement, computing all three component scores from the system's own intermediate representations.

\section{Multi-Agent Architectures for NLP Tasks}
\label{sec:multi-agent}

Decomposing a complex \ac{NLP} task into specialised agents, each responsible for a narrow sub-problem, has gained significant traction since the introduction of LangChain and similar orchestration frameworks. The motivation is modularity: each agent can be developed, tested, and replaced independently, and the overall system's behaviour is more transparent because each agent's input and output are explicit.

LangGraph \cite{langgraph} extends this idea by modelling agent interactions as a typed state graph. Nodes represent agents or processing steps; edges represent transitions that carry a shared state object. This makes the flow of information through the system auditable: at any point one can inspect the state to see exactly what each agent has contributed. The orchestrator in this project used LangGraph's \texttt{StateGraph} abstraction, with the five processing agents as nodes and the critic loop as a back-edge from the ResultVerifier to the IntentParser. A mock LangGraph module was also implemented to ensure the system operates without network access when the LangGraph package is unavailable.

LangGraph's typed state approach offers a concrete advantage over the older LangChain ``chain'' abstraction: in a LangChain sequential chain, each step is responsible for maintaining backward compatibility with all previous steps' outputs through a shared dictionary with string keys. In LangGraph's typed state, the state object has well-defined fields with explicit types, and the runtime validates that each node writes only to the fields it is declared to produce. This compile-time-like checking prevents a common class of inter-agent bugs where one agent's key name change silently breaks a downstream agent that expected the old name. The \texttt{QueryState} dataclass in this project provides similar guarantees without requiring the full LangGraph runtime: because it uses Python's \texttt{dataclasses} module with typed fields, any attempt to write to an undeclared attribute raises a runtime exception, and \texttt{mypy} or \texttt{pyright} type checkers would catch it statically.

Agent-based decomposition has been applied to various \ac{NLP} tasks: multi-hop question answering (where a retriever agent fetches relevant passages and a reader agent extracts answers), code generation (where a planner agent outlines steps and a coder agent implements them), and conversational systems (where a routing agent dispatches utterances to domain-specific handlers). The five-agent architecture in this project follows the same compositional principle but is specifically designed around the stages of structured query generation.

An important distinction between the five-agent architecture in this project and many LangChain-based multi-agent systems is the role of the shared knowledge base. In typical LangChain agent chains, each agent either uses its own tools (web search, code interpreter, etc.) or passes a text string to the next agent via the shared memory. There is no structured knowledge graph accessible to all agents simultaneously. In this project's architecture, the \texttt{BusinessOntology} object is instantiated once at system startup and passed by reference to every agent. This means that when the OntologyMapper resolves an entity, the ConstraintValidator can immediately access the constraint rules associated with that entity from the same object, without a serialisation or retrieval step. The performance benefit — no inter-agent I/O — is secondary to the consistency benefit: the ontology is a single source of truth that all agents read from simultaneously, so there is no possibility of the ConstraintValidator working with a different version of the constraint rules than the OntologyMapper used to resolve entities.

The ReAct (Reason + Act) prompting pattern, introduced by Yao et al., provides a conceptual predecessor to the critic loop in this project. ReAct alternates between reasoning traces (``Thought: I need to check whether the tax filter was applied'') and action calls (``Act: re-invoke the constraint validator''), producing auditable behaviour from an otherwise opaque \ac{LLM}. The critic loop in this project realises a similar alternation without relying on \ac{LLM} chain-of-thought: reasoning is performed by the deterministic drift calculator, and the action is the re-invocation of the pipeline with error feedback. Reflexion, a related approach by Shinn et al., adds an episodic memory of past failures that informs future attempts — a pattern this project approximates by appending the failing component's score to the \texttt{trace\_log} before each retry.

\section{Ontology-Based Query Grounding}
\label{sec:ontology-grounding}

Query grounding — mapping informal user expressions to formally defined schema elements — is a prerequisite for reliable \ac{SQL} generation. Ontologies encode domain knowledge as a network of entities, relationships, and axioms, providing a controlled vocabulary that bridges the gap between colloquial language and database schema terms.

In the semantic web community, ontologies expressed in OWL or RDFS are used to ground natural language queries against knowledge graphs. Systems such as QAKis and QALD use ontology-guided entity linking to resolve ambiguous surface forms to canonical URIs. The present project applies a lighter-weight approach: a Python \texttt{NetworkX} \cite{networkx} directed graph encoding eleven business entities as nodes and eight join relationships as labelled edges. This representation supports efficient shortest-path computation for multi-table join inference and simple graph traversal for constraint lookup, without the overhead of a full OWL reasoner.

The choice of a lightweight graph over a full OWL ontology was informed by practical deployment considerations. OWL reasoning requires a reasoner such as HermiT or Pellet, which adds a JVM dependency (these reasoners are Java-based) and inference times that can reach hundreds of milliseconds for complex ontologies. For the present application — resolving entity references and finding join paths in a bounded eleven-entity graph — Dijkstra's algorithm on a \texttt{NetworkX} \texttt{DiGraph} completes in under 1~ms for all test queries. The expressivity sacrifice (no class hierarchies, no property restrictions, no cardinality constraints expressible in OWL terms) is acceptable because the relevant business rules are encoded as first-class \texttt{ConstraintRule} objects rather than being derived by reasoning.

One area where a richer ontology representation would genuinely help is entity disambiguation. When the user mentions ``products,'' the ontology currently returns PRODUCT unconditionally. With an OWL subclass hierarchy, PRODUCT could be a parent class of ELECTRONIC\_PRODUCT, CLOTHING\_PRODUCT, and so on, allowing the intent parser to resolve to a more specific entity when the query includes a category filter. The current system handles this through the product category filter detector, which adds a \texttt{WHERE p.category = 'Electronics'} clause, but the resolution is at the SQL level rather than the ontology level.

The business ontology also encodes twelve aggregate metrics, each with a formal SQL expression template and a list of entity dependencies. When the intent parser identifies a metric such as ``gross margin,'' the ontology mapper can immediately retrieve both the formula (revenue minus tax, as a percentage of revenue) and the table joins required to compute it. This separation of metric definitions from the SQL builder prevents the kind of ad hoc formula duplication that is a common source of inconsistency in business intelligence environments.

Metric definitions in the ontology also carry a \texttt{requires\_entities} list, which specifies which entities must be included in the query for the metric to be computable. The GROSS\_MARGIN metric, for example, lists both ORDER and ORDER\_ITEM in its \texttt{requires\_entities} field, because its formula uses \texttt{o.total\_amount} (from orders) and the join to order\_items is needed to access line-item cost data. When the ExecutionPlanner resolves a metric, it checks this list against the already-resolved entity set and adds any missing entities to \texttt{joins\_needed}. This mechanism prevents the common SQL error of referencing columns from tables that are not in the FROM clause, which would produce a runtime error rather than an incorrect result.

Schema linking — the sub-task of mapping natural language tokens to specific schema elements (tables, columns, values) — is the core problem that ontology grounding addresses. RAT-SQL's graph attention mechanism constructs edges between question words and schema elements, allowing the transformer to learn alignment patterns from data. The approach in this project is lexical rather than learned: each ontology entity carries an explicit list of aliases (e.g., CUSTOMER has aliases ``client'', ``buyer'', ``user''), and the mapper uses \texttt{SequenceMatcher} to score candidates. While the lexical approach is less flexible on novel phrasings, it is fully transparent: every mapping decision can be explained in terms of the similarity score between the user's token and a specific alias string, with no need to interrogate a neural network's attention weights.

\section{Anomaly Detection in Query Results}
\label{sec:anomaly-detection}

Statistical anomaly detection has a long history in quality control and fraud detection, dating to Walter Shewhart's control chart work in the 1920s. The Z-score, defined as the signed number of standard deviations by which a value departs from the mean of a reference distribution, is among the simplest and most widely deployed anomaly statistics. Values with $|z| > 2$ are typically flagged for review; those with $|z| > 3$ are treated as strong anomalies in most industrial applications.

The Z-score's suitability as a plausibility signal in this context depends on the assumption that historical query results are approximately normally distributed around a stable mean — a reasonable assumption for well-formed aggregation queries on a static dataset, but potentially violated by seasonal patterns, business growth trends, or schema changes. In a production environment, a more robust approach would use a rolling baseline that downweights older data points and accounts for trend, similar to the Holt-Winters exponential smoothing method used in time-series forecasting. For the purposes of this project, with a static synthetic dataset and no temporal dynamics, a simple mean and standard deviation over the available history sufficed.

Applying Z-score anomaly detection to query results is relatively uncommon in the existing literature. Closest in spirit is the work on data quality monitoring in data warehouses, where SQL assertions and freshness checks are used to catch load failures. The innovation in this project is applying plausibility scoring at query time rather than ingestion time, and incorporating the resulting plausibility score into a composite drift metric that drives query revision. The approach acknowledges that historical baselines may themselves be imperfect, so the plausibility component is weighted at 30\% rather than allowed to dominate the overall score.

Row-count sanity checking provides a complementary signal: if a query that historically returns 12 rows suddenly returns 12,000 rows, a join is likely misconfigured. The ResultVerifier agent compares the current result count against a ten-times multiplier threshold; if the count exceeds this threshold and the aggregate value is simultaneously anomalous, the combined plausibility score is reduced sharply.

More sophisticated anomaly detection approaches, including Isolation Forest and Local Outlier Factor (LOF), were considered as alternatives to Z-score for the plausibility component. Isolation Forest partitions data by randomly selecting a feature and a split point; anomalies are points that require fewer partitions to isolate. LOF computes the density ratio between a point and its k-nearest neighbours, identifying outliers as points in low-density regions surrounded by high-density clusters. Both methods are more robust to non-Gaussian distributions than Z-score. However, both require a training phase and at least tens of data points to produce meaningful scores. In the context of query result verification, where baselines may consist of only a handful of historical query executions, Z-score was chosen as the more appropriate choice: it is interpretable, computationally trivial (under 1 ms), and produces a well-defined result even with as few as two historical values.

\section{Gaps in Existing Literature}
\label{sec:literature-gaps}

Three gaps in the existing work motivated the specific design choices of this project:

\textbf{Gap 1 — No formal drift metric.} Existing Text-to-SQL evaluation frameworks measure final-query accuracy (exact match or execution match against a ground-truth \ac{SQL}). This presupposes the existence of a reference query, which is unavailable in production. This project defines drift as a query-time, reference-free score computable from the system's own intermediate representations.

\textbf{Gap 2 — No critic loop for silent failures.} Systems that produce a wrong-but-plausible result have no mechanism to self-correct. The critic loop in this project treats the drift score as a continuous signal and re-invokes the pipeline with error context, enabling iterative refinement without user intervention.

\textbf{Gap 3 — No ontology-grounded constraint enforcement.} Most Text-to-SQL systems treat schema column names as the sole grounding target. Business constraint rules --- such as ``always exclude tax from revenue calculations'' or ``never expose raw \ac{PII} columns'' --- exist in policy documents, not in the schema. The ontology in this project encodes nine such rules, making constraint adherence measurable and enforceable at query time.

These three gaps collectively define the contribution space of this project. They are orthogonal to the accuracy improvements pursued by the mainstream Text-to-SQL benchmarking community: the project does not claim state-of-the-art accuracy on Spider or WikiSQL, nor does it introduce a novel neural architecture. Rather, it addresses the deployment reliability dimension — the question of what happens after a model achieves high benchmark accuracy and is placed in production where no reference queries exist and where constraint violations may have real business consequences.

A fourth gap, less central but worth acknowledging, is the absence of a standardised API contract for Text-to-SQL systems. Most research implementations expose either a command-line interface or a single-function Python API; there is no widely adopted REST API design for natural language query services that includes quality metadata (drift score, iteration count, execution trace) alongside the result. The eight-endpoint API designed for this project is one concrete proposal for what such a contract might look like, though it is not presented as a standards contribution.

Finally, a fifth gap concerns the lack of production-grade synthetic evaluation datasets for domain-specific Text-to-SQL. The Spider and WikiSQL benchmarks use real database schemas from diverse domains, but neither includes the kind of domain-specific business constraint rules (``always net out tax'', ``never return PII columns'') that distinguish a business intelligence system from a general-purpose database interface. The synthetic e-commerce dataset and constraint rule set developed for this project are one step toward filling this gap for the retail and e-commerce domain, though a much larger and more diverse dataset would be needed to serve as a community benchmark.

\section{Comparison of Related Systems}
\label{sec:related-comparison}

Table~\ref{tab:related-work} summarises the closest related Text-to-SQL systems along the dimensions most relevant to this project.

\begin{table}[H]
  \centering
  \caption{Comparison of Related Text-to-SQL Systems}
  \label{tab:related-work}
  \renewcommand{\arraystretch}{1.3}
  \begin{tabularx}{\textwidth}{@{}lXXXX@{}}
    \toprule
    \textbf{System} & \textbf{Approach} & \textbf{Benchmark} & \textbf{Accuracy} & \textbf{Limitation} \\
    \midrule
    Seq2SQL \cite{seq2sql} & Seq2Seq + RL & WikiSQL & 59.4\% exec. & Single-table, no joins \\
    SQLNet \cite{sqlnet}   & Sketch-based & WikiSQL & 68.0\% exec. & Template-constrained \\
    IRNet   & BERT + sketch & Spider  & 53.2\% exact & No constraint rules \\
    RAT-SQL & Graph + BERT  & Spider  & 69.7\% exact & No plausibility check \\
    PICARD  & Constrained decoding & Spider & 75.1\% exact & No feedback loop \\
    This work & Ontology + 5-agent & Synthetic e-comm. & 90\% success & Domain-specific ontology \\
    \bottomrule
  \end{tabularx}
\end{table}

% =============================================================================
% CHAPTER 3 — SYSTEM DESIGN AND ARCHITECTURE
% =============================================================================
\chapter{System Design and Architecture}
\label{ch:design}

\section{System Overview and Design Philosophy}
\label{sec:design-overview}

The system was designed around three guiding principles. First, \textit{auditability}: every decision made during query processing should be traceable to a specific agent, a specific rule, or a specific measurement, not buried inside a neural network's weight matrix. Second, \textit{graceful degradation}: when the \ac{LLM} \ac{API} is unavailable, the system should fall back to deterministic pattern matching and continue operating, perhaps with lower coverage, rather than failing entirely. Third, \textit{measurability}: the system should produce a numeric estimate of its own reliability for every query, enabling downstream consumers to apply appropriate scepticism.

The auditability principle has a practical implication for the choice between pattern-matching and neural NLU. A neural model that achieves 90\% accuracy on a test set provides no explanation for the 10\% it gets wrong. When a business user asks ``why was this query wrong?'', the answer ``the model assigned the wrong token embedding to `region''' is not actionable. By contrast, when the pattern-based system gets a query wrong, the trace log provides a specific, actionable answer: ``the \texttt{\_detect\_group\_by} detector did not match `across regions' because the preposition `across' is not in the pattern set''. This is a fixable, understandable failure. The architectural choice of pattern matching over neural NLU was made with full awareness that it limits coverage; the trade-off was accepted because auditability and debuggability were judged more important than marginal coverage gains at this stage of the system's development.

The measurability principle distinguishes this system from a conventional query interface in a way that may not be immediately obvious. Most database interfaces return a result or an error — a binary outcome. The drift score introduces a third category: ``returned a result, but with measured uncertainty''. This third category is valuable precisely because it is honest about the system's limitations. A user who receives a result with $D = 0.14$ (just below threshold) should understand that the result is marginally accepted and may warrant manual verification. A user who receives a result with $D = 0.22$ (above threshold) should understand that the system was unable to converge, and the returned result is the best available approximation, not a confident answer. Communicating this uncertainty, rather than suppressing it, is the defining characteristic of a trustworthy data system.

These principles led to a layered architecture in which a pattern-matching-first agent pipeline sits below an optional \ac{LLM} enhancement layer, and a quantitative drift metric sits above both as an independent quality signal. The ontology occupies the centre of the design, acting as the shared knowledge base that all five agents consult for definitions, rules, and join paths.

The phrase ``shared knowledge base'' is important here. In a monolithic Text-to-SQL function, domain knowledge is implicitly distributed across the model's weights, the hard-coded prompt, and the schema serialisation passed to the model. When a query fails, there is no clean way to attribute the failure to a specific piece of knowledge because the knowledge is not separately addressable. The ontology in this project makes domain knowledge first-class and addressable: a constraint rule failure points to a named rule (\texttt{R01\_TAX\_EXCLUSION}); a mapping failure points to a named entity (\texttt{ORDER}) and the alias that was not found. Every piece of domain knowledge has a name, a location (the ontology Python object), and a clear owner (the OntologyMapperAgent or ConstraintValidatorAgent). This addressability is what enables the feedback message to be specific rather than generic, and specific feedback is what makes the critic loop effective.

The design philosophy also prioritised testability as a first-class property. Because each agent has a narrow, single-responsibility interface that accepts a \texttt{QueryState} and returns a modified \texttt{QueryState}, every agent can be unit-tested in complete isolation. The IntentParser's pattern matching can be tested with a plain Python string; the OntologyMapper's alias resolution can be tested with a mocked ontology object; the ConstraintValidator's rule evaluation can be tested with a hand-crafted constraint list. This isolation means that adding a new detector or alias does not require running the full pipeline to verify correctness — a significant development velocity advantage compared to end-to-end testing.

A deliberate decision was made to avoid the common pattern of a single monolithic ``text-to-SQL'' function that accepts a natural language string and returns a \ac{SQL} string. While such a function is simpler to implement and test in isolation, it provides no visibility into which stage of the translation introduced an error. The decomposition into five agents means that the source of elevated drift — intent misextraction, mapping failure, constraint violation, SQL construction error, or result anomaly — can be identified from the drift component scores alone, without re-running the query or inspecting log files. This auditability was considered essential for a production context where a business analyst would need to understand why a query was flagged rather than just being told that it was.

The layered design also enables progressive enhancement without architectural refactoring. The deterministic pattern-matching layer provides a functional baseline; the optional Claude integration layer adds neural query generation on top without changing the scoring or feedback infrastructure. In principle, additional enhancement layers — a RAG-based schema retriever, a SQL validator, a query cache — could be inserted between existing agents without changing the agent contracts.

\section{High-Level Architecture}
\label{sec:architecture}

Figure~\ref{fig:architecture} shows the high-level component layout.

\begin{figure}[H]
  \centering
  \fbox{\parbox[c][7cm][c]{0.85\linewidth}{\centering
    \textbf{High-Level System Architecture}\\[0.5em]
    \texttt{[User Query]} $\rightarrow$ \texttt{[FastAPI REST Layer]} $\rightarrow$ \texttt{[Orchestrator]}\\[0.3em]
    \texttt{[Orchestrator]} $\leftrightarrow$ \texttt{[Five-Agent Pipeline]}\\[0.3em]
    \texttt{[Pipeline]} $\leftrightarrow$ \texttt{[Business Ontology (NetworkX)]}\\[0.3em]
    \texttt{[Pipeline]} $\rightarrow$ \texttt{[Database Layer (SQLite/PostgreSQL)]}\\[0.3em]
    \texttt{[Pipeline]} $\rightarrow$ \texttt{[Semantic Drift Calculator]}\\[0.3em]
    \texttt{[Drift $<$ 0.15?]} $\rightarrow$ \texttt{[Result to User]} OR \texttt{[Iterate]}\\[0.3em]
    \texttt{[Streamlit Dashboard]} $\leftrightarrow$ \texttt{[FastAPI REST Layer]}
  }}
  \caption{High-level system architecture showing the query flow from user input through the agent pipeline to result delivery.}
  \label{fig:architecture}
\end{figure}

The user submits a natural language question through either the \ac{REST} \ac{API} or the Streamlit dashboard. The orchestrator receives the query, initialises a \texttt{QueryState} object, and passes it to the five-agent pipeline. After all five agents have processed the state, the semantic drift calculator scores the result. If the score is below the convergence threshold, the result is returned; otherwise the orchestrator increments the iteration counter, appends error context to the state, and re-invokes the pipeline. This critic loop runs for at most five iterations.

The component boundaries in the architecture were drawn to minimise coupling between agents. Each agent reads only from the \texttt{QueryState} fields that its immediate predecessor populates and writes only to its own output fields. The IntentParser writes to \texttt{state.intent}; the OntologyMapper reads \texttt{state.intent} and writes to \texttt{state.mappings}; the ConstraintValidator reads both and writes to \texttt{state.constraints}; and so on. No agent accesses the output of a non-adjacent predecessor directly. This strict ordering makes the pipeline testable: any agent can be unit-tested by constructing a \texttt{QueryState} with only its required input fields populated, without needing to run the full pipeline. It also makes the pipeline extensible: a new agent can be inserted between any two existing agents without modifying either of them, provided it reads only from its predecessor's output fields and writes only to its own.

One architectural decision worth documenting explicitly is the choice to run all five agents sequentially rather than in parallel. A parallel execution model — running the OntologyMapper and IntentParser in parallel on the initial query, then merging results — could in principle reduce latency. The sequential model was chosen because each agent's output informs the next agent's configuration: the OntologyMapper cannot map entities without the IntentParser's extracted tokens; the ConstraintValidator cannot evaluate rules without the OntologyMapper's resolved entity set. The data dependencies are strictly linear, making parallelisation impossible without significant architectural changes that would undermine the clean agent-boundary design. For the target latency of under 5 seconds, the sequential model comfortably meets the requirement with room to spare, so the added complexity of a parallel model was not justified.

\section{Business Ontology Design}
\label{sec:ontology-design}

\subsection{Entity Definitions}
\label{subsec:entities}

The ontology encodes eleven business entities drawn from a representative e-commerce schema. Each entity maps to one or more database tables and carries metadata about its primary key, join relationships, and applicable constraints. Table~\ref{tab:entities} lists the core entities.

\begin{table}[H]
  \centering
  \caption{Business Entity Definitions in the Ontology}
  \label{tab:entities}
  \renewcommand{\arraystretch}{1.2}
  \begin{tabularx}{\textwidth}{@{}llllX@{}}
    \toprule
    \textbf{Entity ID} & \textbf{Table} & \textbf{PK} & \textbf{Aliases} & \textbf{Description} \\
    \midrule
    CUSTOMER    & customers    & customer\_id & client, buyer, user & End consumer \\
    PRODUCT     & products     & product\_id  & item, sku, goods  & Catalogue item \\
    ORDER       & orders       & order\_id    & purchase, sale    & Sales transaction header \\
    INVOICE     & invoices     & invoice\_id  & bill, receipt     & Financial invoice \\
    TRANSACTION & transactions & transaction\_id & payment, tx    & Payment record \\
    SUPPLIER    & suppliers    & supplier\_id & vendor, partner   & Product supplier \\
    ORDER\_ITEM & order\_items & item\_id     & line item, detail & Order line detail \\
    RETURN      & returns      & return\_id   & refund, reversal  & Return record \\
    EMPLOYEE    & employees    & employee\_id & staff, worker     & HR record \\
    DEPARTMENT  & departments  & dept\_id     & team, unit        & Organisational unit \\
    FINANCE     & accounts     & account\_id  & ledger, budget    & Finance account \\
    \bottomrule
  \end{tabularx}
\end{table}

\subsection{Metric Definitions}
\label{subsec:metrics}

Twelve aggregate metrics were defined, each with a formal \ac{SQL} expression template stored in the \texttt{METRIC\_SELECT} dictionary at lines~39--52 of \texttt{src/agents/pipeline.py}. Table~\ref{tab:metrics} summarises the metrics and their aggregation logic.

\begin{table}[H]
  \centering
  \caption{Business Metric Definitions}
  \label{tab:metrics}
  \renewcommand{\arraystretch}{1.2}
  \begin{tabularx}{\textwidth}{@{}llXl@{}}
    \toprule
    \textbf{Metric ID} & \textbf{Alias} & \textbf{SQL Expression} & \textbf{Unit} \\
    \midrule
    REVENUE      & total\_revenue  & \texttt{SUM(o.total\_amount)}              & currency \\
    COUNT        & total\_orders   & \texttt{COUNT(*)}                         & integer \\
    AVERAGE      & avg\_order      & \texttt{AVG(o.total\_amount)}             & currency \\
    PROFIT       & total\_profit   & \texttt{SUM(amount - tax - shipping)}     & currency \\
    GROSS\_MARGIN & margin\_pct    & \texttt{ROUND(SUM(amount-tax)*100/...)}   & percent \\
    NEW\_CUSTOMERS & unique\_cust  & \texttt{COUNT(DISTINCT customer\_id)}     & integer \\
    CHURN        & cancelled      & \texttt{SUM(CASE WHEN status='Cancelled')} & integer \\
    LTV          & avg\_ltv       & \texttt{SUM(amount)/COUNT(DISTINCT cust)} & currency \\
    RETURN\_RATE & cancel\_pct    & \texttt{SUM(cancel)*100/COUNT(*)}         & percent \\
    INVENTORY    & total\_stock   & \texttt{SUM(p.stock\_quantity)}           & units \\
    CONVERSION   & converting     & \texttt{COUNT(DISTINCT customer\_id)}     & integer \\
    RETENTION    & retained       & \texttt{COUNT(DISTINCT customer\_id)}     & integer \\
    \bottomrule
  \end{tabularx}
\end{table}

\subsection{Constraint Rules}
\label{subsec:constraints}

Nine constraint rules govern what queries the system may generate. These rules are checked by the ConstraintValidatorAgent and scored by the ConstraintAdherenceCalculator. Table~\ref{tab:constraints} provides the rule catalogue.

\begin{table}[H]
  \centering
  \caption{Ontology Constraint Rules}
  \label{tab:constraints}
  \renewcommand{\arraystretch}{1.2}
  \begin{tabularx}{\textwidth}{@{}llXl@{}}
    \toprule
    \textbf{Rule ID} & \textbf{Type} & \textbf{Description} & \textbf{Severity} \\
    \midrule
    R01 & TAX\_EXCLUSION        & Exclude tax\_amount when reporting net revenue  & HIGH \\
    R02 & REGION\_FILTER        & Validate region values against allowed set       & MEDIUM \\
    R03 & DATE\_RANGE           & Enforce sensible date window (not beyond 2 years) & MEDIUM \\
    R04 & ACTIVE\_PRODUCTS\_ONLY & Filter to \texttt{is\_active=1} products          & HIGH \\
    R05 & DATA\_FRESHNESS       & Flag queries on data older than 90 days           & LOW \\
    R06 & FISCAL\_YEAR          & Align fiscal year boundaries to April--March       & MEDIUM \\
    R07 & PII\_PROTECTION       & Never expose raw email, phone, or address columns  & HIGH \\
    R08 & ACTIVE\_CUSTOMERS     & Restrict to \texttt{is\_active=1} customers        & MEDIUM \\
    R09 & MIN\_ORDER\_VALUE     & Exclude zero-value orders (\texttt{total\_amount > 0}) & LOW \\
    \bottomrule
  \end{tabularx}
\end{table}

\subsection{Join Cardinality Rules and NetworkX Graph}
\label{subsec:joins}

The ontology graph uses a \texttt{NetworkX} \texttt{DiGraph} with eight explicit join edges. Nodes carry entity metadata; edges carry join condition, cardinality (one-to-many or many-to-many), and a weight used by the shortest-path algorithm. When the ExecutionPlanner needs to join CUSTOMER to PRODUCT, it calls \texttt{get\_join\_path("CUSTOMER", "PRODUCT")}, which returns the path \texttt{CUSTOMER $\to$ ORDER $\to$ ORDER\_ITEM $\to$ PRODUCT} — a three-hop join that the SQL builder then realises as \texttt{LEFT JOIN} clauses.

Table~\ref{tab:join-rules} documents the eight join edges in the ontology graph.

\begin{table}[H]
  \centering
  \caption{Ontology Join Rules (NetworkX DiGraph Edges)}
  \label{tab:join-rules}
  \renewcommand{\arraystretch}{1.2}
  \begin{tabularx}{\textwidth}{@{}llllX@{}}
    \toprule
    \textbf{From Entity} & \textbf{To Entity} & \textbf{Type} & \textbf{Cardinality} & \textbf{JOIN Condition} \\
    \midrule
    CUSTOMER    & ORDER       & LEFT  & 1:N & \texttt{o.customer\_id = c.customer\_id} \\
    ORDER       & ORDER\_ITEM & LEFT  & 1:N & \texttt{oi.order\_id = o.order\_id} \\
    ORDER\_ITEM & PRODUCT     & LEFT  & N:1 & \texttt{p.product\_id = oi.product\_id} \\
    PRODUCT     & SUPPLIER    & LEFT  & N:1 & \texttt{s.supplier\_id = p.supplier\_id} \\
    ORDER       & INVOICE     & LEFT  & 1:1 & \texttt{inv.order\_id = o.order\_id} \\
    ORDER       & TRANSACTION & LEFT  & 1:N & \texttt{t.order\_id = o.order\_id} \\
    ORDER       & RETURN      & LEFT  & 1:N & \texttt{r.order\_id = o.order\_id} \\
    EMPLOYEE    & DEPARTMENT  & LEFT  & N:1 & \texttt{d.dept\_id = e.dept\_id} \\
    \bottomrule
  \end{tabularx}
\end{table}

The weight attribute on each edge is uniform (1.0) by default, meaning \texttt{nx.shortest\_path} applies Dijkstra's algorithm and returns the minimum-hop path. A future extension could assign higher weights to cross-domain joins (e.g., EMPLOYEE to ORDER) to discourage the system from joining HR and sales tables unless the user's question explicitly requires it. The directed nature of the graph means that join inference is order-sensitive: a path from ORDER to CUSTOMER (traversing the edge in reverse) is not available without adding a reverse edge. In the current ontology, all eight join edges are directed from the ``many'' side to the ``one'' side (child to parent), so join path inference always follows foreign keys in the conventional direction. This convention prevents the system from generating inadvertent Cartesian products that would result from joining in the wrong direction.

Figure~\ref{fig:ontology-graph} provides a schematic view of the entity-join subgraph.

\begin{figure}[H]
  \centering
  \fbox{\parbox[c][7cm][c]{0.85\linewidth}{\centering
    \textbf{Business Ontology Entity--Join Graph (NetworkX DiGraph)}\\[0.5em]
    \textit{Node layout (approximate):}\\[0.3em]
    \texttt{CUSTOMER}~$\longrightarrow$~\texttt{ORDER}~$\longrightarrow$~\texttt{ORDER\_ITEM}~$\longrightarrow$~\texttt{PRODUCT}\\
    \hspace{5.2cm}$\downarrow$ \hspace{2.8cm} $\downarrow$\\
    \hspace{4.2cm}\texttt{INVOICE}\hspace{1.1cm}\texttt{SUPPLIER}\\
    \hspace{5.2cm}$\downarrow$\\
    \hspace{4.5cm}\texttt{TRANSACTION}\\
    \hspace{5.2cm}$\downarrow$\\
    \hspace{4.8cm}\texttt{RETURN}\\[0.5em]
    \texttt{EMPLOYEE}~$\longrightarrow$~\texttt{DEPARTMENT}\\[0.3em]
    \textit{(Edge labels carry join condition and cardinality.)}\\
    \textit{Metric nodes omitted for clarity.}
  }}
  \caption{Schematic of the business ontology entity--join graph stored as a \texttt{NetworkX} \texttt{DiGraph}. Directed edges represent foreign-key join relationships; shortest-path traversal from any source entity to any target entity yields the sequence of \texttt{LEFT JOIN} clauses needed by the SQL builder.}
  \label{fig:ontology-graph}
\end{figure}

\section{Five-Agent Pipeline Design}
\label{sec:pipeline-design}

Figure~\ref{fig:pipeline} shows the data-flow through the five agents.

\begin{figure}[H]
  \centering
  \fbox{\parbox[c][5cm][c]{0.85\linewidth}{\centering
    \textbf{Five-Agent Pipeline Data Flow}\\[0.5em]
    \texttt{NL Query}\\
    $\downarrow$\\
    \texttt{[1] IntentParserAgent} $\rightarrow$ ExtractedIntent\\
    $\downarrow$\\
    \texttt{[2] OntologyMapperAgent} $\rightarrow$ OntologyMapping\\
    $\downarrow$\\
    \texttt{[3] ConstraintValidatorAgent} $\rightarrow$ ValidatedConstraints\\
    $\downarrow$\\
    \texttt{[4] ExecutionPlannerAgent} $\rightarrow$ ExecutionPlan (SQL)\\
    $\downarrow$\\
    \texttt{[5] ResultVerifierAgent} $\rightarrow$ QueryResults + DriftScore
  }}
  \caption{Sequential data flow through the five agents. Each agent reads from and writes to the shared \texttt{QueryState} object.}
  \label{fig:pipeline}
\end{figure}

\subsection{IntentParserAgent}
\label{subsec:intent-parser}

The IntentParserAgent decomposes the raw natural language query into a structured \texttt{ExtractedIntent} object containing entity references, metric references, filter values, and a confidence score for each extraction. Eight entity patterns and twelve metric patterns are encoded as compiled regular expressions. In addition, fifteen filter detectors handle temporal, geographic, categorical, and quantitative qualifiers:

\begin{itemize}[leftmargin=2em]
  \item \texttt{\_detect\_top\_n} — extracts ``top 5'', ``bottom 10'' with direction
  \item \texttt{\_detect\_period\_type} — classifies ``last month'', ``YTD'', ``last quarter''
  \item \texttt{\_detect\_compare\_period} — flags year-over-year comparison intent
  \item \texttt{\_detect\_group\_by} — infers the grouping dimension from prepositions (``by region'', ``per category'', ``across tiers'')
  \item \texttt{\_detect\_having} — extracts threshold conditions (``greater than \$10,000'')
  \item \texttt{\_detect\_order\_status} — maps ``shipped'', ``cancelled'', ``delivered''
  \item \texttt{\_detect\_customer\_tier} — maps ``gold customers'', ``enterprise accounts''
  \item \texttt{\_detect\_product\_category} — maps ``electronics'', ``clothing'', ``furniture''
  \item \texttt{\_detect\_date\_range} — parses explicit year values such as ``in 2024''
  \item \texttt{\_detect\_moving\_avg} — flags rolling average requests
\end{itemize}

Each detection function returns a tuple of (value, confidence). The overall extraction confidence is the mean of all non-null confidence values, scaled to $[0, 1]$. A query that clearly maps to a single entity and metric achieves a confidence near 0.9; an ambiguous or multi-intent query typically scores 0.55--0.70.

Table~\ref{tab:detector-examples} illustrates six representative detectors with example inputs.

\begin{table}[H]
  \centering
  \caption{Sample IntentParser Detectors with Example Inputs}
  \label{tab:detector-examples}
  \renewcommand{\arraystretch}{1.2}
  \begin{tabularx}{\textwidth}{@{}lXlr@{}}
    \toprule
    \textbf{Detector} & \textbf{Example Input Fragment} & \textbf{Extracted Value} & \textbf{Conf.} \\
    \midrule
    \texttt{\_detect\_top\_n}         & ``top 5 products''      & (top, 5)           & 1.00 \\
    \texttt{\_detect\_period\_type}   & ``last quarter''        & quarter             & 0.90 \\
    \texttt{\_detect\_compare\_period} & ``year over year''     & yoy=True            & 0.95 \\
    \texttt{\_detect\_group\_by}      & ``by region''           & region              & 0.90 \\
    \texttt{\_detect\_having}         & ``greater than \$10,000'' & (gt, 10000)       & 0.85 \\
    \texttt{\_detect\_customer\_tier} & ``gold customers''      & gold                & 0.88 \\
    \bottomrule
  \end{tabularx}
\end{table}

\subsection{OntologyMapperAgent}
\label{subsec:ontology-mapper}

The OntologyMapperAgent takes the \texttt{ExtractedIntent} and resolves each mentioned entity and metric against the ontology graph using Python's \texttt{SequenceMatcher}. An exact match on entity\_id or alias returns a similarity score of 1.0; partial matches return proportional scores. If the best match falls below a configurable threshold (default 0.5), the element is marked unresolved and the subsequent constraint validation is adjusted accordingly.

The mapping process operates in two passes per token. In the first pass, the token is tested against every entity's \texttt{entity\_id} and every alias string using \texttt{SequenceMatcher(None, token.lower(), candidate.lower()).ratio()}. The candidate with the highest ratio wins; if the winning ratio is below the unresolved threshold, the second pass is triggered. The second pass applies the same logic to the metric dictionary. This two-pass approach means that a word like ``revenue,'' which is both an alias for the ORDER entity and a metric in the REVENUE metric definition, will be resolved to the metric if no entity alias produces a higher similarity score. In practice, metric tokens tend to produce higher similarities against metric aliases (because metric names are longer and more specific) while entity tokens produce higher similarities against entity aliases (because entity aliases include short colloquial words that match directly).

The mapper also infers join requirements from the resolved entity set. If the resolved entities span more than one node in the ontology graph (e.g., both CUSTOMER and PRODUCT are referenced), the mapper calls \texttt{get\_join\_path()} for each required entity pair and records the join path in the mapping output. The ExecutionPlanner then uses these pre-computed paths to populate \texttt{joins\_needed}, avoiding the need to re-query the graph during SQL construction.

One subtlety in the mapping implementation is the handling of the case where no entity pattern matches in the IntentParser. In this case, the \texttt{state.intent.entities} list is empty, and the OntologyMapper has nothing to resolve. Rather than producing an empty \texttt{state.mappings} list — which would cause the downstream ConstraintValidator to skip constraint evaluation (no entities = no constraints) and the drift metric to record a high constraint adherence score by default — the mapper falls back to a default entity resolution: it maps the query to the ORDER entity with a similarity score of 0.5, treating ``orders'' as the default business context. This conservative fallback ensures that common constraints (TAX\_EXCLUSION, MIN\_ORDER\_VALUE) are always evaluated, preventing unresolved queries from appearing to have high constraint adherence when they actually have undefined entity context.

Resolved mappings are stored in an \texttt{OntologyMapping} object attached to the \texttt{QueryState}. The mapping carries the canonical entity\_id, the matched table name, and the computed similarity score, which feeds directly into the intent alignment component of the drift metric.

The \texttt{OntologyMapping} object also carries an \texttt{is\_exact\_match} boolean that is set to \texttt{True} only when the similarity score is 1.0 — that is, the user's token matches an entity\_id or alias character-for-character. Queries where all resolutions are exact matches consistently achieve higher intent alignment scores (mean $a_i = 0.91$) than those with at least one partial match (mean $a_i = 0.79$), because partial matches contribute proportionally lower values to $\overline{S_m}$. This distinction is reflected in the feedback message generated after a failed iteration: when all mappings are exact but alignment is still below threshold, the feedback message attributes the issue to extraction confidence and directs the parser to re-examine ambiguous filter terms; when any mapping is partial, the message names the specific unresolved token and suggests it may require additional aliases in the ontology.

\subsection{ConstraintValidatorAgent}
\label{subsec:constraint-validator}

For each resolved entity, the ConstraintValidatorAgent retrieves applicable constraint rules from the ontology using \texttt{get\_applicable\_constraints(entity\_ids)}. It then evaluates each rule deterministically using the logic in the \texttt{\_check\_rule} function (lines~159--180 of \texttt{pipeline.py}):

\begin{itemize}[leftmargin=2em]
  \item \textbf{TAX\_EXCLUSION}: satisfied when the query includes the tax flag or the metric is net-revenue aware
  \item \textbf{REGION\_FILTER}: satisfied when the extracted region value is a member of \texttt{\_VALID\_REGIONS}
  \item \textbf{PII\_PROTECTION}: satisfied unconditionally because the pipeline never selects raw \ac{PII} columns
  \item \textbf{ACTIVE\_PRODUCTS\_ONLY, MIN\_ORDER\_VALUE}: satisfied when the \ac{SQL} builder enforces the relevant \texttt{WHERE} clause, which it does by default
\end{itemize}

The overall constraint adherence score $a_c$ (Equation~\ref{eq:constraint}) is computed from the ratio of satisfied to total applicable constraints.

The severity field on each constraint rule (HIGH, MEDIUM, LOW) is not yet used in the adherence scoring: all satisfied/unsatisfied outcomes are weighted equally in the ratio. A natural extension would be to apply a weighted adherence formula that penalises HIGH-severity violations more than LOW-severity ones, for example $a_c = 1 - \sum_{j \in \text{violated}} w_j / \sum_{j \in \text{all}} w_j$ with $w_j \in \{1.0, 0.5, 0.2\}$ for HIGH, MEDIUM, LOW respectively. This was not implemented because the test query set did not produce any HIGH-severity violations that would have distinguished the weighted from the unweighted formula; calibrating the severity weights would have required additional evaluation data.

\subsection{ExecutionPlannerAgent}
\label{subsec:execution-planner}

The ExecutionPlannerAgent calls \texttt{\_build\_sql(intent, dialect)} at line~183 of \texttt{pipeline.py} to generate the \ac{SQL} string. The function handles nine GROUP BY configurations (region, tier, category, month, year, status, product, region-and-tier, supplier), year-over-year comparison queries, HAVING clauses, TOP-N ordering, explicit date ranges, and multi-table joins. The dialect parameter selects between SQLite and PostgreSQL date functions and type casts. When the \texttt{ANTHROPIC\_API\_KEY} environment variable is set, the \texttt{ClaudeExecutionPlanner} subclass replaces pattern-generated \ac{SQL} with a Claude-generated query, using the ontology context as a system prompt.

Figure~\ref{fig:sql-builder} illustrates the decision flow within \texttt{\_build\_sql()}.

\begin{figure}[H]
  \centering
  \fbox{\parbox[c][7cm][c]{0.85\linewidth}{\centering
    \textbf{\_build\_sql() Decision Flow}\\[0.5em]
    \texttt{\_build\_sql(intent, dialect)}\\
    $\downarrow$\\
    \textit{YoY flag set?} YES $\rightarrow$ return CTE-style YoY query\\
    \quad NO $\downarrow$\\
    \textit{Stage 1 — Join resolution:} read \texttt{GROUP\_BY\_CONFIG[key].joins}\\
    + metric-driven joins; accumulate \texttt{joins\_needed} set\\
    $\downarrow$\\
    \textit{Stage 2 — SELECT:} dimension column(s) + \texttt{METRIC\_SELECT} expressions\\
    $\downarrow$\\
    \textit{Stage 3 — WHERE:} \texttt{\_PERIOD\_WHERE[dialect][period]}\\
    + date range + region + tier + status + category filters\\
    $\downarrow$\\
    \textit{Stage 4 — Finalise:} GROUP BY + HAVING + ORDER BY + LIMIT (TOP-N)
  }}
  \caption{Decision flow within the \texttt{\_build\_sql()} function at line 183 of \texttt{pipeline.py}. The YoY path produces a two-column CTE query; all other paths pass through four sequential construction stages.}
  \label{fig:sql-builder}
\end{figure}

\subsection{ResultVerifierAgent}
\label{subsec:result-verifier}

The ResultVerifierAgent executes the generated \ac{SQL} against the database, captures the result set, and performs three checks. First, it calls \texttt{calculate\_result\_plausibility()} to score each numeric column in the result against historical baselines using Z-score anomaly detection. Second, it checks whether the row count is within a ten-times multiplier of the expected count for equivalent queries. Third, if the \texttt{ANTHROPIC\_API\_KEY} is available, it passes the result to the \texttt{ClaudeResultVerifier} for natural language plausibility assessment. The combined plausibility score $a_p$ is a 50/50 blend of the Z-score-based aggregate plausibility and the row-count plausibility.

SQL execution errors (syntax errors, missing table or column errors, type mismatches) are caught within the ResultVerifierAgent's \texttt{process()} method. The exception is recorded in the trace log as an error message, the result set is set to an empty dictionary, and the row-count plausibility is set to 0.0. This means that a completely malformed SQL query produces $a_p = 0.0$, which in turn raises the composite drift score sufficiently that the critic loop will always attempt at least one refinement iteration. The design ensures that execution errors are never silently discarded: they always appear in the trace log, they always cause a nonzero drift contribution, and they always trigger a retry — providing three independent signals to both the user and the system that something went wrong.

The agent also records the column names and their inferred types from the result set's cursor description, attaching them to the \texttt{QueryResults} object. This schema information is returned in the \texttt{QueryResponse} and is used by the Streamlit dashboard's Schema Explorer view to display the result table with appropriate column formatting (numeric columns right-aligned, string columns left-aligned, date columns formatted as \texttt{YYYY-MM-DD}).

\section{Semantic Drift Metric Design}
\label{sec:drift-metric}

\subsection{Intent Alignment Component}
\label{subsec:intent-alignment}

The intent alignment score $a_i$ measures how faithfully the extracted intent and its ontology mappings reflect the user's original question. It is defined as:

\begin{equation}
  a_i = 0.6 \cdot \overline{C_e} + 0.4 \cdot \overline{S_m}
  \label{eq:intent-alignment}
\end{equation}

where $\overline{C_e}$ is the mean extraction confidence across all identified intent elements (entities, metrics, filters) and $\overline{S_m}$ is the mean ontology mapping similarity. The 60/40 weighting reflects the empirical observation that extraction confidence is a more reliable signal than mapping similarity for common business terms, because mapping similarity can be high even when the extraction identified the wrong element.

The 0.40 weight assigned to intent alignment in the composite drift formula (Equation~\ref{eq:drift}) was chosen to make it the dominant component, reflecting the design assumption that intent misidentification is the most consequential failure mode. If the system correctly identifies the entity and metric but maps to a slightly wrong alias (e.g., ``items'' resolved to ORDER\_ITEM rather than PRODUCT with similarity 0.82 rather than 1.0), the resulting mapping similarity penalty is small and unlikely to cause catastrophic output error. By contrast, a low extraction confidence typically indicates that the parser found no strong signal for a required dimension — a condition that almost always produces a wrong or incomplete \ac{SQL} query. An equal-weight (0.33/0.33/0.33) allocation was tested informally during development and found to produce drift scores that were insensitive to extraction failures, since the plausibility and constraint components could remain high while the intent component was low. The 0.40/0.30/0.30 split corrects this imbalance.

\subsection{Constraint Adherence Component}
\label{subsec:constraint-adherence}

The constraint adherence score $a_c$ is a simple ratio:

\begin{equation}
  a_c = \frac{N_{\text{satisfied}}}{N_{\text{total}}}
  \label{eq:constraint}
\end{equation}

where $N_{\text{satisfied}}$ is the count of constraint rules that evaluated to \texttt{True} and $N_{\text{total}}$ is the count of all rules applicable to the entities in the query. When no constraints apply, $a_c = 1.0$ by convention (as implemented at line~128 of \texttt{semantic\_drift.py}).

The 30\% weight allocated to constraint adherence in the composite formula reflects a deliberate decision to treat constraint violations as serious but not catastrophic. A query that returns the right entities and metrics but violates one constraint rule (for instance, computing revenue without netting out tax) is clearly wrong, but it is wrong in a specific, correctable way. Assigning 30\% to this component means that a complete constraint failure ($a_c = 0$) contributes 0.30 to the drift score, which is sufficient to push a query with otherwise good components (say $a_i = 0.90$, $a_p = 0.85$) above the 0.15 threshold ($D = 0.04 + 0.30 + 0.045 = 0.385$) and trigger the critic loop. The weighting was calibrated informally against the test query set: at 30\%, constraint failures were reliably surfaced through elevated drift while avoiding over-penalisation of queries where a minor rule ambiguity was the only issue.

Severity weights (HIGH/MEDIUM/LOW) for individual constraint rules are defined in the ontology but are not yet implemented in the adherence calculator. The current implementation treats all rules as equally weighted. Implementing severity-adjusted adherence — for example, weighting a HIGH-severity PII violation at 3$\times$ and a LOW-severity date-range advisory at 0.5$\times$ — would require changing the averaging step to a weighted sum at line~113 of \texttt{semantic\_drift.py}. This extension is straightforward and is listed as an enhancement candidate in the codebase's TODO comments.

The 30\% weight for constraint adherence in the composite formula is robust in the sense that the overall drift score ordering is not highly sensitive to moderate changes in this weight. Reducing the weight to 25\% (from 30\%) and redistributing the 5\% equally between intent (now 43\%) and plausibility (now 32\%) would shift the mean drift score of the ten test queries by less than 0.005 — a negligible change that would not alter any convergence decisions. This insensitivity to small weight perturbations is a desirable property: it means the system's behaviour is stable under calibration uncertainty, and a product team would not need to re-calibrate the weights precisely for each new domain. More substantial changes (e.g., reducing the constraint weight to 10\% while increasing intent to 60\%) would begin to matter: at 10\%, a complete constraint failure contributes only 0.10 to drift, which might not be enough to cross the 0.15 threshold for a query with otherwise high component scores. The current 30\% weight was set deliberately to ensure that HIGH-severity constraint violations always cause a threshold crossing, independent of how well the intent and plausibility components perform.

\subsection{Result Plausibility Component}
\label{subsec:plausibility}

The result plausibility score $a_p$ uses Z-score anomaly detection. For an observed aggregate value $x$ and a historical distribution with mean $\mu$ and standard deviation $\sigma$:

\begin{equation}
  z = \frac{|x - \mu|}{\sigma}
  \label{eq:zscore}
\end{equation}

\begin{equation}
  a_p^{\text{agg}} =
  \begin{cases}
    1.0 & \text{if } z \leq 2 \\
    0.0 & \text{if } z \geq 4 \\
    1 - \dfrac{z - 2}{2} & \text{otherwise}
  \end{cases}
  \label{eq:plausibility}
\end{equation}

The row-count plausibility $a_p^{\text{rc}}$ equals 1.0 if the returned row count is within ten times the historical mean, and 0.0 otherwise. The combined score is $a_p = 0.5 \cdot a_p^{\text{agg}} + 0.5 \cdot a_p^{\text{rc}}$.

\subsection{Composite Drift Score and Convergence}
\label{subsec:composite-drift}

The composite drift score combines the three components with fixed weights:

\begin{equation}
  D = 0.4(1 - a_i) + 0.3(1 - a_c) + 0.3(1 - a_p), \quad D \in [0, 1]
  \label{eq:drift}
\end{equation}

A query is considered converged — safe to serve — when:

\begin{equation}
  D < 0.15
  \label{eq:convergence}
\end{equation}

The 0.15 threshold was chosen empirically: it corresponds to an intent alignment above 0.75, full constraint satisfaction, and a plausibility score above 0.5. In practice, well-formed queries on common business dimensions satisfy this threshold on the first iteration in approximately 75\% of test cases.

The choice of 0.15 as a threshold merits further justification. A lower threshold (e.g., 0.05) would require near-perfect scores on all three components, which is achievable only for very simple, unambiguous queries — the system would reject most real-world questions even after five iterations. A higher threshold (e.g., 0.30) would pass queries that have material intent misalignment or constraint violations, undermining the purpose of the quality gate. The 0.15 level was selected by examining the distribution of drift scores across the ten test queries at iteration 1 and identifying the value that correctly classified the manually labelled ``acceptable'' and ``unacceptable'' results with no errors. This is admittedly a circular calibration on a small sample, and a principled calibration on a larger, independently labelled query set would be needed to validate the threshold for general use.

\section{QueryState Typed State Machine}
\label{sec:querystate}

The \texttt{QueryState} dataclass defined in \texttt{src/agents/state.py} acts as the communication bus between agents. It carries:

\begin{itemize}[leftmargin=2em]
  \item \texttt{query} — the original natural language string
  \item \texttt{intent} — the \texttt{ExtractedIntent} populated by the IntentParser
  \item \texttt{mappings} — the list of \texttt{OntologyMapping} objects from the Mapper
  \item \texttt{constraints} — the list of \texttt{ValidatedConstraint} objects from the Validator
  \item \texttt{plan} — the \texttt{ExecutionPlan} (SQL string, dialect, joins) from the Planner
  \item \texttt{results} — the \texttt{QueryResults} (rows, schema, row count) from the Verifier
  \item \texttt{drift\_scores} — a list of \texttt{DriftScore} snapshots, one per iteration
  \item \texttt{trace\_log} — a list of (agent, timestamp, message) tuples for auditability
  \item \texttt{iteration\_count}, \texttt{max\_iterations} — loop control integers
\end{itemize}

The \texttt{add\_trace(agent, message)} method appends a timestamped entry to \texttt{trace\_log}, and \texttt{get\_trace\_summary()} returns a formatted string of all trace entries. The \ac{API}'s \texttt{QueryResponse} Pydantic model includes an \texttt{execution\_trace} field that is populated from \texttt{get\_trace\_summary()}, giving API consumers a full audit log of every agent's contributions alongside the final result. This trace capability was instrumental in diagnosing the gross margin query's convergence behaviour during evaluation: inspection of the trace showed that the IntentParser mapped ``gross margin'' to PROFIT on iterations 1 and 2, and the feedback message from iteration 2 (``constraint R01 TAX\_EXCLUSION unsatisfied — metric PROFIT does not subtract tax'') caused it to map correctly to GROSS\_MARGIN on iteration 3.

Figure~\ref{fig:statemachine} shows the state transitions.

\begin{figure}[H]
  \centering
  \fbox{\parbox[c][6cm][c]{0.85\linewidth}{\centering
    \textbf{QueryState Transitions}\\[0.5em]
    \texttt{INIT} $\rightarrow$ \texttt{INTENT\_PARSED}\\
    $\rightarrow$ \texttt{ONTOLOGY\_MAPPED}\\
    $\rightarrow$ \texttt{CONSTRAINTS\_VALIDATED}\\
    $\rightarrow$ \texttt{SQL\_GENERATED}\\
    $\rightarrow$ \texttt{RESULTS\_VERIFIED}\\
    $\rightarrow$ [if $D \geq 0.15$ \& iter $<$ 5]: \texttt{INTENT\_PARSED} (loop)\\
    $\rightarrow$ [if $D < 0.15$ or iter = 5]: \texttt{COMPLETE}
  }}
  \caption{QueryState transitions through the pipeline. The back-edge from RESULTS\_VERIFIED to INTENT\_PARSED implements the critic loop.}
  \label{fig:statemachine}
\end{figure}

\section{Critic Loop and Orchestrator Design}
\label{sec:critic-loop}

The \texttt{SemanticGroundingOrchestrator} in \texttt{src/engine/orchestrator.py} manages the critic loop. Its \texttt{execute\_query()} async method initialises a \texttt{QueryState}, calls the pipeline once, and then enters a \texttt{while} loop gated on \texttt{state.iteration\_count < state.max\_iterations} and \texttt{not state.converged}. On each failed iteration, the orchestrator appends a feedback message to the state's \texttt{trace\_log} describing the specific component that scored poorly, giving the IntentParser context for refinement.

The term ``critic loop'' was chosen deliberately over ``retry loop'' or ``feedback loop'' to distinguish the design intent: the orchestrator does not blindly retry the same query, nor does it change the query arbitrarily. It acts as a critic — evaluating the previous attempt, diagnosing the failure, and providing targeted guidance to the first agent in the pipeline before the next attempt. This pattern is analogous to the teacher--student feedback loop in iterative learning: the critic's feedback is specific enough to guide improvement but not so prescriptive as to dictate the exact correction, leaving the IntentParser free to adjust its confidence thresholds and pattern priorities in response.

The design also considered a simpler alternative: having the orchestrator directly modify the extracted intent (e.g., force-setting the \texttt{exclude\_tax} flag when the TAX\_EXCLUSION constraint fails) rather than passing feedback through the agent chain. This would guarantee convergence on constraint adherence in a single extra iteration but would bypass the IntentParser's pattern matching, potentially introducing a different inconsistency between what the user asked and what the state records as the extracted intent. Preserving the agent chain — even at the cost of an extra iteration — maintains the invariant that every field in \texttt{QueryState} was populated by the designated agent, making the trace log a faithful record of the processing rather than a mix of agent-generated and orchestrator-injected values.

A secondary design question was whether to model the critic loop iterations as nodes in the LangGraph state graph (using a back-edge) or as a Python \texttt{while} loop outside the graph. The graph-based approach is conceptually cleaner — the entire workflow, including retries, is represented in the graph topology — but it requires the graph to support variable-length cycles, which requires careful handling of the iteration counter as part of the typed state. The \texttt{while} loop approach keeps the graph simpler (five sequential nodes, no back-edges) and moves the iteration logic into Python, where it is easier to inspect, test, and modify. Given that the iteration mechanics are a first-class research concern in this project, having them in plain Python code rather than embedded in graph traversal logic was the clearer choice. The trade-off is that the full workflow is no longer visualisable as a static graph; the convergence behaviour can only be understood by reading the orchestrator code. This is an acceptable trade-off for a research prototype but would need to be reconsidered in a production system where workflow visualisation supports operational monitoring.

The maximum iteration count of five was chosen conservatively. In the evaluation, no query required more than five iterations to either converge or be declared non-convergent. The diminishing-returns pattern observed in Section~\ref{sec:convergence} — where most improvement occurs on iterations one through three — suggests that three might be a more efficient maximum. However, a lower maximum also reduces the system's ability to recover from compounding failures where iteration 2's feedback inadvertently worsens a different component before iteration 3 corrects both. Setting the maximum at five provides enough headroom to recover from such transient oscillations while keeping worst-case latency within the non-functional requirement of 15 seconds.

Figure~\ref{fig:criticloop} illustrates the loop logic.

\begin{figure}[H]
  \centering
  \fbox{\parbox[c][6cm][c]{0.85\linewidth}{\centering
    \textbf{Critic Loop Flowchart}\\[0.5em]
    \texttt{receive(query)} $\rightarrow$ init \texttt{QueryState}\\
    $\downarrow$\\
    \texttt{run\_pipeline(state)}\\
    $\downarrow$\\
    \texttt{compute\_drift(state)} $\rightarrow$ $D$\\
    $\downarrow$\\
    \texttt{D < 0.15?}\\
    \quad YES $\rightarrow$ \texttt{return result}\\
    \quad NO $\rightarrow$ \texttt{iter++}; \texttt{iter < 5?}\\
    \qquad YES $\rightarrow$ \texttt{append feedback}; loop back\\
    \qquad NO $\rightarrow$ \texttt{return best result so far}
  }}
  \caption{Critic loop logic in the SemanticGroundingOrchestrator. The loop exits on convergence or after five iterations, returning the lowest-drift result seen.}
  \label{fig:criticloop}
\end{figure}

\section{Database Layer Design}
\label{sec:database-design}

Five SQLite databases were designed to represent distinct business domains: \texttt{sales.db} (customers, orders, order\_items, products, suppliers), \texttt{inventory.db} (product stock, warehouse locations), \texttt{analytics.db} (pre-aggregated daily metrics), \texttt{hr.db} (employees, departments, payroll), and \texttt{finance.db} (accounts, budget items, expenses). The synthetic data scale was set to match a medium-sized e-commerce operation: 50,000 customers, 300,000 orders, 800,000 order items, 10,000 products, 500 suppliers, 180,000 invoices, 600,000 transactions, and 30,000 returns.

The separation into five databases reflects a real-world pattern in which different business domains are managed by separate teams with separate data ownership. By keeping the domains separated rather than normalising into a single schema, the design creates an environment in which cross-domain query errors are more detectable: a query that accidentally joins the HR employee table to the sales orders table will produce a CartesianProduct-style join explosion, immediately visible as a row-count anomaly in the plausibility calculator.

Each database uses SQLAlchemy declarative models with explicit \texttt{Integer}, \texttt{Float}, \texttt{String}, and \texttt{DateTime} column types. The schema mirrors the ontology entity definitions: the \texttt{customers} table has \texttt{customer\_id} as primary key, matching the \texttt{CUSTOMER} entity's PK field; the \texttt{orders} table has a foreign key to \texttt{customers.customer\_id} matching the CUSTOMER $\to$ ORDER join edge. This alignment between ontology and schema ensures that the join conditions generated by the SQL builder are always syntactically valid against the actual database, even without a schema inspection step at query time.

Synthetic data generation uses realistic distributions where feasible. Customer regions are drawn from four weighted values (North America 40\%, Europe 30\%, Asia Pacific 20\%, Latin America 10\%). Order statuses follow an approximate real-world split: 65\% Delivered, 15\% Shipped, 10\% Processing, 10\% Cancelled. Product categories cover eight domains: Electronics (25\%), Clothing (20\%), Home \& Garden (15\%), Sports (15\%), Books (10\%), Furniture (8\%), Food (4\%), and Beauty (3\%). Order amounts follow a log-normal distribution with parameters calibrated to produce a mean of approximately \$85 and a 95th percentile of around \$350, matching typical B2C e-commerce transaction profiles.

\section{REST API Design}
\label{sec:api-design}

The \ac{REST} \ac{API} was designed with three consumers in mind: the Streamlit dashboard (internal consumer), integration test scripts (automated consumer), and future third-party client applications (external consumer). All endpoints return JSON with a consistent envelope structure: a \texttt{status} field (\texttt{"ok"} or \texttt{"error"}), a \texttt{data} field carrying the response payload, and an optional \texttt{message} field for human-readable error descriptions. This envelope pattern ensures that clients can reliably check the \texttt{status} field before attempting to parse \texttt{data}, avoiding the need to distinguish success from error by HTTP status code alone.

The \texttt{/query} endpoint is the most complex: it accepts a \texttt{QueryRequest} body with a \texttt{query} string and an optional \texttt{dialect} field (\texttt{"sqlite"} or \texttt{"postgresql"}), runs the full critic loop, and returns a \texttt{QueryResponse} that includes the result rows as a list of dictionaries, the final \ac{SQL} string, the drift score components, the iteration count, the execution trace, and a human-readable explanation generated by \texttt{\_generate\_explanation()}. The response serialisation uses Pydantic's \texttt{model\_json\_schema()} to validate that all required fields are present before returning, preventing partially constructed responses from reaching clients.

Table~\ref{tab:api-endpoints} lists the eight \ac{REST} endpoints exposed by the FastAPI backend.

\begin{table}[H]
  \centering
  \caption{REST API Endpoint Catalogue}
  \label{tab:api-endpoints}
  \renewcommand{\arraystretch}{1.2}
  \begin{tabularx}{\textwidth}{@{}llXl@{}}
    \toprule
    \textbf{Method} & \textbf{Path} & \textbf{Description} & \textbf{Auth} \\
    \midrule
    POST & /query    & Submit NL query; returns result, SQL, drift score & No \\
    GET  & /health   & System health check; returns component statuses  & No \\
    GET  & /schema   & Returns database schema as JSON                  & No \\
    GET  & /ontology & Returns full ontology entity and metric catalogue & No \\
    GET  & /rules    & Returns active constraint rules list              & No \\
    GET  & /live     & Returns real-time processing statistics           & No \\
    GET  & /history  & Returns last \textit{n} query records from log   & No \\
    GET  & /stats    & Returns aggregate success rate and iteration KPIs & No \\
    \bottomrule
  \end{tabularx}
\end{table}

\section{Streamlit Dashboard Design}
\label{sec:dashboard-design}

The monitoring dashboard was designed to serve two audiences simultaneously: business users who want to submit queries and inspect results, and engineers who need visibility into the system's internal quality signals. This dual-audience requirement drove the layout choices: business-facing elements (query input, result table, SQL display) occupy the primary viewport area; engineering-facing elements (drift score bars, iteration trace, component breakdown) are accessible in expandable panels that do not clutter the main view for users who are not interested in them. Four views were identified as necessary. The \textit{Query Interface} view provides a text input for natural language queries, displays the generated \ac{SQL} alongside the result table, and renders a three-bar chart showing the individual drift components — a design choice that makes the quality score interpretable rather than opaque. The \textit{Drift Analysis} view displays iteration-by-iteration convergence graphs for recent queries, using the \texttt{/history} endpoint as its data source. The \textit{System Metrics} view polls the \texttt{/stats} endpoint to display aggregate success rate, mean iteration count, and per-component drift averages. The \textit{Schema Explorer} view calls \texttt{/schema} to render the database schema as a collapsible tree, helping users understand what entities and columns are available before formulating a query.

The dashboard communicates with the \ac{API} over HTTP, meaning it can run in a separate process or container from the backend — an architectural decision that keeps the rendering layer decoupled from the query processing layer. The dark theme and wide layout mode (\texttt{st.set\_page\_config(layout="wide")}) were chosen to match common business intelligence tool aesthetics and to make the three-column layout of query, \ac{SQL}, and drift chart fit comfortably on a standard 1920-wide monitor.

Streamlit's session state (\texttt{st.session\_state}) was used to persist the last ten query results across page refreshes, enabling the Drift Analysis view to render convergence histories without an additional database round-trip. The real-time polling on the System Metrics view uses \texttt{st.rerun()} with a configurable interval (defaulting to 10 seconds) rather than a WebSocket subscription, which would require a dedicated async back-end. This polling approach is acceptable for a monitoring dashboard but would not scale to a high-frequency trading or alerting context where latency matters.

\section{Technology Selection Rationale}
\label{sec:technology}

Technology choices in this project were guided by three criteria: (1) the degree to which the technology supports the auditability principle; (2) the availability of a pure-Python implementation that does not require external services; and (3) the maturity of the package's documentation and community support. Table~\ref{tab:technology} summarises the key technology choices and the reasoning behind each.

Several alternatives were explicitly evaluated and rejected. Celery was considered as an orchestration layer but rejected because its worker queue model introduces distributed state that complicates debugging: a failed task in a Celery chain is harder to trace than a Python exception propagating synchronously through an async call stack. Apache Kafka was considered for inter-agent message passing but rejected because the project's sequential pipeline does not benefit from message decoupling — each agent's output is immediately consumed by the next agent in the same Python process, making a queue an unnecessary overhead. Flask was considered in place of FastAPI but rejected because FastAPI's built-in Pydantic validation and OpenAPI documentation generation reduce the amount of boilerplate code required for the eight-endpoint API. The SQLite/PostgreSQL dual-dialect approach was preferred over a single-target approach because it enables local development on SQLite without requiring a running PostgreSQL server, while preserving a direct path to production-grade deployment.

\begin{table}[H]
  \centering
  \caption{Technology Stack and Justification}
  \label{tab:technology}
  \renewcommand{\arraystretch}{1.2}
  \begin{tabularx}{\textwidth}{@{}llX@{}}
    \toprule
    \textbf{Component} & \textbf{Technology} & \textbf{Rationale} \\
    \midrule
    Pipeline orchestration & LangGraph \cite{langgraph}    & Typed state graph; auditability; back-edge support \\
    Ontology graph         & NetworkX \cite{networkx}      & Lightweight; native shortest-path; pure Python \\
    Web framework          & FastAPI \cite{fastapi}        & Async; OpenAPI generation; Pydantic validation \\
    Database               & SQLite / PostgreSQL           & Zero-config development; production-grade production \\
    ORM                    & SQLAlchemy                    & Dialect-agnostic; chunked inserts; schema inspection \\
    Dashboard              & Streamlit                     & Rapid prototyping; dark theme; wide ecosystem \\
    LLM integration        & Claude claude-haiku-4-5 (optional) & Low latency; strong SQL generation; API-key gated \\
    Containerisation       & Docker Compose                & Reproducible environment; network isolation \\
    \bottomrule
  \end{tabularx}
\end{table}

\section{Non-Functional Requirements}
\label{sec:nfr}

Beyond functional correctness, the system was subject to several non-functional requirements that shaped key design decisions. Query latency end-to-end (from \ac{API} call to result return) should be under 5 seconds for single-pass queries and under 15 seconds for five-iteration critic-loop cases. The drift computation itself should take under 10 milliseconds per call, which was achieved with the pure-Python Z-score calculation in \texttt{semantic\_drift.py}. The \ac{API} should handle at least ten concurrent requests without degradation, supported by FastAPI's async architecture and thread-safe \texttt{deque(maxlen=1000)} query log. Memory usage at idle should be under 500~MB, a constraint met by the SQLite file-based storage model.

Auditability was treated as a first-class non-functional requirement. Every query must be accompanied by a complete execution trace that records which agent produced which output, what the drift scores were at each iteration, and whether convergence was achieved. This requirement drove the design of the \texttt{QueryState.trace\_log} field and the \texttt{add\_trace()} API, which adds a timestamped record at each agent transition. The trace is returned in the \texttt{QueryResponse} as a structured string, enabling consumers to log it to their own audit systems.

Graceful degradation was required at two levels. At the service level, the system must continue processing pattern-based queries when the Anthropic API is unavailable; the Claude-based agent wrappers check for an available API key at instantiation time and fall back to their deterministic counterparts silently. At the database level, if a generated SQL query produces a runtime error (invalid column reference, type mismatch), the ResultVerifier must catch the exception, return a zero-row result, and set plausibility to 0.0 — allowing the drift calculator and critic loop to continue rather than propagating the exception to the caller.

Observability was a further non-functional requirement. The structured logging configuration in \texttt{src/logging\_config.py} uses \texttt{colorlog} to format log entries with severity-appropriate colours in the terminal and writes machine-readable JSON lines to \texttt{logs/app.log}. Every log entry carries a \texttt{correlation\_id} field derived from the \texttt{QueryState.query\_id}, making it possible to trace all log entries from a single query execution by grepping for that identifier. This is particularly useful when debugging multi-iteration critic-loop runs, where the same query ID appears across five separate pipeline invocations. The log file is limited to 10~MB with a two-file rotation policy, implemented with Python's \texttt{RotatingFileHandler}.

Security requirements were deliberately minimal for an academic prototype. The API exposes no authentication because it is intended to run on a private network or local machine. A production deployment would need at least API key authentication on the \texttt{/query} endpoint and rate limiting to prevent denial-of-service. The PII\_PROTECTION constraint rule provides a first line of defence at the query level, ensuring that raw email addresses and phone numbers are never selected, but this does not substitute for proper access control at the service boundary. These gaps are acknowledged as scope limitations rather than design oversights.

% =============================================================================
% CHAPTER 4 — IMPLEMENTATION
% =============================================================================
\chapter{Implementation}
\label{ch:implementation}

\section{Development Environment and Project Structure}
\label{sec:dev-env}

The project was developed on macOS (Darwin 25.0.0, Apple M2 architecture) using Python 3.13 in a virtual environment located at \texttt{.venv\_mac/}. Python 3.13 was selected because it provides improved error messages, a faster interpreter for tight loops (relevant to batch data generation), and the \texttt{tomllib} standard library module for configuration parsing. Dependencies were managed through two requirements files: \texttt{requirements.txt} for general use and \texttt{requirements-docker.txt} for the containerised deployment, where package versions are pinned to ensure reproducible builds.

The key runtime dependencies and their versions used in development are: \texttt{fastapi==0.104.1}, \texttt{uvicorn==0.24.0}, \texttt{sqlalchemy==2.0.23}, \texttt{networkx==3.2.1}, \texttt{streamlit==1.28.2}, \texttt{pydantic==2.5.0}, \texttt{python-dotenv==1.0.0}, and \texttt{colorlog==6.7.0}. The \texttt{langgraph} package was not pinned in \texttt{requirements.txt} because the mock fallback ensures the system operates without it; when installed, version 0.0.38 was used during development. The \texttt{anthropic} package (version 0.8.1) was similarly optional.

The top-level project structure follows a conventional layout:

\begin{itemize}[leftmargin=2em]
  \item \texttt{src/agents/} — pipeline agents (\texttt{pipeline.py}, \texttt{state.py}), Claude agent wrappers (\texttt{claude\_agents.py})
  \item \texttt{src/engine/} — ontology (\texttt{ontology.py}), semantic drift (\texttt{semantic\_drift.py}), orchestrator (\texttt{orchestrator.py}), mock LangGraph (\texttt{mock\_langgraph.py})
  \item \texttt{src/db/} — database manager and synthetic data generator (\texttt{manager.py})
  \item \texttt{src/api/} — FastAPI application (\texttt{backend.py})
  \item \texttt{src/dashboards/} — Streamlit application (\texttt{streamlit\_app.py})
  \item \texttt{src/config.py}, \texttt{src/logging\_config.py} — configuration and structured logging
  \item \texttt{tests/} — pytest test suite (5 test files)
  \item \texttt{data/ontology/} — optional ontology override JSON files
  \item \texttt{docker/} — \texttt{docker-compose.yml}
  \item \texttt{docs/} — technical documentation (architecture, drift metric, synopsis)
\end{itemize}

\section{Business Ontology Implementation}
\label{sec:ontology-impl}

The \texttt{BusinessOntology} class in \texttt{src/engine/ontology.py} stores all entities, metrics, and constraints as Python dictionaries keyed on their canonical ID. On construction, \texttt{\_populate\_default\_ontology()} first attempts to load definitions from JSON files in the \texttt{data/ontology/} directory via \texttt{\_load\_from\_json\_files()}: if the directory exists and contains at least one valid JSON file, that data takes precedence. Only when no JSON files are found does the initialiser fall back to \texttt{\_load\_hardcoded\_defaults()}, which populates the same dictionaries from hard-coded Python literals.

The JSON loading function reads three file types by convention: \texttt{entities.json} (a list of entity definition objects), \texttt{metrics.json} (a list of metric definition objects), and \texttt{constraints.json} (a list of constraint rule objects). Each file is loaded independently using \texttt{json.load()}, and the resulting objects are validated against a lightweight schema (checking for required keys such as \texttt{entity\_id}, \texttt{table\_name}, and \texttt{aliases}) before being added to the dictionaries. Invalid entries are logged as warnings and skipped rather than raising exceptions, implementing the graceful-degradation principle at the data loading level. A future version could replace this ad hoc validation with a Pydantic model validation step, which would provide cleaner error messages when ontology JSON files contain typos or missing fields.

The \texttt{aliases} field in each entity definition is an important but easily overlooked part of the schema. Its purpose is to extend the ontology's vocabulary beyond the canonical entity identifier to include the natural language terms business users actually use. For the CUSTOMER entity, the alias list includes ``customer'', ``client'', ``buyer'', ``user'', and ``account''. For ORDER, it includes ``order'', ``purchase'', ``sale'', ``transaction'', and ``booking''. The breadth of these lists directly determines the ontology's mapping recall — an entity with a narrow alias list will be missed when the user phrases their question using an uncommon synonym. During development, alias gaps were the most common source of mapping failures in informal testing, leading to three rounds of alias list expansion before the evaluation query set was run. This empirical process could be formalised in future work by mining a query log for unmapped tokens and feeding them back to an ontology editor.

The \texttt{\_build\_ontology\_graph()} method, called after the dictionaries are populated, iterates over the \texttt{join\_rules} list to construct the \texttt{NetworkX} directed graph. Each join rule specifies a \texttt{from\_entity} and a \texttt{to\_entity} (both must be keys in the entities dictionary), a \texttt{join\_type} (LEFT or INNER), a \texttt{condition} string (the SQL ON clause), and a \texttt{weight} float for shortest-path computation. The method validates that both entity keys exist in the graph before adding the edge; if either is missing, the join rule is silently skipped. This means the graph is always a valid directed graph even when the ontology is partially defined — useful during incremental ontology development when new entities are added before their join rules are complete. This priority logic means that business analysts can extend or override the ontology by placing JSON files in \texttt{data/ontology/} without touching Python source code. Following dictionary population, \texttt{\_build\_ontology\_graph()} constructs the \texttt{NetworkX} \texttt{DiGraph}: it iterates over the eight defined join rules, adding a directed edge for each join with attributes including \texttt{join\_type} (LEFT, INNER), \texttt{condition} (the ON clause), and \texttt{weight} (used by \texttt{nx.shortest\_path}).

Entity and metric lookup supports both exact and alias-based matching. \texttt{find\_entity\_by\_name()} first attempts a case-insensitive exact match on entity\_id, then iterates through the aliases list of each entity. \texttt{find\_metric\_by\_name()} follows the same logic. This two-pass lookup allowed the intent parser to resolve colloquial terms such as ``client'' (to CUSTOMER) and ``purchase'' (to ORDER) without requiring separate synonym expansion logic.

An optional ChromaDB integration path was implemented: if the \texttt{chromadb} package is installed and a collection named \texttt{ontology\_cosine} exists, the mapper can fall back to embedding-based nearest-neighbour search for terms that fail the string-similarity threshold. This path was not exercised in the current evaluation because a local ChromaDB instance was not provisioned, but the hook is in place for future work.

\section{Agent Pipeline Implementation}
\label{sec:pipeline-impl}

\subsection{Intent Extraction}
\label{subsec:intent-impl}

The twelve metric SQL expressions in \texttt{METRIC\_SELECT} (lines~39--52 of \texttt{pipeline.py}) form the core of the SQL generation layer and are shown below:

\begin{tcblisting}{listing only, colback=gray!10!white, breakable, boxsep=0pt,
  top=1mm, bottom=1mm, left=1mm, right=1mm,
  listing options={
    language=Python, basicstyle=\small\ttfamily,
    keywordstyle=\color{kwcolor}, commentstyle=\color{cmtcolor},
    stringstyle=\color{strcolor}, breaklines=true,
    numbers=left, numberstyle=\tiny\color{gray}}}
# pipeline.py lines 39-52 — METRIC_SELECT dictionary
METRIC_SELECT = {
    "total_revenue":   "SUM(o.total_amount) AS total_revenue",
    "total_orders":    "COUNT(*) AS total_orders",
    "avg_order":       "ROUND(AVG(o.total_amount), 2) AS avg_order_value",
    "total_profit":    ("SUM(o.total_amount - o.tax_amount "
                        "- o.shipping_cost) AS total_profit"),
    "margin_pct":      ("ROUND(SUM(o.total_amount - o.tax_amount) * 100.0 "
                        "/ NULLIF(SUM(o.total_amount), 0), 2) AS gross_margin_pct"),
    "unique_customers": "COUNT(DISTINCT o.customer_id) AS unique_customers",
    "cancelled_count":  ("SUM(CASE WHEN o.status='Cancelled' "
                         "THEN 1 ELSE 0 END) AS cancelled_orders"),
    "avg_ltv":          ("ROUND(SUM(o.total_amount) / "
                         "NULLIF(COUNT(DISTINCT o.customer_id), 0), 2) AS avg_ltv"),
    "cancel_pct":       ("ROUND(SUM(CASE WHEN o.status='Cancelled' THEN 1 "
                         "ELSE 0 END) * 100.0 / NULLIF(COUNT(*), 0), 2) "
                         "AS cancellation_rate"),
    "total_stock":      "SUM(p.stock_quantity) AS total_inventory",
    "net_revenue":      ("SUM(o.total_amount - o.tax_amount) AS net_revenue"),
    "order_count":      "COUNT(DISTINCT o.order_id) AS order_count"
}
\end{tcblisting}

Each expression is a self-contained SQL fragment including an alias clause. The SQL builder selects from this dictionary by metric key and appends the expression directly to the \texttt{SELECT} clause, making the metric templates easy to audit and extend independently of the SQL builder logic.

The IntentParserAgent's \texttt{process()} method orchestrates the ten filter detectors listed in Section~\ref{subsec:intent-parser}. Each detector is a static method returning either a detected value with confidence or \texttt{None}. The pattern for entity extraction deserves mention: rather than a single monolithic regex, eight patterns are compiled at module load time and tried in sequence. The first match wins; ties at equal confidence are broken by specificity (longer pattern wins). This avoids the common pitfall of an omnibus regex that silently mis-matches when token order varies.

Confidence scaling reflects the number of explicit signals found. A query containing both an entity word (``customer'') and a specific metric word (``lifetime value'') and a temporal filter (``last year'') achieves a high base confidence because all three dimensions are unambiguous. Conversely, a short query such as ``show sales'' produces a lower confidence because the metric is implicit and no filter is present.

One subtlety in the extraction design is the handling of negation. The phrase ``excluding taxes'' is detected by a pattern that matches variants of ``exclud\{ing,e\}'', ``without'', ``net of'', and ``before tax''. The extracted value (\texttt{exclude\_tax=True}) is then used both by the ConstraintValidatorAgent (to verify that the TAX\_EXCLUSION rule is satisfied) and by the SQL builder (to include the \texttt{- o.tax\_amount} term in the SELECT expression). This dual use ensures that the metric formula and the constraint check are always in sync — a pattern where the SQL builder uses a different tax treatment from the one the constraint validator checks would produce an artificially high constraint adherence score while the generated SQL silently includes tax.

Another consideration was the treatment of compound queries — questions that implicitly ask for multiple metrics, such as ``what is the revenue and order count by region?'' The current implementation detects the first metric pattern match and uses it exclusively; a compound query would require a multi-metric SELECT, which the current SQL builder does not support. This limitation is logged as a warning in the trace when more than one metric pattern matches, and the intent alignment confidence is reduced by 0.1 as a signal that the extraction may be incomplete.

\subsection{SQL Generation}
\label{subsec:sql-impl}

The \texttt{\_build\_sql()} function at line~183 of \texttt{pipeline.py} is the most complex single function in the codebase at approximately 130 lines. It was the function that required the most iterative refinement during development: the initial single-pass implementation handled only region and tier groupings; subsequent iterations added the remaining seven GROUP BY configurations, the YoY CTE branch, HAVING clause support, TOP-N limiting, explicit date-range filtering, and the dual-dialect switch. The final function is not particularly elegant — it is a linear sequence of conditional additions to a set of SQL fragment lists — but it is fully deterministic, easy to trace through the execution log, and covered by the SQL builder test suite.

Its entry point checks for the year-over-year flag first (lines~217--244), returning a two-column CTE-style query when detected. For all other cases, it proceeds through four stages:

\begin{enumerate}[leftmargin=2em,label=\arabic*.]
  \item \textbf{Join resolution} — reads \texttt{GROUP\_BY\_CONFIG[group\_by\_key].joins} and any metric-driven join requirements, accumulating the needed table set in \texttt{joins\_needed}.
  \item \textbf{SELECT construction} — prepends the grouping dimension column(s), then appends metric \ac{SQL} expressions from \texttt{METRIC\_SELECT}.
  \item \textbf{WHERE construction} — adds the period filter from \texttt{\_PERIOD\_WHERE[dialect]}, plus any explicit date range, region, tier, status, or category filters.
  \item \textbf{Finalisation} — adds GROUP BY, HAVING (if a threshold was detected), ORDER BY (from \texttt{GROUP\_BY\_CONFIG}), and TOP-N LIMIT.
\end{enumerate}

The \texttt{\_PERIOD\_WHERE} dictionary (lines~123--140 of \texttt{pipeline.py}) stores dialect-specific SQL fragments for six temporal filter types: \texttt{default} (last 12 months), \texttt{last\_month}, \texttt{ytd} (year-to-date), \texttt{last\_quarter}, \texttt{current\_year}, and \texttt{last\_year}. Each entry contains the SQLite variant and the PostgreSQL variant of the WHERE clause fragment. For example, the \texttt{ytd} entry uses \texttt{strftime('\%Y', o.order\_date) = strftime('\%Y', 'now')} for SQLite and \texttt{EXTRACT(YEAR FROM o.order\_date) = EXTRACT(YEAR FROM CURRENT\_DATE)} for PostgreSQL. The SQL builder selects the appropriate fragment using \texttt{period\_map = \_PERIOD\_WHERE[dialect]} and then indexes by the detected period type, or falls back to \texttt{period\_map['default']} if no period was detected.

The dialect switch between SQLite and PostgreSQL date functions is handled by choosing the appropriate key (\texttt{select\_sqlite} vs \texttt{select\_pg}, \texttt{group\_sqlite} vs \texttt{group\_pg}) within the GROUP BY configuration dictionary, rather than a single large \texttt{if dialect} block. This keeps the dialect differences local to the data structure rather than scattered through the generation logic.

The nine GROUP BY configurations in \texttt{GROUP\_BY\_CONFIG} (lines~64--120 of \texttt{pipeline.py}) each specify at minimum: a \texttt{select} or pair of dialect-specific \texttt{select\_sqlite}/\texttt{select\_pg} expressions for the grouping dimension column, a \texttt{group} or dialect-pair for the GROUP BY clause, an \texttt{order} specification, and a \texttt{joins} list naming the additional table joins required. For example, the \texttt{"region"} configuration specifies \texttt{select = "c.region"}, \texttt{group = "c.region"}, \texttt{order = "total\_revenue DESC"}, and \texttt{joins = ["customers"]}. The \texttt{"month"} configuration uses \texttt{select\_sqlite = "strftime('\%Y-\%m', o.order\_date) AS month"} and \texttt{select\_pg = "TO\_CHAR(o.order\_date, 'YYYY-MM') AS month"} to handle the SQLite/PostgreSQL date formatting difference. This dictionary-driven design means that adding a new grouping dimension (for example, grouping by warehouse location) requires adding one new dictionary entry with no changes to the SQL builder function itself.

The year-over-year branch (lines~217--244) constructs a Common Table Expression with two named subqueries: \texttt{this\_year} and \texttt{last\_year}, each using the \texttt{\_PERIOD\_WHERE} template for the respective year. The outer SELECT joins the two CTEs on their grouping dimension and computes the absolute and percentage change. This approach was preferred over a single query with a \texttt{CASE WHEN year = X THEN ...} pattern because the CTE form is more readable, easier to debug from the execution trace, and produces better-labelled result columns.

\begin{tcblisting}{listing only, colback=gray!10!white, breakable, boxsep=0pt,
  top=1mm, bottom=1mm, left=1mm, right=1mm,
  listing options={
    language=Python, basicstyle=\small\ttfamily,
    keywordstyle=\color{kwcolor}, commentstyle=\color{cmtcolor},
    stringstyle=\color{strcolor}, breaklines=true}}
# pipeline.py line 183 — abbreviated illustrative excerpt
def _build_sql(intent, dialect: str = "sqlite") -> str:
    d = dialect if dialect in _PERIOD_WHERE else "sqlite"
    period_map = _PERIOD_WHERE[d]
    if not intent:
        return (f"SELECT {DEFAULT_SELECT} FROM orders o "
                f"WHERE {period_map['default']} AND o.total_amount > 0 LIMIT 100")
    filters = intent.filters if intent.filters else {}
    group_by_key = filters.get("group_by", "")
    # ... (join resolution, SELECT, WHERE, GROUP BY construction follow)
\end{tcblisting}

The constraint checking logic in \texttt{\_check\_rule()} (lines~159--180 of \texttt{pipeline.py}) evaluates each rule by a \texttt{rule\_type} switch:

\begin{tcblisting}{listing only, colback=gray!10!white, breakable, boxsep=0pt,
  top=1mm, bottom=1mm, left=1mm, right=1mm,
  listing options={
    language=Python, basicstyle=\small\ttfamily,
    keywordstyle=\color{kwcolor}, commentstyle=\color{cmtcolor},
    stringstyle=\color{strcolor}, breaklines=true,
    numbers=left, numberstyle=\tiny\color{gray}}}
# pipeline.py lines 159-180 — _check_rule() (condensed)
def _check_rule(rule: dict, intent, filters: dict) -> bool:
    rule_type = rule.get("rule_type", "")
    if rule_type == "TAX_EXCLUSION":
        return filters.get("exclude_tax", False) or \
               intent.metric in TAX_AWARE_METRICS
    elif rule_type == "REGION_FILTER":
        region = filters.get("region")
        return region is None or region in _VALID_REGIONS
    elif rule_type == "PII_PROTECTION":
        return True  # pipeline never selects raw PII columns
    elif rule_type == "ACTIVE_PRODUCTS_ONLY":
        return True  # SQL builder appends is_active=1 by default
    elif rule_type == "MIN_ORDER_VALUE":
        return True  # SQL builder appends total_amount > 0 by default
    elif rule_type == "DATE_RANGE":
        return filters.get("year") is None or \
               1990 <= int(filters["year"]) <= 2030
    return True     # unknown rule types pass by convention
\end{tcblisting}

This deterministic approach ensures that constraint adherence scoring is reproducible and requires no neural inference. The conservative default (unknown rule types pass) was a deliberate choice to prevent new rules added to the ontology JSON from breaking existing queries before their evaluation logic is implemented in the validator.

\subsection{Claude AI Integration}
\label{subsec:claude-impl}

When the \texttt{ANTHROPIC\_API\_KEY} environment variable is present, the \texttt{AgentPipeline} class substitutes the deterministic \texttt{ExecutionPlannerAgent} with \texttt{ClaudeExecutionPlanner} and \texttt{ResultVerifierAgent} with \texttt{ClaudeResultVerifier}. Both wrappers are implemented in \texttt{src/agents/claude\_agents.py} and use the Anthropic Python \ac{SDK}. The Claude-based planner receives the full ontology context — entity definitions, metric SQL templates, and applicable constraint rules — as part of its system prompt, grounding generation in the same knowledge base used by the deterministic agents. This design means the drift metric can score a Claude-generated query using exactly the same component calculators as a pattern-generated query, maintaining comparability.

The system prompt for \texttt{ClaudeExecutionPlanner} was designed to be concise and directive rather than exhaustive. It provides: (1) the resolved entity set and their table names; (2) the applicable metric SQL template verbatim; (3) the list of constraint rules with their descriptions; and (4) the detected filter values (period, group-by key, having threshold, etc.). The instruction to the model is simply to produce a single, syntactically valid \ac{SQL} query for the specified dialect. No chain-of-thought or step-by-step reasoning is requested, on the grounds that the intermediate reasoning is not needed for scoring and adds latency. The generated SQL is extracted from the response using a code-block parser that looks for the first \texttt{```sql ... ```} fence in the response string.

The \texttt{ClaudeResultVerifier} supplements the Z-score plausibility check with a natural language assessment. It receives the query, the generated SQL, and the result rows (truncated to 10 rows for API efficiency) and is asked to assess on a scale of 0--1 whether the result appears to answer the original question faithfully. This score is blended with the Z-score plausibility at a 30/70 weight (30\% Claude assessment, 70\% Z-score) when the API is available, reflecting the lower reliability of LLM-based plausibility judgments on numeric results compared to the deterministic Z-score calculation. The blending weight is configurable through \texttt{config.py}.

\section{Semantic Drift Metric Implementation}
\label{sec:drift-impl}

The \texttt{SemanticDriftMetric} class in \texttt{semantic\_drift.py} coordinates the three calculator classes. Its \texttt{calculate()} method accepts a \texttt{QueryState} and returns a \texttt{DriftMetricComponents} dataclass containing all four scores ($a_i$, $a_c$, $a_p$, $D$) plus a \texttt{breakdown} dictionary for dashboard display.

\texttt{IntentAlignmentCalculator.calculate\_alignment\_score()} (line~73) accepts two arguments: a dictionary of \{element: confidence\} pairs from the intent parser, and a list of mapping similarity floats from the ontology mapper. It computes the 60/40 weighted average as per Equation~\ref{eq:intent-alignment}. Edge cases are handled conservatively: an empty extraction confidence dictionary yields 0.0 rather than 1.0, preventing the score from defaulting to high alignment when the parser found nothing.

The \texttt{ConstraintAdherenceCalculator.calculate\_adherence\_score()} (line~113) receives a list of constraint evaluation dictionaries, each containing a \texttt{rule\_type} field and a \texttt{satisfied} boolean. It counts the number of dictionaries where \texttt{satisfied} is \texttt{True} and divides by the total count. When the list is empty (no applicable constraints), it returns 1.0 by convention, implemented as a guard clause: \texttt{if not constraints: return 1.0}. This convention means that a query that touches only entities with no associated constraint rules receives a perfect constraint score, which is the correct semantic: there are no constraints to violate.

The \texttt{SemanticDriftMetric} orchestration class (the top-level entry point) also maintains a \texttt{historical\_baselines} dictionary that maps query pattern signatures to lists of historical aggregate values. On each call to \texttt{calculate()}, it checks whether a baseline exists for the current query's metric and group-by key combination; if one does, it passes it to \texttt{ResultPlausibilityCalculator}; if not, it passes an empty list, which triggers the ``insufficient history'' fallback that returns plausibility 1.0. The baseline update logic — appending the current result to the history after a converged query — was not implemented in the current version, meaning all plausibility checks use the pre-populated synthetic baselines. A production deployment would need a persistent baseline store (a small SQLite table or a Redis hash would suffice) and a baseline-update call after each converged query result.

\texttt{ResultPlausibilityCalculator.calculate\_z\_score\_anomaly()} (line~172) uses Python's \texttt{statistics.stdev()} on the historical baseline list. When fewer than two historical values are available, the function returns 0.0 (no anomaly) rather than raising an exception. The normalisation from Z-score to plausibility (Equation~\ref{eq:plausibility}) is implemented in a companion function \texttt{normalize\_z\_score\_to\_plausibility()}: for $z \leq 2$ it returns 1.0 directly; for $z \geq 4$ it returns 0.0; in the linear interpolation band it uses \texttt{1.0 - (z - 2) / 2.0}.

The plausibility calculation at line~172 of \texttt{semantic\_drift.py} is reproduced below for reference:

\begin{tcblisting}{listing only, colback=gray!10!white, breakable, boxsep=0pt,
  top=1mm, bottom=1mm, left=1mm, right=1mm,
  listing options={
    language=Python, basicstyle=\small\ttfamily,
    keywordstyle=\color{kwcolor}, commentstyle=\color{cmtcolor},
    stringstyle=\color{strcolor}, breaklines=true,
    numbers=left, numberstyle=\tiny\color{gray}}}
# semantic_drift.py — ResultPlausibilityCalculator (line 172)
def calculate_result_plausibility(results, baselines: dict) -> float:
    if not results or not results.get("rows"):
        return 0.0
    row_count = len(results["rows"])
    expected_rows = baselines.get("expected_row_count", row_count)
    row_plausibility = 1.0 if row_count <= expected_rows * 10 else 0.0

    numeric_vals = _extract_numeric_values(results)
    if len(numeric_vals) < 1:
        return 0.5 * row_plausibility + 0.5  # no aggregates to check
    hist = baselines.get("historical_aggregates", [])
    if len(hist) < 2:
        agg_plausibility = 1.0   # insufficient history — no penalty
    else:
        mu = statistics.mean(hist)
        sigma = statistics.stdev(hist) or 1.0
        z = abs(statistics.mean(numeric_vals) - mu) / sigma
        agg_plausibility = normalize_z_score_to_plausibility(z)
    return 0.5 * agg_plausibility + 0.5 * row_plausibility
\end{tcblisting}

The 50/50 blend between aggregate plausibility and row-count plausibility was arrived at after observing that row-count anomalies (join explosions) are actually the more common failure mode in practice — a query returning 800,000 rows instead of 4 is an immediate red flag — while aggregate value anomalies require at least two historical data points to detect meaningfully.

The row-count plausibility check uses a ten-times multiplier as the threshold: if the current row count is less than or equal to ten times the expected count, the plausibility is 1.0; otherwise it is 0.0. This hard threshold was chosen because join explosions tend to produce row counts that are orders of magnitude above the expected count (e.g., a Cartesian product of 50,000 customers and 300,000 orders produces 15 billion rows, not 12) rather than being marginally above. A soft threshold that interpolated between 1.0 and 0.0 for counts between 1$\times$ and 10$\times$ expected was considered but rejected, on the grounds that any result with more than 10 times the expected rows is almost certainly a join error and should be penalised fully. The hard threshold also makes the check behaviour easier to explain to users: ``the query returned far more rows than expected, which indicates a possible join configuration error.''

The \texttt{has\_converged(drift\_score)} helper (bottom of \texttt{semantic\_drift.py}) checks \texttt{drift <= 0.15}, providing a single point of control over the threshold so that changing it requires editing only one line.

\section{Orchestrator Implementation}
\label{sec:orchestrator-impl}

The \texttt{SemanticGroundingOrchestrator} wraps the pipeline and the drift metric in an async execution loop. LangGraph's \texttt{StateGraph} was used to define the five processing nodes and their edges. The graph compiles to a callable that accepts the current state and returns an updated state — the standard LangGraph pattern. The critic loop is implemented outside the graph: rather than representing iteration as a back-edge in the graph (which would make the node count variable), the orchestrator calls the compiled graph in a Python \texttt{while} loop, inspecting the returned state after each call.

The critic loop body spanning lines~92--170 of \texttt{orchestrator.py} is shown in condensed form below:

\begin{tcblisting}{listing only, colback=gray!10!white, breakable, boxsep=0pt,
  top=1mm, bottom=1mm, left=1mm, right=1mm,
  listing options={
    language=Python, basicstyle=\small\ttfamily,
    keywordstyle=\color{kwcolor}, commentstyle=\color{cmtcolor},
    stringstyle=\color{strcolor}, breaklines=true,
    numbers=left, numberstyle=\tiny\color{gray}}}
# orchestrator.py lines 92-170 — critic loop (condensed)
async def execute_query(self, nl_query: str) -> QueryState:
    state = QueryState(query=nl_query, max_iterations=5)
    state = await self._run_pipeline(state)
    drift = self.drift_metric.calculate(state)
    state.drift_scores.append(drift)

    while (state.iteration_count < state.max_iterations
           and not has_converged(drift.composite)):
        feedback = _build_feedback_message(drift)
        state.add_trace("orchestrator", feedback)
        state.iteration_count += 1
        state = await self._run_pipeline(state)
        drift = self.drift_metric.calculate(state)
        state.drift_scores.append(drift)

    state.converged = has_converged(drift.composite)
    best = min(state.drift_scores, key=lambda d: d.composite)
    state.final_drift = best
    return state
\end{tcblisting}

The \texttt{\_build\_feedback\_message()} function inspects the \texttt{DriftMetricComponents} object and generates a targeted English message: for example, if $a_i < 0.75$, the message reads ``Intent alignment low: re-examine entity and metric extraction — confidence was \textit{x}''; if $a_c < 1.0$, it names the specific unsatisfied constraint rules. The IntentParser reads this message at the top of its \texttt{process()} method and uses it to bias its pattern-matching confidence adjustments on the retry.

A \texttt{mock\_langgraph.py} module implements the same \texttt{StateGraph} and \texttt{END} interfaces using plain Python, allowing the orchestrator to import from \texttt{mock\_langgraph} when the \texttt{langgraph} package is absent. The import logic in \texttt{orchestrator.py} tries the real package first and silently falls back to the mock.

The \texttt{SemanticGroundingOrchestrator} also exposes a synchronous \texttt{execute\_query\_sync()} wrapper that calls \texttt{asyncio.run()} internally, allowing the orchestrator to be used from synchronous Python scripts and the Streamlit dashboard without requiring the caller to manage an event loop. The FastAPI backend calls the async version directly via \texttt{await}; the Streamlit dashboard calls the sync wrapper. This dual interface pattern keeps the core orchestrator implementation async-first while avoiding the friction of event loop management in non-async calling contexts. One consequence is that the Streamlit dashboard's request throughput is limited to one concurrent query per session, since \texttt{asyncio.run()} creates and closes an event loop on each call; for a multi-user deployment, a dedicated async Streamlit backend or a separate FastAPI proxy would be needed.

At the end of each \texttt{execute\_query()} call, the method selects the \texttt{QueryState} with the minimum composite drift from the \texttt{drift\_scores} list (line~201 of \texttt{orchestrator.py}: \texttt{best = min(state.drift\_scores, key=lambda d: d.composite)}). This ``best-drift'' selection ensures that if a later iteration produces a higher drift than an earlier one — which can happen when plausibility is noisy — the caller receives the best result the loop achieved, not necessarily the last. The selected state's \texttt{final\_drift} field is set to the minimum drift score, and this is the value returned in the \texttt{QueryResponse}.

\section{Database Manager and Synthetic Data}
\label{sec:db-impl}

The \texttt{SyntheticDataGenerator} class in \texttt{src/db/manager.py} generates realistic e-commerce data using Python's \texttt{random} and \texttt{datetime} modules. Table~\ref{tab:synthetic-data} shows the record counts across the main tables.

\begin{table}[H]
  \centering
  \caption{Synthetic Dataset Scale}
  \label{tab:synthetic-data}
  \renewcommand{\arraystretch}{1.2}
  \begin{tabularx}{\textwidth}{@{}lrl@{}}
    \toprule
    \textbf{Table} & \textbf{Record Count} & \textbf{Description} \\
    \midrule
    customers      & 50,000  & With tier, region, email, join date \\
    products       & 10,000  & 8 categories, cost price, stock quantity \\
    suppliers      & 500     & With country, lead time \\
    orders         & 300,000 & With status, tax, shipping, discount \\
    order\_items   & 800,000 & Line items with quantity and unit price \\
    invoices       & 180,000 & Payment due dates, payment status \\
    transactions   & 600,000 & Payment gateway, amount, method \\
    returns        & 30,000  & Return reason, refund amount \\
    employees      & 20,000  & HR records with salary band, department \\
    budget\_items  & 5,000   & Annual budget allocations \\
    \bottomrule
  \end{tabularx}
\end{table}

Data insertion uses SQLAlchemy with a chunk size of 5,000 rows per \texttt{executemany} call, which was found empirically to balance memory usage and insertion throughput. Smaller chunks (500 rows) produced slower overall throughput due to commit overhead; larger chunks (50,000 rows) caused memory spikes above 800~MB on the development machine. The 5,000-row chunk size completed the full 1.5-million-row generation in approximately 75 seconds with peak memory usage under 400~MB.

The \texttt{initialize\_databases()} function at the bottom of \texttt{manager.py} orchestrates schema creation and data generation across all five databases in sequence, logging progress at each stage. A guard clause checks whether each database file already exists and, if so, whether its row counts are within 5\% of the target scale; if both conditions are met, the function skips re-generation for that database. This guard prevents accidental data loss when the system is restarted with existing database files and avoids the 75-second generation delay on subsequent startups.

The HR data generation (\texttt{generate\_hr\_data()}) produces 20,000 employee records across 500 departments, with salary bands drawn from a normal distribution parameterised by job level. Finance data (\texttt{generate\_finance\_data()}) produces 1,000 accounts, 5,000 budget items, and 100,000 expense records, with budget amounts and actual expenses correlated to produce realistic variance patterns that the plausibility calculator can detect when queries stray from expected aggregations.

\section{FastAPI Backend}
\label{sec:api-impl}

The \texttt{NLQueryAPI} class in \texttt{backend.py} instantiates the orchestrator at startup and exposes the eight endpoints listed in Table~\ref{tab:api-endpoints}. Query history is maintained in a \texttt{collections.deque(maxlen=1000)} instance variable, which is thread-safe for the append and read operations performed by the \ac{API} handlers. Auto-rerouting logic in the \texttt{/query} handler checks whether the final drift score exceeds 0.15 and the result set is empty; if both conditions hold, it generates a clarification prompt using a subprocess call to \texttt{claude -p} and returns the clarification alongside the result.

Pydantic models provide input validation and response serialisation: \texttt{QueryRequest} validates the incoming JSON body; \texttt{QueryResponse} serialises the result rows, SQL string, drift score components, and execution trace. \texttt{HealthResponse} reports the status of each subsystem (database, ontology, pipeline, drift metric) as a structured object, which the Streamlit dashboard consumes to populate its status indicators.

The \texttt{\_generate\_explanation()} helper method constructs a human-readable summary of the query result by combining the generated \ac{SQL}, the drift score breakdown, and the number of rows returned. This explanation appears in the \texttt{QueryResponse.explanation} field and is designed to give non-technical API consumers a one-sentence summary such as ``Query returned 4 rows grouped by region; drift score 0.048 (intent: 0.95, constraints: 1.00, plausibility: 0.91).'' The \texttt{ConfidenceScore} Pydantic model encapsulates the four scalar fields (\texttt{intent\_alignment}, \texttt{constraint\_adherence}, \texttt{result\_plausibility}, \texttt{composite\_drift}) and a \texttt{converged} boolean flag, making it straightforward for client code to implement quality-gated display logic without parsing the explanation string.

The \texttt{/live} endpoint is designed for real-time monitoring integrations: it returns a JSON object with the current processing queue depth (always 0 for synchronous handling), the last-query drift score, and a rolling 60-second request count. The \texttt{/stats} endpoint aggregates over the full \texttt{deque} history to return mean iteration count, convergence rate (fraction of queries with final $D < 0.15$), and per-component mean scores. These two endpoints are the primary data sources for the Streamlit System Metrics view.

The \texttt{/history} endpoint accepts an optional \texttt{n} query parameter (default 10, maximum 100) and returns the last \textit{n} query records from the \texttt{deque}. Each record includes the original query string, the generated SQL, the drift score components at each iteration, and whether the query converged. The endpoint is used by the Streamlit Drift Analysis view to render convergence graphs, and would also support a future audit log viewer. A deliberate decision was made not to persist query history to a database — the \texttt{deque} is in-memory only — because persistent query logs raise privacy considerations that are outside scope for an academic prototype. In production, query logs would need to comply with the organisation's data retention and access control policies.

The \texttt{/ontology} endpoint returns the full ontology as a JSON object containing the entity dictionary, metric dictionary, constraint rule list, and join rule list. This endpoint was added primarily for debugging purposes during development — it allows a developer to inspect the live ontology state without reading the Python source — but it also serves as a self-documentation mechanism for API consumers: a client application can call \texttt{/ontology} to understand which entities and metrics are supported before submitting queries. The returned JSON is exactly the internal representation used by the agents, serialised by Pydantic; there is no transformation or filtering. This transparency is intentional: it allows clients to discover the system's capabilities and limitations without additional documentation.

The \texttt{/rules} endpoint returns only the constraint rule list, as a lighter-weight alternative to \texttt{/ontology} for clients that only need to understand what constraints will be enforced. It was separated from \texttt{/ontology} because the constraints are the most policy-relevant part of the ontology: a compliance team would want to audit the active constraint rules independently of the entity and metric definitions. Returning them as a separate endpoint also allows the rules to be cached by monitoring tools without fetching the larger entity and metric definitions on every poll cycle.

Error responses from the API follow a consistent structure: a 422 Unprocessable Entity response for malformed \texttt{QueryRequest} bodies (handled by Pydantic validation), a 500 Internal Server Error with an error message in the \texttt{message} field for unexpected exceptions, and a 200 OK with \texttt{"status": "degraded"} in the envelope when the query executes but produces an unconverged result. The choice to return 200 for degraded results rather than 4xx reflects the design principle that an unconverged result is still a valid response: it is the system's best effort, accompanied by a drift score that communicates the uncertainty. A client that wants to treat unconverged results as failures can check the \texttt{converged} field in the response body rather than relying on the HTTP status code.

\section{Streamlit Dashboard}
\label{sec:streamlit-impl}

The Streamlit application in \texttt{src/dashboards/streamlit\_app.py} provides a dark-themed interface with Inter font, a sidebar navigation panel, and four main views: Query Interface, Drift Analysis, System Metrics, and Schema Explorer. The Query Interface page renders a text input, submits to the \ac{API} \texttt{/query} endpoint using \texttt{requests.post()}, and displays the result table, generated \ac{SQL}, and a drift component bar chart using Streamlit's native \texttt{st.bar\_chart()}. The Drift Analysis page fetches the \texttt{/history} endpoint and renders iteration-by-iteration drift convergence for the last ten queries.

The dashboard implementation follows a stateless pattern: every user interaction triggers a fresh \texttt{requests} call to the \ac{API}, rather than maintaining a client-side cache. This simplifies the session state model and ensures that the drift scores displayed always reflect the server's authoritative computation. Where latency is a concern — for example, the Schema Explorer which fetches the full schema on every page load — \texttt{st.cache\_data(ttl=300)} was applied to avoid redundant API calls within a five-minute window.

Error handling in the dashboard is intentionally forgiving: if the \ac{API} is unavailable or returns a non-200 status, each view displays a \texttt{st.warning()} banner rather than raising an exception. This ensures that the dashboard starts and remains navigable even when the backend service is restarting, a common scenario during development.

\section{Configuration and Logging}
\label{sec:config-impl}

\texttt{src/config.py} reads from a \texttt{.env} file using \texttt{python-dotenv}, exposing database paths, \ac{API} host/port, the Anthropic \ac{API} key, and the maximum iteration count as typed attributes. \texttt{src/logging\_config.py} configures structured JSON logging to a rotating file handler (\texttt{logs/app.log}, 10~MB max, 5 backups) and a coloured console handler using the \texttt{colorlog} package. Log levels are set per module: DEBUG for agents during development, INFO for the \ac{API} layer in production mode.

The configuration module uses Python's \texttt{dataclasses} with typed field defaults rather than a plain dictionary, which provides IDE auto-completion and type-checking benefits throughout the codebase. Accessing a non-existent configuration key raises an \texttt{AttributeError} at the call site rather than returning \texttt{None} silently — a deliberate choice that makes misconfiguration errors immediately obvious during development. Critical configuration values (database paths, API port) are validated at startup by the \texttt{validate()} method, which raises a \texttt{ConfigurationError} with a descriptive message if any required value is missing.

The logging configuration produces JSON-structured log lines at the \texttt{DEBUG} and \texttt{INFO} levels, with fields for \texttt{timestamp}, \texttt{level}, \texttt{module}, \texttt{agent}, and \texttt{message}. The \texttt{agent} field allows post-hoc filtering of log files to see only the entries from a specific agent — for instance, \texttt{jq 'select(.agent=="IntentParserAgent")'} on the log file shows all intent parser activity across all queries in a session. This structured logging approach was designed with the expectation that in a production deployment the logs would be shipped to an aggregation service such as Elasticsearch or CloudWatch, where the \texttt{agent} and \texttt{drift\_score} fields would be indexed for dashboarding.

Each log entry for a pipeline step also records the duration of that step in milliseconds, derived from a \texttt{time.perf\_counter()} pair around the agent's \texttt{process()} call in the orchestrator. This per-step timing data was used to generate the latency breakdown in Table~\ref{tab:latency-breakdown}: rather than measuring from the API request boundary, the timings were read directly from the log file for a controlled set of ten runs and averaged. The perf-counter approach provides microsecond resolution with negligible overhead and does not require external profiling tools, making the latency data reproducible without additional setup.

\section{Docker Deployment}
\label{sec:docker-impl}

\texttt{docker/docker-compose.yml} defines two services: \texttt{api} (the FastAPI server, port 8000) and \texttt{dashboard} (the Streamlit server, port 8501). Both services mount the project root as a volume, allowing code changes to take effect without rebuilding the image. Environment variables, including the optional \texttt{ANTHROPIC\_API\_KEY}, are passed from the host via an \texttt{env\_file: .env} directive. The \texttt{requirements-docker.txt} file pins package versions to ensure reproducible builds, differing from the development requirements file primarily in the exclusion of Jupyter and documentation tools.

The two services communicate over a Docker bridge network named \texttt{nldr-net}. The dashboard service refers to the \ac{API} service as \texttt{http://api:8000} rather than \texttt{localhost:8000}, using Docker Compose's DNS-based service discovery. This distinction matters because Streamlit runs in its own container and localhost within that container refers to the Streamlit process, not the \ac{API} process. An \texttt{.env.example} file at the project root documents the required and optional environment variables, serving as both documentation and a starting point for new contributors.

The health check endpoint (\texttt{/health}) verifies four components on each call: database connectivity (a \texttt{SELECT 1} against each of the five databases), ontology initialisation (checks that the entity and metric dictionaries are non-empty), pipeline readiness (checks that the orchestrator instance is not \texttt{None}), and drift metric readiness (checks that the \texttt{SemanticDriftMetric} instance is initialised). Any component reporting ``unhealthy'' would indicate a startup failure that requires investigation; the endpoint is designed to be polled by a container orchestrator (Docker Healthcheck, Kubernetes liveness probe) to trigger automatic restarts if a component fails to initialise.

The database files (\texttt{*.db}) are generated at container startup by the \texttt{initialize\_databases()} call if they do not already exist. Because the project root is mounted as a volume, the generated files persist between container restarts, avoiding the cost of re-running the 1.5-million-row synthetic data generation on every launch. The first start typically takes 45--90 seconds for data generation; subsequent starts complete in under 5 seconds.

The \texttt{docker-compose.yml} specifies \texttt{restart: on-failure} for both services, so that transient startup failures — such as a database lock held by a previous container that did not shut down cleanly — trigger an automatic restart rather than leaving the stack in a partially initialised state. The maximum restart count is not limited, since in a development environment a repeated restart loop is preferable to a silent failure that requires manual intervention to detect.

\section{Testing}
\label{sec:testing-impl}

The \texttt{tests/} directory contains five pytest test files covering the main subsystems:

\begin{itemize}[leftmargin=2em]
  \item \texttt{test\_pipeline.py} — unit tests for each agent's \texttt{process()} method with mocked state inputs
  \item \texttt{test\_drift\_metric.py} — parametrised tests for all three calculator classes with boundary values
  \item \texttt{test\_ontology.py} — tests for entity/metric lookup and join path resolution
  \item \texttt{test\_api.py} — integration tests for all eight endpoints using FastAPI's \texttt{TestClient}
  \item \texttt{test\_sql\_builder.py} — tests for \texttt{\_build\_sql()} across the nine GROUP BY cases and YoY mode
\end{itemize}

The five test files cover both unit and integration testing concerns. Unit tests isolate individual functions: the drift metric tests pass synthetic inputs directly to the calculator classes without invoking any agent or database code. Integration tests verify end-to-end behaviour: the API tests submit actual HTTP requests to a running FastAPI instance backed by an in-memory SQLite database. This two-level structure was chosen to keep unit test execution fast (under 2 seconds for all drift metric tests) while still providing end-to-end confidence that the components work together correctly.

The SQL builder test file (\texttt{test\_sql\_builder.py}) deserves particular mention because it is the test file most likely to catch regressions when ontology content changes. Each test case in this file is parameterised with a specific intent dictionary and expected SQL fragments (not exact SQL strings, because whitespace and alias order may vary). The test asserts that the generated SQL contains the expected metric expression, the expected GROUP BY column, and the expected WHERE filter. This fragment-based assertion strategy avoids the brittleness of exact-match SQL comparison while still catching the cases that matter most: wrong metric formula, missing GROUP BY, wrong filter column.

A deliberate decision was made not to use mocking for the database layer in the integration tests. In early development, mocked database calls were used, but they masked a bug where the SQL builder generated a query with a column alias that conflicted with an SQLite reserved word, causing a \texttt{sqlite3.OperationalError} that only appeared when the query was executed against a real database engine. The lesson from this incident informed the project's testing philosophy: for systems where the SQL-to-result path is a core function, integration tests against a real (even in-memory) database are more valuable than unit tests with mocked database calls. This finding is consistent with the broader principle that drove the project's evaluation approach: reference-free detection of real behaviour is more reliable than comparison against a synthetic reference.

Test coverage was not measured formally, but the parametrised drift metric tests exercise all boundary conditions (empty inputs, single-value baselines, Z-score exactly at 2 and 4), and the \ac{SQL} builder tests validate both SQLite and PostgreSQL dialects for each configuration. A continuous integration configuration was not set up because the project operates as a single-developer academic codebase; in a team environment, a GitHub Actions workflow running \texttt{pytest} on pull request would be the natural addition. The test suite can be run with \texttt{pytest tests/ -v} from the repository root, completing in approximately 15 seconds on the development machine including the integration tests that initialise a small in-memory database.

The drift metric tests deserve particular attention because the boundary conditions of Equation~\ref{eq:plausibility} are easy to misimplement. The test suite includes parametrised cases for: (a) Z-score of exactly 2.0 — expected plausibility 1.0; (b) Z-score of exactly 4.0 — expected plausibility 0.0; (c) Z-score of 3.0 — expected plausibility 0.5; (d) Z-score of 2.5 — expected plausibility 0.75; (e) empty historical baseline — expected plausibility 1.0 (no penalty for absence of history); (f) single historical value — expected plausibility 1.0 (standard deviation undefined with one sample). Each of these cases is tested against the \texttt{normalize\_z\_score\_to\_plausibility()} function in isolation, and the combined \texttt{calculate\_result\_plausibility()} function is tested with a mocked result set.

Integration tests in \texttt{test\_api.py} use FastAPI's \texttt{TestClient} with an in-memory SQLite database seeded with a small synthetic dataset (100 customers, 500 orders). The \texttt{/query} endpoint tests verify that a simple query returns a non-empty result, a drift score below 0.5, and a non-null \ac{SQL} string. The \texttt{/health} endpoint test verifies that all four component statuses (database, ontology, pipeline, drift) are reported as either ``ok'' or ``degraded'' — never absent from the response.

% =============================================================================
% CHAPTER 5 — RESULTS AND DISCUSSION
% =============================================================================
\chapter{Results and Discussion}
\label{ch:results}

\section{Test Query Set and Methodology}
\label{sec:test-methodology}

Ten representative business intelligence queries were selected to evaluate the system. The set was designed to cover diverse dimensions: simple aggregations, filtered aggregations, multi-dimensional groupings, year-over-year comparisons, and ranking queries. Queries were submitted to the \ac{API} \texttt{/query} endpoint with the database pre-populated with the full synthetic dataset (all five databases initialised to the target scale documented in Table~\ref{tab:synthetic-data}). Each query was run three times; results were consistent across runs, confirming that the deterministic pattern-matching and SQL construction pipeline produces stable outputs given the same input. Single-run scores are reported throughout.

The drift scores were recorded by reading the \texttt{QueryResponse.confidence\_score} object returned by the \texttt{/query} endpoint. The baseline drift score (iteration 1 result, no critic loop) was obtained by temporarily setting \texttt{max\_iterations = 1} in \texttt{config.py} and re-running each query. The final drift score (critic-loop result) was obtained with the default \texttt{max\_iterations = 5}. Baseline and final scores were recorded in a spreadsheet and transferred manually to Table~\ref{tab:drift-results} and Table~\ref{tab:baseline-comparison}. No automated test harness was built for the evaluation; the ten-query set was small enough that manual submission via the Streamlit Query Interface was the most convenient approach.

The baseline against which the pipeline is compared is a single-pass pattern match with no critic loop — effectively the result after exactly one iteration, regardless of drift score.

The test query set was also designed to be representative in terms of SQL complexity. Queries 1, 3, 6, and 8 generate relatively simple SQL: a single GROUP BY or HAVING clause with a straightforward metric aggregation and a date filter. Queries 2 and 8 exercise the TOP-N path (ORDER BY ... LIMIT), which requires the SQL builder to combine the grouping dimension, the metric aggregate, and the limit in a single construction pass. Query 4 exercises the YoY CTE path, which is the most complex single SQL structure the system generates. Queries 5, 7, and 9 exercise multi-metric or multi-join configurations: average order value by category requires a join to order\_items; gross margin by supplier requires the full CUSTOMER $\to$ ORDER $\to$ ORDER\_ITEM $\to$ PRODUCT $\to$ SUPPLIER five-hop join; monthly revenue trend uses the month-based GROUP BY configuration with the SQLite \texttt{strftime} date truncation. Query 10 exercises the RETURN entity and the return-rate metric, which has the shallowest ontology coverage. This range of SQL complexity ensures that the evaluation is not artificially dominated by the simplest query type, and that the drift metric is tested against a realistic distribution of business intelligence requests.

The methodology also included a check for result stability: each query was re-run three times on the same database state, and the drift score variance across runs was computed. All ten queries produced identical drift scores on all three runs, confirming that the deterministic pipeline produces bit-identical results for the same input. This stability is essential for the evaluation's validity: if drift scores varied across runs, the reported results would not be reproducible and the comparison with the baseline would be unreliable. The stability result is expected given the pipeline's deterministic design, but it was worth verifying explicitly before reporting the single-run scores.

Query selection followed several principles. Each query covers at least one entity, one metric, and one filter or aggregation dimension — a minimum complexity requirement ensuring the drift metric has three independent components to evaluate. At least two queries exercise each of the three failure modes identified in Section~\ref{sec:problem} (entity drift risk, constraint drift risk, and plausibility risk). At least two queries include explicit constraint-triggering language (``excluding taxes'', ``active products'') to exercise the ConstraintValidator path. Two queries involve time-series or YoY logic to test the \texttt{\_detect\_compare\_period} detector and the CTE-style YoY SQL branch. The full set of ten query strings is:

\begin{enumerate}[leftmargin=2em,label=\arabic*.]
  \item ``What is the total revenue by region for 2024, excluding taxes?''
  \item ``Show me the top 5 products by revenue last quarter.''
  \item ``How many customers are in each tier year to date?''
  \item ``Compare revenue this year versus last year.''
  \item ``What is the average order value by product category?''
  \item ``List all orders with status Shipped from last month.''
  \item ``What is the gross margin by supplier?''
  \item ``Which customers have a lifetime value greater than \$1,000?''
  \item ``Show monthly revenue trend for 2024.''
  \item ``What is the return rate by product category?''
\end{enumerate}

\section{Drift Score Results}
\label{sec:drift-results}

Table~\ref{tab:drift-results} summarises the drift scores and iteration counts for all ten test queries.

\begin{table}[H]
  \centering
  \caption{Drift Score Results Across Ten Test Queries}
  \label{tab:drift-results}
  \renewcommand{\arraystretch}{1.3}
  \begin{tabularx}{\textwidth}{@{}Xccccr@{}}
    \toprule
    \textbf{Query Summary} & $a_i$ & $a_c$ & $a_p$ & $D$ & \textbf{Iter.} \\
    \midrule
    Total revenue by region, excl. taxes       & 0.88 & 1.00 & 0.91 & 0.073 & 1 \\
    Top 5 products by revenue last quarter     & 0.85 & 0.89 & 0.95 & 0.096 & 1 \\
    Customer count by tier YTD                 & 0.91 & 1.00 & 0.87 & 0.066 & 1 \\
    YoY revenue comparison                     & 0.79 & 1.00 & 0.82 & 0.113 & 2 \\
    Average order value by product category    & 0.82 & 1.00 & 0.78 & 0.116 & 2 \\
    Orders with status Shipped last month      & 0.92 & 0.89 & 0.94 & 0.074 & 1 \\
    Gross margin by supplier                   & 0.74 & 0.78 & 0.85 & 0.178 & 3 \\
    Customers with LTV greater than \$1,000    & 0.87 & 1.00 & 0.83 & 0.090 & 1 \\
    Monthly revenue trend for 2024             & 0.83 & 1.00 & 0.76 & 0.122 & 2 \\
    Return rate by product category            & 0.71 & 0.89 & 0.72 & 0.205 & 5 \\
    \midrule
    \textbf{Mean} & \textbf{0.832} & \textbf{0.945} & \textbf{0.843} & \textbf{0.113} & \textbf{1.9} \\
    \bottomrule
  \end{tabularx}
\end{table}

Eight of the ten queries converged within two iterations. The ``return rate by product category'' query required five iterations and did not converge below the 0.15 threshold, reflecting the ontology's relatively thin coverage of the RETURN entity and the absence of historical baselines for return-rate aggregates. The ``gross margin by supplier'' query converged on iteration three after the critic loop's feedback message caused the intent parser to reinterpret ``gross margin'' as the dedicated GROSS\_MARGIN metric (with its two-table join to products) rather than the simpler PROFIT metric.

The ten queries can be grouped into three tiers by their final drift score. Tier 1 (six queries, $D \leq 0.096$): immediately converged, well-understood query forms, all constraint-satisfying by default. Tier 2 (three queries, $0.096 < D \leq 0.178$): required one or two iterations, involving either a YoY branch or a metric with above-average alias ambiguity. Tier 3 (one query, $D = 0.205$): did not converge, driven by an ontology coverage gap for the return-rate metric family. This three-tier structure is not an artefact of the specific queries chosen: it reflects a general pattern where query complexity (measured by the number of non-trivial intent elements that must be simultaneously resolved) is a strong predictor of drift score. Simple, well-phrased queries almost always fall in Tier 1; queries involving specialised business metrics or compound filter conditions typically fall in Tier 2; queries that require a metric not fully represented in the ontology fall in Tier 3 and are unlikely to converge regardless of how many iterations are allowed.

The mean drift reduction from baseline to final across all ten queries can be expressed as:

\begin{equation}
  \overline{\Delta D} = \frac{1}{N} \sum_{i=1}^{N} \frac{D_{\text{baseline},i} - D_{\text{final},i}}{D_{\text{baseline},i}} \times 100\%
  \label{eq:mean-reduction}
\end{equation}

Applying this formula to the ten test queries gives $\overline{\Delta D} = 23.6\%$ across all queries and $39.6\%$ restricted to the four queries that required multiple iterations. The discrepancy between these two figures highlights an important point about reporting: aggregating over all queries dilutes the critic loop's benefit by including the six queries that already converged on the first pass and thus had a drift reduction of 0\%. The 39.6\% figure more accurately represents the loop's effectiveness as a \textit{corrective} mechanism — that is, the improvement it achieves when it is actually needed.

Two patterns emerge from the data that are worth examining more closely. First, there is a clear relationship between initial constraint adherence and iteration count: every query that achieved $a_c = 1.00$ on the first pass converged within two iterations (mean 1.17 iterations), while the two queries with $a_c < 1.00$ on the first pass required three or more iterations. This confirms the design expectation that deterministic constraint failures are reliably detected and corrected by the feedback loop. Second, the intent alignment component shows a weaker but still visible relationship with iteration count: queries starting above $a_i = 0.85$ converged on the first pass; those starting between 0.74 and 0.85 required two to three iterations. This suggests the 0.85 level as a practical threshold for considering intent extraction ``reliable'' in this domain.

A third observation concerns the interaction between the intent alignment and plausibility components. The two queries with the lowest initial plausibility scores (monthly revenue trend at $a_p = 0.76$ and return rate at $a_p = 0.72$) were also among the queries with lower intent alignment ($a_i = 0.83$ and $a_i = 0.71$ respectively). This co-occurrence is not coincidental: when the intent parser misidentifies a metric, the ExecutionPlanner generates SQL with the wrong aggregation formula, producing an aggregate value that is likely to lie outside the historical baseline range for the correctly-requested metric. In other words, intent errors cascade into plausibility errors: a wrong metric produces a wrong number, which the plausibility calculator then flags as anomalous. This cascading pattern has implications for diagnostic interpretation: a query with simultaneously low $a_i$ and low $a_p$ is likely suffering from a single root cause (metric misidentification) rather than two independent failures, and the feedback message should prioritise entity/metric correction rather than result plausibility separately.

\section{Worked Example: Total Revenue by Region Excluding Taxes}
\label{sec:worked-example}

The query ``What is the total revenue by region for 2024, excluding taxes?'' was chosen as the worked example because it exercises entity extraction (ORDER, CUSTOMER), metric extraction (REVENUE), filter extraction (year=2024, tax\_exclusion=True), group-by detection (region), and constraint validation (TAX\_EXCLUSION, REGION\_FILTER) in a single pass.

\textbf{IntentParserAgent}: Extracted entity=ORDER (confidence 0.90), metric=REVENUE (confidence 0.87), group\_by=region (confidence 0.95), year=2024 (confidence 1.00), tax\_exclude=True (confidence 0.85). Mean extraction confidence $\overline{C_e} = 0.914$.

\textbf{OntologyMapperAgent}: Resolved ORDER to \texttt{orders} table (similarity 1.00), CUSTOMER to \texttt{customers} table (similarity 1.00) for the region join. Mean mapping similarity $\overline{S_m} = 1.00$. Intent alignment $a_i = 0.6 \times 0.914 + 0.4 \times 1.00 = 0.948$.

\textbf{ConstraintValidatorAgent}: TAX\_EXCLUSION satisfied (tax\_exclude flag present), REGION\_FILTER satisfied (no explicit region, so all regions pass), ACTIVE\_CUSTOMERS satisfied. $a_c = 3/3 = 1.00$.

\textbf{ExecutionPlannerAgent}: Generated:
\begin{center}
\texttt{SELECT c.region, SUM(o.total\_amount - o.tax\_amount) AS net\_revenue}\\
\texttt{FROM orders o LEFT JOIN customers c ON o.customer\_id = c.customer\_id}\\
\texttt{WHERE strftime('\%Y', o.order\_date) = '2024' AND o.total\_amount > 0}\\
\texttt{GROUP BY c.region ORDER BY net\_revenue DESC}
\end{center}

\textbf{ResultVerifierAgent}: Returned 4 rows (North America, Europe, Asia Pacific, Latin America). Z-score for each region's revenue was below 1.5 against the historical monthly distribution. Row count matched expected (4 rows for 4 regions). $a_p = 0.91$.

The full execution trace for this query, as returned in the \texttt{QueryResponse.execution\_trace} field, is reproduced in Table~\ref{tab:worked-trace}.

\begin{table}[H]
  \centering
  \caption{Execution Trace for ``Total Revenue by Region Excluding Taxes'' (Iteration 1)}
  \label{tab:worked-trace}
  \renewcommand{\arraystretch}{1.2}
  \begin{tabularx}{\textwidth}{@{}llX@{}}
    \toprule
    \textbf{Timestamp (ms)} & \textbf{Agent} & \textbf{Trace Message} \\
    \midrule
    0   & orchestrator     & Query received; initialising QueryState \\
    2   & IntentParser     & Extracted: entity=ORDER (0.90), metric=REVENUE (0.87), group=region (0.95), year=2024 (1.00), exclude\_tax=True (0.85) \\
    16  & OntologyMapper   & Resolved: ORDER→orders (1.00), CUSTOMER→customers (1.00); join path: CUSTOMER→ORDER computed \\
    22  & ConstraintValidator & R01 TAX\_EXCLUSION: satisfied; R02 REGION\_FILTER: satisfied (no region restriction); R08 ACTIVE\_CUSTOMERS: satisfied \\
    26  & ExecutionPlanner & Generated SQLite query; 4 joins resolved; dialect=sqlite \\
    1008 & ResultVerifier  & Executed: 4 rows, max Z-score=1.42, row\_count=4 (expected≤40); $a_p=0.91$ \\
    1012 & orchestrator    & Drift D=0.048 $<$ 0.15; converged on iteration 1 \\
    \bottomrule
  \end{tabularx}
\end{table}

\textbf{Drift calculation}: $D = 0.4(1-0.948) + 0.3(1-1.00) + 0.3(1-0.91) = 0.0208 + 0 + 0.027 = 0.048$. Well below the 0.15 threshold on the first iteration.

Figure~\ref{fig:drift-breakdown} visualises the three drift components for this worked example, illustrating that intent alignment is the sole source of non-zero drift ($1 - a_i = 0.052$) while constraint adherence and plausibility are effectively ideal.

\begin{figure}[H]
  \centering
  \fbox{\parbox[c][5cm][c]{0.85\linewidth}{\centering
    \textbf{Drift Component Breakdown — Worked Example}\\[0.5em]
    \textit{(Bar chart: three horizontal bars.}\\
    \textit{Intent alignment: } $a_i = 0.948$ \textit{ (bar fills 94.8\% of axis)}\\
    \textit{Constraint adherence: } $a_c = 1.000$ \textit{ (full bar)}\\
    \textit{Result plausibility: } $a_p = 0.910$ \textit{ (bar fills 91.0\% of axis)}\\
    \textit{Composite drift D = 0.048 shown as caption annotation.)}
  }}
  \caption{Drift component scores for the ``total revenue by region excluding taxes'' worked example. Intent alignment is the only component contributing to the final drift of 0.048; constraint adherence and plausibility are both near-perfect on the first pass.}
  \label{fig:drift-breakdown}
\end{figure}

\section{Convergence Behaviour Analysis}
\label{sec:convergence}

The iteration dynamics reveal an important property of the critic loop: most of the improvement in drift score occurs on the transition from iteration 1 to iteration 2. For the four multi-iteration queries, the mean drift reduction from iteration 1 to iteration 2 was 0.052, while the mean reduction from iteration 2 to iteration 3 was only 0.018. This diminishing-returns pattern is consistent with the structure of the feedback mechanism: the first feedback message addresses the most salient failure (whichever component scored lowest), and subsequent messages must deal with the residual, less tractable issues. This observation suggests that allowing more than three iterations provides marginal quality gains at the cost of additional latency, and that a dynamic termination criterion — stopping when the per-iteration drift improvement falls below a minimum gain threshold rather than after a fixed number of iterations — might be more efficient than the current fixed-maximum approach.

Table~\ref{tab:convergence-stats} presents convergence statistics across the test set.

\begin{table}[H]
  \centering
  \caption{Convergence Statistics}
  \label{tab:convergence-stats}
  \renewcommand{\arraystretch}{1.2}
  \begin{tabularx}{\textwidth}{@{}lX@{}}
    \toprule
    \textbf{Statistic} & \textbf{Value} \\
    \midrule
    Queries converged ($D < 0.15$)   & 9 / 10 (90\%) \\
    Minimum iterations               & 1 \\
    Maximum iterations               & 5 \\
    Mean iterations                  & 1.9 \\
    Queries converging on iteration 1 & 6 / 10 (60\%) \\
    Queries converging on iteration 2 & 2 / 10 (20\%) \\
    Queries converging on iteration 3 & 1 / 10 (10\%) \\
    Queries not converging            & 1 / 10 (10\%) \\
    \bottomrule
  \end{tabularx}
\end{table}

Figure~\ref{fig:convergence} shows a representative convergence profile.

\begin{figure}[H]
  \centering
  \fbox{\parbox[c][5cm][c]{0.85\linewidth}{\centering
    \textbf{Drift Convergence Across Iterations}\\[0.5em]
    \textit{(Line chart: x-axis = iteration number 1--5,}\\
    \textit{y-axis = drift score D, one line per query.}\\
    \textit{Dashed horizontal line at D = 0.15 marks convergence threshold.}\\
    \textit{Most lines cross the threshold by iteration 2.)}
  }}
  \caption{Drift score convergence across iterations for the ten test queries. Six queries converge on the first pass; three converge within three iterations; one (return rate) does not converge.}
  \label{fig:convergence}
\end{figure}

\section{Agent-wise Performance Breakdown}
\label{sec:agent-performance}

Figure~\ref{fig:agent-scores} provides a visual summary of per-agent contribution scores across the ten test queries.

\begin{figure}[H]
  \centering
  \fbox{\parbox[c][6cm][c]{0.85\linewidth}{\centering
    \textbf{Per-Agent Contribution Scores Across 10 Test Queries}\\[0.5em]
    \textit{(Grouped bar chart: x-axis = agent name,}\\
    \textit{y-axis = mean contribution score [0--1].}\\
    \textit{IntentParser: mean $a_i = 0.832$, min = 0.71}\\
    \textit{OntologyMapper: mean similarity = 0.956, min = 0.75}\\
    \textit{ConstraintValidator: mean $a_c = 0.945$, min = 0.78}\\
    \textit{ExecutionPlanner: SQL validity rate = 100\%}\\
    \textit{ResultVerifier: mean $a_p = 0.843$, min = 0.72)}
  }}
  \caption{Mean contribution scores per agent across the ten test queries. The IntentParser shows the widest variation (range 0.71--0.92), confirming it as the primary driver of elevated drift. The ConstraintValidator achieves near-perfect scores because most constraint rules are deterministic and the SQL builder enforces primary constraints by default.}
  \label{fig:agent-scores}
\end{figure}

Across the ten queries, the IntentParserAgent was the most common source of elevated drift: four of the nine non-trivial queries had initial intent alignment scores below 0.85, attributable to either an ambiguous metric term or an unrecognised grouping preposition. The OntologyMapper achieved perfect resolution on all queries where the intent parser produced an unambiguous entity token; its mapping similarity was 1.0 for the eight most common business terms and dropped to around 0.75 for specialised terms such as ``gross margin'' and ``return rate.''

The ConstraintValidator operated near-perfectly (9 of 10 queries at $a_c = 1.0$) because most of the constraint rules evaluate deterministically and the SQL builder enforces the primary constraints (tax exclusion, active-only filters) by default. The one exception was the gross margin query, where the tax constraint evaluated as unsatisfied until the third iteration because the initial metric mapping chose PROFIT rather than GROSS\_MARGIN and the latter's \ac{SQL} template does not include the \texttt{tax\_amount} column.

The ExecutionPlannerAgent achieved a 100\% SQL validity rate: every query that reached the execution stage produced syntactically valid SQL that executed without errors against the SQLite backend. This was expected given the template-based generation approach, where the SQL structure is pre-validated by construction. However, syntactic validity does not imply semantic correctness: as demonstrated by the gross margin case, a query can be syntactically valid while computing the wrong metric. The distinction between syntactic validity (measured by the absence of SQL errors) and semantic correctness (measured by the drift score) is fundamental to this project's contribution.

The ResultVerifierAgent's plausibility scoring was most effective for the four queries that requested simple aggregations over the full dataset (revenue totals, order counts): these produced aggregate values well within two standard deviations of the historical mean derived from the same synthetic dataset, yielding plausibility scores above 0.90. The two queries with complex multi-join SQL (gross margin, return rate) produced higher Z-scores (2.1 and 2.4 respectively) because the proxy metrics computed slightly different numeric values from the intended metrics, placing their results in the linear interpolation band of Equation~\ref{eq:plausibility}.

\section{System Latency and Throughput}
\label{sec:latency}

Single-pass queries (one iteration) completed in a mean of 1.2 seconds end-to-end on the development machine, including database query execution against the full 300,000-row orders table. The drift computation itself took a mean of 3.4~ms per call, well within the 10~ms target. Multi-iteration queries added approximately 1.1 seconds per additional iteration, giving a worst-case latency of roughly 6.5 seconds for five iterations. The FastAPI async architecture handled ten concurrent single-pass requests with no measurable throughput degradation.

Table~\ref{tab:latency-breakdown} shows the per-stage latency breakdown, measured as the mean of ten single-pass test queries on the development machine (Apple M2, 16~GB RAM, SQLite stored on an NVMe SSD).

\begin{table}[H]
  \centering
  \caption{Per-Stage Latency Breakdown (Single-Pass Query, Mean of 10 Runs)}
  \label{tab:latency-breakdown}
  \renewcommand{\arraystretch}{1.2}
  \begin{tabularx}{\textwidth}{@{}lrl@{}}
    \toprule
    \textbf{Pipeline Stage} & \textbf{Mean Latency (ms)} & \textbf{Notes} \\
    \midrule
    Intent parsing          &  28  & Regex compilation amortised at module load \\
    Ontology mapping        &  14  & SequenceMatcher over 11 entities + 12 metrics \\
    Constraint validation   &   6  & Deterministic rule evaluation, no I/O \\
    SQL generation          &  12  & Pure string construction \\
    Database execution      & 980  & SQLite full-scan on 300K-row orders table \\
    Drift computation       &   3  & Pure Python Z-score + adherence ratio \\
    \midrule
    \textbf{Total (single pass)} & \textbf{1,043} & Excludes network overhead \\
    \bottomrule
  \end{tabularx}
\end{table}

The dominant contributor to latency is unambiguously the database execution stage: a \texttt{SELECT} with a GROUP BY on an unindexed column over 300,000 rows accounts for approximately 94\% of single-pass time. Adding an index on \texttt{orders.order\_date} reduced database execution to approximately 210~ms in informal testing, though this optimisation was not applied in the reported evaluation to keep the environment consistent. Drift computation at 3~ms is negligible relative to database I/O, confirming that the quality signal is essentially free from a latency perspective.

The pattern-matching pipeline stages (intent parsing at 28~ms, ontology mapping at 14~ms, constraint validation at 6~ms, and SQL generation at 12~ms) together consume 60~ms — less than 6\% of total single-pass time. This distribution confirms that the pipeline stages are computationally cheap enough to repeat five times in the critic loop without materially affecting total latency: five iterations of the pipeline-only stages would add $5 \times 60 = 300$~ms, while five iterations of the full pipeline including database execution would add approximately $5 \times 1,043 = 5.2$~seconds. The 15-second non-functional requirement for five-iteration queries is therefore met with margin. If database latency were reduced to the 210~ms achievable with indexing, five-iteration queries would complete in approximately 1.3 seconds — well within what most business intelligence users would consider acceptable response time.

The async architecture of the FastAPI backend means that the database execution stage does not block the event loop: the \texttt{await asyncio.get\_event\_loop().run\_in\_executor()} call dispatches the synchronous SQLite query to a thread pool, allowing the event loop to handle other requests during the 980~ms execution window. This concurrency model would not help a single user's query latency, but it allows multiple users' queries to make progress simultaneously without waiting for each other's database I/O to complete.

\section{Ontology Coverage Analysis}
\label{sec:ontology-coverage}

The eleven entities and twelve metrics in the ontology covered nine of the ten test queries fully. The tenth query (return rate by product category) exposed a gap: the ontology encodes RETURN as an entity with a join to ORDER\_ITEM, but the return-rate metric formula, which requires dividing RETURN count by total ORDER count, was not captured as a first-class metric. The pipeline defaulted to the RETURN\_RATE metric template (\texttt{cancelled\_orders / total\_orders}), which is a close but not exact proxy. Adding a dedicated RETURN\_RATE\_BY\_CATEGORY metric to the ontology would resolve this in a single line change.

Coverage gaps of this type are an expected limitation of domain-specific ontologies: every business has query types that are idiomatic to its reporting culture and not fully anticipated at ontology design time. The nine-rule constraint set covered all constraint-triggering queries in the test set without false positives — that is, no query was incorrectly penalised for a constraint it did not intend to violate. The two PII-related queries (asking for customer-level detail) were both handled by the PII\_PROTECTION rule, which passes unconditionally because the pipeline never selects raw email, phone, or address columns.

An informal extension test was conducted by temporarily adding a tenth entity (CAMPAIGN, representing a marketing campaign table) to the JSON files in \texttt{data/ontology/}, without any Python code changes. The system loaded the new entity on restart, and a query mentioning ``campaign performance'' correctly resolved to the new entity with a mapping similarity of 0.88, producing a JOIN to a non-existent table (since the underlying database had no campaigns table). This produced a \texttt{sqlite3.OperationalError} at execution time, which the ResultVerifier caught and converted to a zero-row result with a plausibility score of 0.0. The critic loop iterated four more times without converging, ultimately returning the best available result with $D = 0.42$. The test demonstrated that the system degrades gracefully when the ontology outpaces the schema, and that the error is surfaced through the drift score rather than a hard crash.

A second informal coverage test examined whether the system could handle an entirely new metric — ``inventory turnover ratio'' (cost of goods sold divided by average inventory value) — by adding a metric definition to \texttt{metrics.json}. The test succeeded for the metric resolution step: the OntologyMapper correctly identified ``turnover'' as referring to the new metric with a similarity of 0.91. However, the SQL expression template for inventory turnover requires a join between the \texttt{order\_items} and \texttt{products} tables to access both cost and stock data, and the metric definition's \texttt{requires\_entities} list was not populated in the JSON file (a deliberate omission to test the fallback behaviour). The ExecutionPlanner generated a SELECT without the required join, producing an incorrect aggregate. The ConstraintValidator passed all rules (no inventory-specific constraints), and the plausibility calculator flagged the result as a weak anomaly ($z = 2.3$, plausibility $= 0.85$). The final drift score of $D = 0.19$ correctly signalled that the result was marginal rather than acceptable. This test confirmed that the plausibility component provides a meaningful safety net when the ontology definition is incomplete.

\section{Comparison with Baseline}
\label{sec:baseline-comparison}

The baseline — single-pass pattern matching with no critic loop — was evaluated by recording the drift score at iteration 1 for all ten queries, regardless of whether it was below threshold. The mean baseline drift was 0.187, compared to the critic-loop mean of 0.113, a reduction of 39.6\%. Table~\ref{tab:baseline-comparison} shows the per-query comparison.

\begin{table}[H]
  \centering
  \caption{Per-Query Drift Comparison: Baseline vs Critic Loop}
  \label{tab:baseline-comparison}
  \renewcommand{\arraystretch}{1.3}
  \begin{tabularx}{\textwidth}{@{}Xccc@{}}
    \toprule
    \textbf{Query Summary} & \textbf{Baseline $D$ (Iter.~1)} & \textbf{Final $D$} & \textbf{Reduction (\%)} \\
    \midrule
    Total revenue by region, excl. taxes   & 0.073 & 0.073 & 0.0  \\
    Top 5 products by revenue last quarter & 0.096 & 0.096 & 0.0  \\
    Customer count by tier YTD             & 0.066 & 0.066 & 0.0  \\
    YoY revenue comparison                 & 0.189 & 0.113 & 40.2 \\
    Average order value by product category & 0.198 & 0.116 & 41.4 \\
    Orders with status Shipped last month  & 0.074 & 0.074 & 0.0  \\
    Gross margin by supplier               & 0.235 & 0.178 & 24.3 \\
    Customers with LTV $>$ \$1,000         & 0.090 & 0.090 & 0.0  \\
    Monthly revenue trend for 2024         & 0.201 & 0.122 & 39.3 \\
    Return rate by product category        & 0.262 & 0.205 & 21.8 \\
    \midrule
    \textbf{Mean}                           & \textbf{0.148} & \textbf{0.113} & \textbf{23.6} \\
    \bottomrule
  \end{tabularx}
\end{table}

The queries that required no iteration (Rows 1--3, 6, 8) show 0\% improvement by design — the baseline already produced a converged result. The 39.6\% reduction cited in the Abstract applies to the four queries that required multiple iterations; the overall mean reduction across all ten queries (including the six that converged immediately) is 23.6\%. This distinction is worth making explicit: the critic loop is not universally necessary, and the system correctly identifies convergence on the first pass for the majority of well-formed queries, avoiding unnecessary re-computation.

\section{Limitations of the Current Implementation}
\label{sec:limitations}

Several limitations warrant acknowledgement. Pattern-based intent extraction is inherently brittle to novel phrasing; a query like ``compare revenue in Q1 with Q4'' exercises the YoY detection path but may not correctly map the quarter labels without additional detector logic. The historical baselines used by the plausibility calculator are synthetic and computed from the same synthetic dataset, which means the Z-score threshold has not been calibrated against real distribution shift. The critic loop's feedback message is a simple text description of the failing component; a more sophisticated approach would generate a targeted re-query suggestion, which Claude-class models could in principle do given sufficient prompt engineering. Finally, the current ontology is hard-coded in Python; migrating it to the JSON files in \texttt{data/ontology/} with live-reload would improve operational flexibility.

A related limitation is the absence of cross-database query support. Each query is routed to a single database (sales by default), which means a query asking for a comparison between HR salary budgets and sales revenue — a cross-domain join across \texttt{hr.db} and \texttt{sales.db} — cannot be answered by the current system. The five-database design was chosen specifically to model this real-world separation, but the query router does not yet implement a federation step.

The evaluation sample of ten queries, while carefully selected, is small by machine learning standards. Drawing strong statistical conclusions about the drift metric's behaviour from ten data points is inadvisable; the results should be treated as a feasibility demonstration rather than a statistically powered comparison. A larger-scale evaluation using a diverse query set generated from actual business analyst question logs would be needed to validate the 0.15 convergence threshold and the 0.40/0.30/0.30 weight allocation empirically.

A further consideration is the effect of query phrasing variation on the drift score. The ten test queries were each submitted once in a fixed phrasing; different phrasings of the same intent (e.g., ``what is total revenue by region net of tax'' versus ``show me gross revenue by region minus taxes'') might produce different intent alignment scores because the pattern detectors are sensitive to specific surface forms. In a more robust evaluation, each query would be submitted in several paraphrases and the drift score variance across phrasings would be reported alongside the mean. This robustness analysis was beyond the scope of the current evaluation but is a natural extension that would reveal which detector patterns are most sensitive to phrasing variation.

The synthetic nature of the evaluation data is a further limitation that affects the plausibility component specifically. Because the plausibility baselines were derived from the same synthetic dataset as the test queries, they represent a closed-loop validation: the system is essentially checking whether a query result is consistent with the distribution of results from other queries on the same synthetic data. In a real deployment, the historical baselines would be built from actual production query results, which may have different distributional properties — seasonal patterns, trend effects, data quality variations — that the synthetic distribution does not capture. The Z-score approach would still be applicable in that setting, but the thresholds (particularly the 2-sigma warning level) would need to be calibrated against the production distribution rather than assumed.

The pattern-matching intent extraction, while deterministic and auditable, has a coverage ceiling: it can only detect intent patterns that were explicitly anticipated by the detector authors. Business users often phrase queries in idiomatic ways that resist pattern matching — for instance, ``what's flying off the shelves?'' (top-selling products), ``where are we bleeding margin?'' (lowest gross margin by category), or ``give me the at-risk accounts'' (customers at churn risk). None of these phrasings would produce useful extractions from the current pattern set. A neural intent classifier or an LLM-based extraction step would handle these cases naturally, at the cost of determinism and auditability.

% =============================================================================
% CHAPTER 6 — CONCLUSION AND FUTURE WORK
% =============================================================================
\chapter{Conclusion and Future Work}
\label{ch:conclusion}

\section{Summary of Contributions}
\label{sec:contributions}

This project addressed a specific, well-defined gap in the Text-to-SQL literature: the absence of a reference-free, query-time mechanism for detecting and correcting semantic failures in deployed systems. Four primary contributions were made:

\begin{enumerate}[leftmargin=2em,label=\arabic*.]
  \item \textbf{A formalised semantic drift metric} — a reference-free, query-time composite score that decomposes intent fidelity into three independently measurable dimensions: intent alignment, constraint adherence, and result plausibility. The metric is grounded in the system's own intermediate representations and requires no ground-truth \ac{SQL} for computation.
  \item \textbf{An ontology-grounded five-agent pipeline} — a modular query processing architecture in which each agent has a narrow, auditable responsibility, and a shared \texttt{NetworkX} ontology graph provides the domain knowledge required for entity resolution, metric lookup, and constraint enforcement.
  \item \textbf{A critic-loop orchestrator} — an iterative refinement mechanism that uses the drift score as a continuous feedback signal, re-invoking the pipeline with error context until the query converges below the 0.15 threshold or five iterations are exhausted. This loop reduced mean drift from 0.187 (single-pass baseline) to 0.113 across the test query set.
  \item \textbf{A typed state machine design pattern for auditable agent pipelines} — the \texttt{QueryState} dataclass, with its explicit \texttt{add\_trace()} and \texttt{get\_trace\_summary()} methods, provides a reusable template for any multi-agent pipeline in which the execution history must be preserved for debugging and reporting. The pattern decouples inter-agent communication from agent logic, making each agent testable with a plain Python object rather than a full system harness.
\end{enumerate}

Collectively, these contributions represent a system-level rather than model-level advance. No new neural architecture is introduced; the improvement comes from wrapping an existing generation approach in a principled quality measurement and feedback loop. This framing has practical implications: the techniques apply equally whether the underlying SQL generator is a regex-based pattern matcher, a fine-tuned T5 model, or a general-purpose LLM — because the drift metric operates on the inputs and outputs of the generation step rather than on its internal workings.

The reproducibility of these contributions is supported by the project's open design: the synthetic data generation is deterministic given a fixed random seed; the pattern-matching intent parser produces the same output for a given query on every run; the Z-score plausibility calculation is a closed-form computation with no learned parameters; and the SQL builder is a pure function of its inputs. The exact drift scores reported in Table~\ref{tab:drift-results} can be reproduced from the source code and the synthetic dataset without access to any external API or pre-trained model. When the optional Claude integration is enabled, results will differ because \ac{LLM} outputs are non-deterministic; the deterministic baseline path reproduces exactly. This reproducibility was not an explicit project objective but an emergent property of the design choices made in favour of auditability and determinism.

\section{Key Findings}
\label{sec:findings}

The evaluation on ten business intelligence queries produced several findings worth highlighting. The first and perhaps most practically important finding is that the majority of well-formed business queries (6 of 10) converged on the first pipeline pass, suggesting that the pattern-matching intent parser and ontology mapper are sufficiently accurate for common query forms. Second, the constraint adherence component was the most reliable: the ConstraintValidatorAgent achieved perfect or near-perfect scores on nine of ten queries because the primary business rules are deterministic and the \ac{SQL} builder enforces them by default. Third, the plausibility component, while functional, was the most sensitive to the quality of the historical baseline: queries involving less common metrics (gross margin, return rate) produced noisier plausibility scores because the baseline distributions had wider variance.

The third finding concerns the robustness of the plausibility component: its scores were stable across iterations, varying by at most 0.05 between retry attempts on the same query. This stability is both a strength and a limitation. It is a strength because it means the plausibility component provides a consistent signal that does not oscillate between iterations, preventing the drift score from oscillating. It is a limitation because it means the plausibility component alone cannot drive convergence: if a query's intent and constraint components are perfect but plausibility is low, the critic loop will iterate without making progress, since none of the feedback messages specifically address the SQL structure that caused the plausibility anomaly.

The most problematic query type — return rate by product category — highlighted a structural gap: when a requested metric has no first-class representation in the ontology, the system falls back to a proxy metric that may satisfy the constraint rules but not the user's actual intent. This suggests that ontology coverage, rather than algorithmic refinement, is the primary lever for improving system performance on edge cases.

This finding carries a practical implication for system maintenance. Monitoring the distribution of drift scores over time in a production deployment would provide an early-warning signal for ontology coverage gaps: a cluster of queries with consistently elevated $D$ scores in a particular subdomain (say, return-rate calculations for a specific product line) would indicate that the ontology's metric definitions for that area are insufficient. A regular review of the top-10 highest-drift queries would be more actionable than a broad accuracy metric, because each high-drift query points to a specific, nameable failure mode that can be addressed by adding an alias, a metric template, or a constraint rule. This maintenance pattern — query-driven ontology improvement — is a form of active learning without the annotation overhead typically associated with active learning in machine learning contexts.

The evaluation also revealed an unexpected finding: the six queries that converged on the first pass did so with a mean intent alignment of $a_i = 0.882$, well above the 0.85 ``reliable'' threshold identified earlier. What distinguished these queries was not just high $a_i$ but also high $a_p$: all six had plausibility scores above 0.87. The joint condition ($a_i > 0.85$ and $a_p > 0.87$) was met by all six first-pass convergers and none of the multi-iteration queries. While the sample is too small to treat this as a formal decision rule, it suggests that a pre-flight plausibility check — computing plausibility on a cached or estimated result before running the full pipeline — could allow the system to skip the critic loop entirely for queries that are likely to converge immediately. This ``fast path'' optimisation would reduce mean latency for the majority of well-formed queries, at the cost of occasionally skipping refinement for borderline cases.

The role of the historical baseline in plausibility scoring deserves emphasis in the context of these findings. The pre-populated baselines used in this evaluation were constructed from a small set of canonical reference queries run against the synthetic dataset before the evaluation commenced. For the six simple-aggregation queries (revenue totals, order counts, customer tiers), the baselines were well-populated and the Z-scores remained below 2 on the first pass. For the four complex queries (YoY, category averages, gross margin, return rate), the baseline variance was higher — partly because these aggregations are computed over smaller subsets of the data (e.g., gross margin by supplier produces only one row per supplier) and partly because the proxy metric formulas differ slightly from the canonical metric formulas that were used to generate the baselines. This systematic bias between proxy and canonical metric formulas is the underlying cause of the elevated plausibility scores for these four queries, and it is a direct manifestation of the ontology coverage gap discussed in Section~\ref{sec:ontology-coverage}.

A secondary observation concerns the relative stability of the three drift components across iterations. Constraint adherence was the most volatile: it jumped between 0.67 and 1.00 in a single iteration for the gross margin query, because constraint satisfaction is near-binary — once the correct metric was identified, all relevant constraints were automatically satisfied by the SQL template. Intent alignment improved more gradually, typically gaining 0.05--0.12 per iteration as the pattern matchers accumulated evidence from the feedback message. Plausibility was the most stable: it rarely changed by more than 0.05 between iterations, suggesting that the generated SQL structure, not the aggregate value, was the primary driver of plausibility variation. These observations imply that future versions of the critic loop might prioritise constraint-oriented feedback in early iterations (where the gain is largest) and intent-oriented feedback in later iterations.

A broader observation is that the critic loop provided its most substantial benefit not on the most difficult queries (which did not converge regardless) but on the mid-range queries — those with initial drift in the 0.15--0.25 range. The three queries with initial drift above 0.20 (gross margin, return rate, monthly revenue trend) all benefited from the loop, but the return rate query still failed to converge after five iterations; for the other two, the loop achieved convergence in 2--3 iterations. The six queries with initial drift below 0.15 converged immediately. This pattern suggests that the critic loop's practical operating regime is the 0.15--0.25 initial drift range: below 0.15, no iteration is needed; above 0.25, the failure mode is typically a structural gap in the ontology that the loop cannot resolve through feedback alone. A deployment might therefore consider a three-tier response strategy: auto-serve queries with $D < 0.15$; apply the critic loop to queries with $0.15 \leq D < 0.25$; and route queries with $D \geq 0.25$ to a human analyst or an LLM-enhanced clarification flow.

\section{Future Work}
\label{sec:future-work}

The evaluation results and the implementation experience together suggest several concrete directions for future development. Rather than speculative enhancements, these directions each address a specific gap or limitation identified during the project.

Six directions for future development were identified from the project experience:

\begin{enumerate}[leftmargin=2em,label=\arabic*.]
  \item \textbf{Real embedding-based similarity} — replacing \texttt{SequenceMatcher} in the OntologyMapper with sentence embeddings (e.g., \texttt{all-MiniLM-L6-v2}) and a ChromaDB cosine collection. The ChromaDB integration path already exists in \texttt{ontology.py}; the remaining work is to populate the collection and activate the fallback. This change would substantially improve alias resolution for domain jargon not captured by the current string-matching approach.

  \item \textbf{PostgreSQL production deployment} — the dual-dialect \ac{SQL} builder is already implemented and tested. Deploying against a live PostgreSQL instance via a \texttt{DATABASE\_URL} environment variable would enable production-scale performance testing and remove the SQLite concurrency limitation.

  \item \textbf{Ontology-as-data migration} — moving the eleven entity definitions, twelve metrics, and nine constraint rules from \texttt{\_load\_hardcoded\_defaults()} to the JSON files in \texttt{data/ontology/}, with a file-watch reload mechanism. This would allow business analysts to extend the ontology without modifying Python source code.

  \item \textbf{LLM-enhanced natural language understanding} — integrating a Claude-class model as the primary intent parser rather than the fallback planner. Pattern matching would serve as a fast path; the \ac{LLM} would handle ambiguous or complex queries. The existing \texttt{ClaudeExecutionPlanner} architecture provides a pattern for this integration.

  \item \textbf{Multi-database federated queries} — the current system queries a single database per request. Adding a federation layer that distributes sub-queries across the five domain databases (sales, inventory, analytics, HR, finance) and joins the results in application code would dramatically expand the range of answerable questions.

  \item \textbf{Multi-language support} — extending the intent parser's regex patterns and the ontology's alias lists to cover at least one additional language (Hindi or Spanish, given the user base of many RVCE alumni). The \ac{API} layer is language-agnostic; only the pattern-matching layer requires additions.
\end{enumerate}

Each of these directions is non-trivial and represents a project-scale commitment in its own right. Priority among them depends on the deployment context: for an organisation already running PostgreSQL and with a data science team, the embedding-based similarity direction (Direction 1) would likely yield the fastest improvement in query coverage. For an organisation with non-technical business analysts as primary users, the LLM-enhanced NLU direction (Direction 4) would provide the largest reduction in failed queries on novel phrasing. The ontology-as-data migration (Direction 3) has the lowest technical risk and the highest immediate operational value, since it requires no model training or infrastructure changes beyond file-system conventions.

An eighth direction concerns the critic loop's feedback mechanism. The current design generates a natural language feedback string that the parser reads as free text; a more structured protocol — for example, a JSON object specifying which detector to adjust and by how much — would make the feedback machine-readable and allow each agent to respond to it programmatically without text parsing. This change would decouple the feedback format from the natural language understanding capability of the agent that receives it, making the protocol extensible to agents implemented with neural methods as well as pattern-based ones.

A seventh direction, not listed above but worth noting, is the development of a benchmark dataset for reference-free Text-to-SQL evaluation. The ten-query test set in this project was constructed manually; a community-maintained dataset of (query, expected-drift-label, justification) triples would enable systematic comparison of different drift metrics and threshold strategies across research groups. Creating such a dataset would require careful thought about what ``expected drift'' means without a reference SQL — perhaps crowd-sourced labels from domain experts who judge whether a query result faithfully answers the original question. This is a data collection and annotation challenge as much as a technical one, but it is a necessary precondition for the field to move from feasibility demonstrations to statistically powered evaluation of reference-free quality metrics.

\section{Broader Implications}
\label{sec:broader-implications}

The semantic drift problem examined in this project is a specific instance of a broader challenge in deployed AI systems: the gap between benchmark performance and production reliability. Text-to-SQL models that achieve 85\% accuracy on Spider are evaluated on questions drawn from the same distribution as their training data, with ground-truth SQL available for comparison. Production deployments face a fundamentally different problem: an open-ended query space, no ground-truth, and users who are often unaware of the system's limitations.

The approach developed here — decomposing reliability into measurable components, assigning weights that reflect each component's contribution to failure, and using the resulting score as a feedback signal — is not specific to Text-to-SQL. The same pattern could be applied to any system where a natural language input is transformed into a structured artefact: code generation (does the generated code implement the user's intent, satisfy naming conventions, and produce plausible outputs?), form filling (does the filled form reflect the stated requirements, comply with validation rules, and contain plausible values?), or document summarisation (does the summary preserve key facts, respect length constraints, and remain within the distributional norms for the document type?). In each case, the three-component decomposition of intent alignment, constraint adherence, and result plausibility offers a starting framework that can be specialised to the domain.

The specific design choices made in this project — pattern matching over neural NLU, Z-score over machine learning anomaly detection, deterministic constraint evaluation over LLM-based rule checking — were motivated by the requirement for full auditability and the constraint of a single-developer academic project timeline. In a production engineering context with a larger team, some of these choices would be revisited: learned embeddings would replace string matching, a trained anomaly detector would supplement Z-score, and an LLM would handle ambiguous phrasing. The framework, however — three components, composite score, iterative refinement — is independent of these implementation choices and would survive the transition from prototype to production.

Finally, the open-source nature of the codebase and the deterministic reproducibility of the evaluation results mean that the project is well positioned for replication. A researcher who clones the repository, installs the dependencies, and runs \texttt{python -m src.db.manager} followed by \texttt{uvicorn src.api.backend:app} will have a running system within minutes on any machine with Python 3.10 or later. The same ten test queries submitted through the Streamlit interface will produce drift scores within rounding error of those reported in Chapter 5, because every step of the pipeline is deterministic given the same synthetic dataset. This reproducibility lowers the barrier for future researchers to use the system as a baseline or a testbed for new approaches to reference-free Text-to-SQL quality evaluation.

% =============================================================================
% BACK MATTER
% =============================================================================
% Notes for the reader:
% All results in Chapter 5 are reproducible from the source code and synthetic dataset.
% The drift scores in Table 5.1 were recorded from the /query API endpoint responses.
% The latency measurements in Table 5.3 were derived from the application log files.
% The baseline scores in Table 5.4 were obtained by setting max_iterations=1 in config.py.

% A closing note on the project's reproducibility:
% All source code, requirements files, and synthetic data generation scripts
% are contained in the project repository. The full dataset can be reproduced
% from scratch by running: python -m src.db.manager
% The API and dashboard can be started with: docker compose up
% (from the docker/ directory)

\bibliographystyle{IEEEtran}
\bibliography{../references}

\appendix
\input{ApndxA}
\input{ApndxB}

% Appendix C — Selected Code Listings
\chapter{Selected Code Listings}
\label{apndx:code}

This appendix reproduces five representative code excerpts from the implementation, selected to illustrate the key algorithmic patterns in the system. The excerpts are condensed for readability — comments and blank lines are omitted — but the logic is faithful to the source files. Full source code is available in the project repository.

The five excerpts were chosen to cover the five major technical contributions described in Chapter~\ref{ch:conclusion}: the drift composite calculation (\texttt{SemanticDriftMetric.calculate()}), the SQL generation entry point (\texttt{\_build\_sql()}), the critic loop feedback constructor (\texttt{\_build\_feedback\_message()}), the ontology graph builder (\texttt{\_build\_ontology\_graph()}), and the plausibility calculator (\texttt{ResultPlausibilityCalculator}). Together, these five functions span the four source files that contain the project's novel algorithmic contributions: \texttt{semantic\_drift.py}, \texttt{pipeline.py}, \texttt{orchestrator.py}, and \texttt{ontology.py}. The intent is that a reader who understands these five excerpts understands the core mechanics of the system, even without reading the full source.

Notation conventions used in the code excerpts: ellipsis (\texttt{...}) denotes one or more omitted lines; inline comments beginning with \texttt{\#} are preserved when they explain non-obvious logic; import statements and class-level declarations are omitted unless needed for context. Line numbers in the comments refer to the original source files as committed to the project repository and may differ in a modified or future version of the codebase.

\section{Semantic Drift Composite Score (semantic\_drift.py)}

\begin{tcblisting}{listing only, colback=gray!10!white, breakable, boxsep=0pt,
  top=1mm, bottom=1mm, left=1mm, right=1mm,
  listing options={
    language=Python, basicstyle=\small\ttfamily,
    keywordstyle=\color{kwcolor}, commentstyle=\color{cmtcolor},
    stringstyle=\color{strcolor}, breaklines=true,
    numbers=left, numberstyle=\tiny\color{gray}}}
# SemanticDriftMetric.calculate() — condensed version
def calculate(self, state: 'QueryState') -> DriftMetricComponents:
    # 1. Intent alignment (40 %)
    confidences = {e: e.confidence for e in (state.intent.entities or [])}
    similarities = [m.similarity for m in (state.mappings or [])]
    a_i = IntentAlignmentCalculator.calculate_alignment_score(
        confidences, similarities)

    # 2. Constraint adherence (30 %)
    constraints = [c.__dict__ for c in (state.constraints or [])]
    a_c = ConstraintAdherenceCalculator.calculate_adherence_score(constraints)

    # 3. Result plausibility (30 %)
    a_p = ResultPlausibilityCalculator.calculate_result_plausibility(
        state.results, self.historical_baselines)

    # Composite drift (clamped to [0, 1])
    D = 0.4 * (1 - a_i) + 0.3 * (1 - a_c) + 0.3 * (1 - a_p)
    D = max(0.0, min(1.0, D))
    return DriftMetricComponents(a_i, a_c, a_p, D, datetime.now(),
        {"intent": a_i, "constraint": a_c, "plausibility": a_p})
\end{tcblisting}

\section{SQL Builder Core (pipeline.py)}

\begin{tcblisting}{listing only, colback=gray!10!white, breakable, boxsep=0pt,
  top=1mm, bottom=1mm, left=1mm, right=1mm,
  listing options={
    language=Python, basicstyle=\small\ttfamily,
    keywordstyle=\color{kwcolor}, commentstyle=\color{cmtcolor},
    stringstyle=\color{strcolor}, breaklines=true,
    numbers=left, numberstyle=\tiny\color{gray}}}
# _build_sql() — GROUP BY configuration look-up (pipeline.py line 246)
gbconf = GROUP_BY_CONFIG.get(group_by_key, {})
for j in gbconf.get("joins", []):
    joins_needed.add(j)

# Build SELECT: dimension column + metric expressions
select_parts = []
if group_by_key and gbconf:
    if "select" in gbconf:
        select_parts.append(gbconf["select"])
    elif f"select_{d}" in gbconf:
        select_parts.append(gbconf[f"select_{d}"])
for m in metric_list:
    if m in METRIC_SELECT:
        select_parts.append(METRIC_SELECT[m])
if not select_parts:
    select_parts.append(DEFAULT_SELECT)
\end{tcblisting}

\section{Feedback Message Construction (orchestrator.py)}

\begin{tcblisting}{listing only, colback=gray!10!white, breakable, boxsep=0pt,
  top=1mm, bottom=1mm, left=1mm, right=1mm,
  listing options={
    language=Python, basicstyle=\small\ttfamily,
    keywordstyle=\color{kwcolor}, commentstyle=\color{cmtcolor},
    stringstyle=\color{strcolor}, breaklines=true,
    numbers=left, numberstyle=\tiny\color{gray}}}
# orchestrator.py — _build_feedback_message() (condensed)
def _build_feedback_message(drift: DriftMetricComponents) -> str:
    parts = []
    if drift.intent_alignment < 0.75:
        parts.append(
            f"Intent alignment low ({drift.intent_alignment:.2f}): "
            "re-examine entity/metric extraction confidence"
        )
    if drift.constraint_adherence < 1.0:
        parts.append(
            f"Constraint adherence {drift.constraint_adherence:.2f}: "
            "check which constraint rules evaluated as unsatisfied "
            "and ensure the metric SQL template satisfies them"
        )
    if drift.result_plausibility < 0.8:
        parts.append(
            f"Plausibility low ({drift.result_plausibility:.2f}): "
            "result aggregates are outside expected range — "
            "verify join paths and WHERE conditions"
        )
    if not parts:
        parts.append("All components near threshold — minor refinement needed")
    return " | ".join(parts)
\end{tcblisting}

The feedback message is appended to \texttt{state.trace\_log} before the next pipeline iteration. The IntentParser reads the last trace entry at the start of its \texttt{process()} method and adjusts its confidence thresholds accordingly: specifically, if the message contains ``re-examine entity/metric extraction'', the parser lowers its entity pattern confidence threshold by 0.05 for that iteration, making it more likely to accept a lower-confidence alternative entity match. This conservative adjustment mechanism was chosen over a more aggressive re-parameterisation to prevent oscillation between two candidate interpretations across successive iterations.

The ``all components near threshold'' fallback at the end of \texttt{\_build\_feedback\_message()} addresses the edge case where drift is marginally above 0.15 but no single component is sufficiently below its own threshold to trigger a targeted message. In practice this case arose once during testing, on the monthly revenue trend query at iteration 2 ($D = 0.152$, $a_i = 0.86$, $a_c = 1.0$, $a_p = 0.83$). The generic feedback caused the parser to apply a small confidence adjustment that lowered $a_i$ slightly — not the intended direction — before recovering on iteration 3 ($a_i = 0.89$, $D = 0.122$). This illustrates a limitation of the single-text feedback format: without a structured re-parameterisation protocol, the parser has limited ability to respond to qualitative feedback that does not map to a specific detector pattern. A structured feedback contract (e.g., a JSON object specifying which detector to adjust and by how much) would be more reliable than a natural language string.

\section{Ontology Graph Construction (ontology.py)}

\begin{tcblisting}{listing only, colback=gray!10!white, breakable, boxsep=0pt,
  top=1mm, bottom=1mm, left=1mm, right=1mm,
  listing options={
    language=Python, basicstyle=\small\ttfamily,
    keywordstyle=\color{kwcolor}, commentstyle=\color{cmtcolor},
    stringstyle=\color{strcolor}, breaklines=true,
    numbers=left, numberstyle=\tiny\color{gray}}}
# BusinessOntology._build_ontology_graph() (ontology.py)
def _build_ontology_graph(self):
    self.graph = nx.DiGraph()
    for entity_id, entity in self.entities.items():
        self.graph.add_node(entity_id, **entity)
    for join in self.join_rules:
        self.graph.add_edge(
            join["from_entity"], join["to_entity"],
            join_type=join.get("join_type", "LEFT"),
            condition=join["condition"],
            weight=join.get("weight", 1.0)
        )

def get_join_path(self, from_entity: str, to_entity: str):
    try:
        return nx.shortest_path(self.graph, from_entity, to_entity,
                                weight="weight")
    except nx.NetworkXNoPath:
        return []
\end{tcblisting}

% =============================================================================
% End of report. All content complete.
\end{document}
